% Leçon 9 — Vocabulaire et compréhension de texte : importance de la pratique du sport — Chinois 1ère A — Cours signé OnBuch+
% Programme officiel MINESEC. Compiler avec : tectonic ZHA09-vocabulaire-importance-de-la-pratique-du-sport.tex
\newcommand{\DOCMATIERE}{Chinois}
\newcommand{\DOCNIVEAU}{1ère A}
\newcommand{\DOCMODULE}{Module 2 : Environnement et santé}
\newcommand{\DOCLECON}{ZHA09}
\newcommand{\DOCTITRE}{Leçon 9 — Vocabulaire et compréhension de texte : importance de la pratique du sport}
% OnBuch+ — préambule « cours signés » ENRICHI (compilé avec tectonic / XeLaTeX)
% Système visuel « Pop » d'OnBuch+ : fond crème (jamais blanc), bordures encre
% épaisses, ombres « dures » décalées, titres Archivo Black en capitales.
% Version MATHÉMATIQUES : graphes (pgfplots/tikz), boîtes pédagogiques
% (définition, propriété, méthode, attention, le savais-tu, exercice résolu,
% exercice, TP, bilan), page de garde et signature OnBuch+ ancrée en bas.
%
% Macros attendues dans main.tex AVANT \input{preamble} :
%   \DOCMATIERE  \DOCNIVEAU  \DOCMODULE  \DOCLECON (numéro)  \DOCTITRE
\documentclass[11pt]{article}

\usepackage{fontspec}
\usepackage[french]{babel} % français = langue principale ; le chinois via \zh{..} / textebox (xeCJK)
\usepackage{xeCJK}
\usepackage{pifont} % \ding{51} \ding{55} utilisés par les modèles
\usepackage[a4paper,margin=2cm,top=3.4cm,bottom=2.8cm,headheight=1.4cm,footskip=1.3cm]{geometry}
\usepackage{amsmath,amssymb,amsfonts}
\usepackage{mathtools} % \xmapsto, \coloneqq... souvent utilisés par les modèles
\usepackage{cancel} % simplifications barrées (bilans de forces)
% \abs{z} (module, valeur absolue) et \norm{u} (norme), utilisés par les modèles
\providecommand{\abs}{}\let\abs\relax\DeclarePairedDelimiter{\abs}{\lvert}{\rvert}
\providecommand{\norm}{}\let\norm\relax\DeclarePairedDelimiter{\norm}{\lVert}{\rVert}
\usepackage[table]{xcolor} % [table] : fournit aussi \rowcolors (alternance de lignes)
\usepackage{pagecolor}
\usepackage{tikz}
\usetikzlibrary{arrows.meta,positioning,calc,shapes.geometric,shapes.misc,shapes.arrows,decorations.pathmorphing,decorations.pathreplacing,decorations.markings,patterns,shadows,mindmap,trees,fit,backgrounds,matrix,intersections,angles,quotes}
\usepackage{pgfplots}
\usepackage[version=4]{mhchem} % équations nucléaires \ce{^{235}_{92}U}
\usepackage[european,siunitx]{circuitikz} % schémas électriques
\pgfplotsset{compat=1.18}
\usepgfplotslibrary{fillbetween,groupplots} % aires entre courbes (calcul intégral)
\usepackage{enumitem}
\usepackage{fancyhdr}
% [most] charge toutes les bibliothèques tcolorbox (listings, minted, raster,
% documentation...) et gonfle inutilement la mémoire TeX sur un document long
% (~300 boîtes) : on ne charge que ce qui est réellement utilisé.
\usepackage[skins,breakable]{tcolorbox}
\usepackage{graphicx}
\usepackage{colortbl}
\usepackage{booktabs}
\usepackage{tabularx}
\usepackage{diagbox} % case d'en-tête diagonale des tableaux à double entrée
% Colonnes extensibles centrée / gauche / droite (C, L, R), souvent utilisées par les modèles
\newcolumntype{C}{>{\centering\arraybackslash}X}
\newcolumntype{L}{>{\raggedright\arraybackslash}X}
\newcolumntype{R}{>{\raggedleft\arraybackslash}X}
\usepackage{array}
\usepackage{multicol}
\usepackage{multirow}
\usepackage{siunitx}
% parse-numbers=false : accepte aussi \SI{2,5 \times 10^{-3}}{...}
\sisetup{locale=FR,output-decimal-marker={,},group-minimum-digits=4,parse-numbers=false}
\DeclareSIUnit\ohm{\ensuremath{\symup{\Omega}}} % U+2126 absent de la police
\DeclareSIUnit\division{div} % réglages d'oscilloscope (ms/div, V/div)
\DeclareSIUnit\tr{tr} % tours (vitesse de rotation, ex. \SI{1500}{\tr\per\minute})
\DeclareSIUnit\molar{M} % concentration molaire (ex. \SI{3}{\molar}, \SI{1}{\micro\molar})
\DeclareSIUnit\an{an} % année (le modèle confond parfois avec \year, un compteur TeX) — ex. \SI{5}{\kilo\meter\per\an}
\DeclareSIUnit\FCFA{FCFA} % franc CFA (ex. \SI{500}{\FCFA\per\kilo\gram})
\DeclareSIUnit\degreeBrix{\textdegree Brix} % teneur en sucre d'un sirop/jus (réfractométrie)
\DeclareSIUnit\month{mois} % le modèle utilise parfois \month, un compteur TeX, comme unité de durée
\usepackage{qrcode}
% \degree : commande fréquemment utilisée par les modèles pour un angle ou une
% température en dehors de \SI{}{} (non fournie par les paquets chargés).
\providecommand{\degree}{^{\circ}}
% \qed (fin de démonstration), utilisé par les modèles sans amsthm
\providecommand{\qed}{\hfill$\square$}
% \barre : double barre des valeurs interdites dans un tableau de variations (tkz-tab)
\providecommand{\barre}{\ensuremath{\|}}
% environnement proof (amsthm non chargé) utilisé par les modèles
\ifdefined\proof\else\newenvironment{proof}[1][Démonstration]{\par\noindent\textit{#1.}\ }{\qed\par}\fi
\usepackage{lastpage}
\usepackage{hyperref}
\hypersetup{hidelinks}

% ============================================================
% Palette « Pop » OnBuch+ — mêmes hex que la classe P (plus_ui.dart)
% ============================================================
\definecolor{poporange}{HTML}{F59321}
\definecolor{popdark}{HTML}{E07A0C}
\definecolor{popcream}{HTML}{FFF6E8}
\definecolor{popink}{HTML}{1C1714}
\definecolor{poppurple}{HTML}{7A5AE0}
\definecolor{popgreen}{HTML}{2E9E6A}
\definecolor{poppink}{HTML}{E0457B}
\definecolor{popblue}{HTML}{2F7DE1}
\definecolor{popgold}{HTML}{C98A12}
\definecolor{popmuted}{HTML}{8A7F76}
\definecolor{popline}{HTML}{EADFD2}
\definecolor{popgray}{HTML}{8A7F76}
\colorlet{popgrey}{popgray}
% Alias tolérants (le modèle nomme parfois les couleurs différemment)
\colorlet{popwhite}{white}
\colorlet{popblack}{popink}
\colorlet{popred}{poppink}
\colorlet{popyellow}{popgold}
% Teintes claires (fonds de tableaux/encadrés) — le modèle nomme parfois
% les couleurs avec un suffixe L (ex. popblueL) sans les avoir définies.
\colorlet{poporangeL}{poporange!20}
\colorlet{popdarkL}{popdark!20}
\colorlet{poppurpleL}{poppurple!20}
\colorlet{popgreenL}{popgreen!20}
\colorlet{poppinkL}{poppink!20}
\colorlet{popblueL}{popblue!20}
\colorlet{popgoldL}{popgold!20}
\colorlet{popredL}{popred!20}
\colorlet{popgrayL}{popgray!20}
% Teintes claires pour fonds de tableaux / zones
\colorlet{popblueL}{popblue!10!white}
\colorlet{poporangeL}{poporange!14!white}
\colorlet{popgreenL}{popgreen!12!white}
\colorlet{poppurpleL}{poppurple!12!white}
\colorlet{poppinkL}{poppink!12!white}
\colorlet{popgoldL}{popgold!14!white}

\pagecolor{popcream}

% ============================================================
% Typographie
% ============================================================
\defaultfontfeatures{Ligatures=TeX}
\setmainfont{PlusJakartaSans-Medium.ttf}[Path=../../../fonts/,BoldFont=PlusJakartaSans-Bold.ttf]
% Symboles, API et emojis absents de la police principale : repli sur DejaVu Sans (caractères actifs)
\newfontface\symfont{DejaVuSans.ttf}[Path=../../../fonts/]
\makeatletter
\newcommand\symactive[1]{\catcode#1=13 \begingroup\lccode`\~=#1 \lowercase{\endgroup\def~}{{\symfont\char#1}}}
\newcommand\symblank[1]{\catcode#1=13 \begingroup\lccode`\~=#1 \lowercase{\endgroup\def~}{}}
\newcommand\symhyph[1]{\catcode#1=13 \begingroup\lccode`\~=#1 \lowercase{\endgroup\def~}{\mbox{-}}}
\newcount\symc
\newcommand\symrange[2]{\symc=#1 \@whilenum\symc<#2 \do{\expandafter\symactive\expandafter{\the\symc}\advance\symc by 1 }}
\newcommand\symblankrange[2]{\symc=#1 \@whilenum\symc<#2 \do{\expandafter\symblank\expandafter{\the\symc}\advance\symc by 1 }}
\makeatother
\symrange{"0250}{"02FF}\symactive{"2713}\symactive{"2714}\symactive{"2717}\symactive{"2718}\symactive{"2610}\symactive{"2611}\symactive{"2612}
\symrange{"2600}{"27BF}
\catcode"2705=13 \begingroup\lccode`\~="2705 \lowercase{\endgroup\def~}{{\symfont\char"2713}}\symblank{"26BD}\symblank{"26A1}
\symrange{"2460}{"257F}\symactive{"223C}\symactive{"2248}\symactive{"2260}
\symhyph{"2011}\symhyph{"2E1A}\symblankrange{"1F300}{"1FAFF}\symblank{"FE0F}
% Caractères chinois : WenQuanYi Zen Hei (sous-ensemble GB2312 + pinyin), gras/italique simulés
\setCJKmainfont{WenQuanYiZenHei-subset.ttf}[Path=../../../fonts/,AutoFakeBold=3,AutoFakeSlant=0.2]
\xeCJKsetup{CJKspace=false}
% Pinyin : voyelles à caron (ǎ ǐ ǒ ǔ ǖ ǘ ǚ ǜ) absentes de la police principale -> caractères actifs
\newfontface\pyfont{WenQuanYiZenHei-subset.ttf}[Path=../../../fonts/]
\newcommand\pyactive[1]{\catcode#1=13 \begingroup\lccode`\~=#1 \lowercase{\endgroup\def~}{{\pyfont\char#1}}}
\pyactive{"01CD}\pyactive{"01CE}\pyactive{"01CF}\pyactive{"01D0}\pyactive{"01D1}\pyactive{"01D2}\pyactive{"01D3}\pyactive{"01D4}\pyactive{"01D5}\pyactive{"01D6}\pyactive{"01D7}\pyactive{"01D8}\pyactive{"01D9}\pyactive{"01DA}\pyactive{"01DB}\pyactive{"01DC}
% Maths en sans-serif (Fira Math) avec chiffres et lettres droites de Plus
% Jakarta, pour que les formules se fondent dans le texte.
\usepackage[math-style=ISO,bold-style=ISO]{unicode-math}
\setmathfont{FiraMath-Regular.otf}[Path=../../../fonts/]
\setmathfont{PlusJakartaSans-Medium.ttf}[Path=../../../fonts/,range={up/num,up/{Latin,latin},\mathup}]
\usepackage{icomma} % virgule décimale sans espace en mode math
% Caractères mathématiques Unicode absents des polices de texte (Plus Jakarta,
% Archivo Black) : rendus par leur équivalent mathématique, en texte comme en
% formule — sinon ils s'affichent en carré vide (ex. « Arithmétique dans ℤ »).
\usepackage{newunicodechar}
\newunicodechar{ℝ}{\ensuremath{\mathbb{R}}}
\newunicodechar{ℤ}{\ensuremath{\mathbb{Z}}}
\newunicodechar{ℂ}{\ensuremath{\mathbb{C}}}
\newunicodechar{ℕ}{\ensuremath{\mathbb{N}}}
\newunicodechar{ℚ}{\ensuremath{\mathbb{Q}}}
\newunicodechar{𝔻}{\ensuremath{\mathbb{D}}}
\newunicodechar{α}{\ensuremath{\alpha}}
\newunicodechar{β}{\ensuremath{\beta}}
\newunicodechar{γ}{\ensuremath{\gamma}}
\newunicodechar{δ}{\ensuremath{\delta}}
\newunicodechar{ε}{\ensuremath{\varepsilon}}
\newunicodechar{θ}{\ensuremath{\theta}}
\newunicodechar{λ}{\ensuremath{\lambda}}
\newunicodechar{μ}{\ensuremath{\mu}}
\newunicodechar{σ}{\ensuremath{\sigma}}
\newunicodechar{τ}{\ensuremath{\tau}}
\newunicodechar{φ}{\ensuremath{\varphi}}
\newunicodechar{ψ}{\ensuremath{\psi}}
\newunicodechar{ω}{\ensuremath{\omega}}
\newunicodechar{Σ}{\ensuremath{\Sigma}}
% majuscules produites par \MakeUppercase dans les titres (α-aminés -> Α)
\newunicodechar{Α}{\ensuremath{\alpha}}
\newunicodechar{Β}{\ensuremath{\beta}}
\newunicodechar{Γ}{\ensuremath{\Gamma}}
\newunicodechar{Δ}{\ensuremath{\Delta}}
\newunicodechar{Ε}{\ensuremath{\varepsilon}}
\newunicodechar{Θ}{\ensuremath{\Theta}}
\newunicodechar{Λ}{\ensuremath{\Lambda}}
\newunicodechar{Μ}{\ensuremath{\mu}}
\newunicodechar{Τ}{\ensuremath{\tau}}
\newunicodechar{Φ}{\ensuremath{\Phi}}
\newunicodechar{Ψ}{\ensuremath{\Psi}}
\newunicodechar{Ω}{\ensuremath{\Omega}}
\newunicodechar{↦}{\ensuremath{\mapsto}}
\newunicodechar{∈}{\ensuremath{\in}}
\newunicodechar{∉}{\ensuremath{\notin}}
\newunicodechar{∧}{\ensuremath{\wedge}}
\newunicodechar{⇔}{\ensuremath{\Leftrightarrow}}
\newunicodechar{⟺}{\ensuremath{\Longleftrightarrow}}
\newunicodechar{⇒}{\ensuremath{\Rightarrow}}
\newunicodechar{⟹}{\ensuremath{\Longrightarrow}}
\newunicodechar{∘}{\ensuremath{\circ}}
\newunicodechar{∪}{\ensuremath{\cup}}
\newunicodechar{∩}{\ensuremath{\cap}}
\newunicodechar{≡}{\ensuremath{\equiv}}
\newunicodechar{⊂}{\ensuremath{\subset}}
\newunicodechar{⊥}{\ensuremath{\perp}}
\newunicodechar{∃}{\ensuremath{\exists}}
\newunicodechar{∀}{\ensuremath{\forall}}
\newunicodechar{∛}{\ensuremath{\sqrt[3]{\,}}}
\newunicodechar{∜}{\ensuremath{\sqrt[4]{\,}}}
\newunicodechar{⃗}{\ensuremath{{}^{\to}}}
\newunicodechar{ⁿ}{\ifmmode^{n}\else\textsuperscript{\ensuremath{n}}\fi}
\newunicodechar{ˣ}{\ifmmode^{x}\else\textsuperscript{\ensuremath{x}}\fi}
\newunicodechar{ᵇ}{\ifmmode^{b}\else\textsuperscript{\ensuremath{b}}\fi}
\newunicodechar{ʳ}{\ifmmode^{r}\else\textsuperscript{\ensuremath{r}}\fi}
\newunicodechar{ᵘ}{\ifmmode^{u}\else\textsuperscript{\ensuremath{u}}\fi}
\newunicodechar{ʸ}{\ifmmode^{y}\else\textsuperscript{\ensuremath{y}}\fi}
\newunicodechar{ᵃ}{\ifmmode^{a}\else\textsuperscript{\ensuremath{a}}\fi}
\newunicodechar{ᵉ}{\ifmmode^{e}\else\textsuperscript{\ensuremath{e}}\fi}
\newunicodechar{ᵏ}{\ifmmode^{k}\else\textsuperscript{\ensuremath{k}}\fi}
\newunicodechar{ᵖ}{\ifmmode^{p}\else\textsuperscript{\ensuremath{p}}\fi}
\newunicodechar{ᴺ}{\ifmmode^{N}\else\textsuperscript{\ensuremath{N}}\fi}
\newunicodechar{ᶜ}{\ifmmode^{c}\else\textsuperscript{\ensuremath{c}}\fi}
\newunicodechar{ʰ}{\ifmmode^{h}\else\textsuperscript{\ensuremath{h}}\fi}
\newunicodechar{ᵠ}{\ifmmode^{q}\else\textsuperscript{\ensuremath{q}}\fi}
\newunicodechar{ₙ}{\ifmmode_{n}\else\textsubscript{\ensuremath{n}}\fi}
\newunicodechar{ₐ}{\ifmmode_{a}\else\textsubscript{\ensuremath{a}}\fi}
\newunicodechar{ᵢ}{\ifmmode_{i}\else\textsubscript{\ensuremath{i}}\fi}
\newunicodechar{ᵣ}{\ifmmode_{r}\else\textsubscript{\ensuremath{r}}\fi}
\newunicodechar{ₖ}{\ifmmode_{k}\else\textsubscript{\ensuremath{k}}\fi}
\newunicodechar{ᵦ}{\ifmmode_{\beta}\else\textsubscript{\ensuremath{\beta}}\fi}
\newunicodechar{ₓ}{\ifmmode_{x}\else\textsubscript{\ensuremath{x}}\fi}
\newfontfamily\popdisplay{ArchivoBlack-Regular.ttf}[Path=../../../fonts/]
\newfontfamily\poplogo{SpaceGrotesk-Bold.ttf}[Path=../../../fonts/]
\newfontfamily\popsemi{PlusJakartaSans-SemiBold.ttf}[Path=../../../fonts/]
\newfontfamily\popxbold{PlusJakartaSans-ExtraBold.ttf}[Path=../../../fonts/]
\newfontfamily\popxboldls{PlusJakartaSans-ExtraBold.ttf}[Path=../../../fonts/,LetterSpace=3.0]

\setlist[enumerate]{leftmargin=1.7em,itemsep=5pt,topsep=4pt}
\setlist[itemize]{leftmargin=1.7em,itemsep=4pt,topsep=4pt,label={\color{poporange}\textbullet}}
\setlength{\parindent}{0pt}
\setlength{\parskip}{5pt}
\renewcommand{\arraystretch}{1.25}

% pgfplots : style Pop par défaut
\pgfplotsset{
  every axis/.append style={axis line style={popink,line width=1pt},tick style={popink},
    grid style={popline,line width=0.5pt},label style={font=\small},tick label style={font=\footnotesize}},
  popcurve/.style={poporange,line width=1.8pt,no markers,smooth},
}
% Même style disponible dans un \draw TikZ simple (hors pgfplots)
\tikzset{popcurve/.style={poporange,line width=1.8pt}}
\tikzset{popgrid/.style={popline,very thin}}
\colorlet{popcurve}{poporange} % le nom du style est parfois utilisé comme couleur

% ============================================================
% Pill / squiggle
% ============================================================
\newcommand{\poppill}[3]{%
  \tikz[baseline=-0.55ex]\node[draw=popink,line width=1.2pt,rounded corners=2.6pt,%
    fill=#2,inner xsep=6.5pt,inner ysep=3pt]{%
    \popxboldls\fontsize{7}{7}\selectfont\color{#3}\MakeUppercase{#1}};%
}
\newcommand{\popsquiggle}[1]{%
  \tikz[baseline=-0.4ex]{
    \draw[#1,line width=2.6pt,line cap=round]
      plot [smooth] coordinates {(0,0.09) (0.34,0.01) (0.68,0.21) (1.02,0.01) (1.36,0.10)};
  }%
}
% Mot-clé surligné (marqueur orange sous le texte)
\newcommand{\cle}[1]{{\popxbold\color{popdark}#1}}

% ============================================================
% En-tête / pied
% ============================================================
\pagestyle{fancy}
\fancyhf{}
\renewcommand{\headrulewidth}{0pt}
\renewcommand{\footrulewidth}{0pt}
\lhead{\raisebox{-0.15ex}{\poplogo\fontsize{13}{13}\selectfont\color{popink}OnBuch\color{poporange}+}}
\rhead{\poppill{\DOCMATIERE\ \textbullet\ \DOCNIVEAU\ \textbullet\ Leçon \DOCLECON}{white}{popink}}
\lfoot{\popxbold\fontsize{7.6}{7.6}\selectfont\color{popmuted}\MakeUppercase{Cours signé OnBuch+}\ \textbullet\ {\popsemi\DOCTITRE}}
\rfoot{\tikz[baseline=-0.6ex]\node[rounded corners=8.5pt,fill=popink,minimum height=17pt,minimum width=17pt,inner xsep=5pt,inner sep=0]{\popxbold\fontsize{8}{8}\selectfont\color{popcream}\thepage\,/\,\pageref*{LastPage}};}

% ============================================================
% Titres
% ============================================================
\newcommand{\chaptitre}[2]{%
  \begin{tcolorbox}[enhanced,colback=poporange,colframe=popink,boxrule=2.6pt,arc=11pt,
      drop shadow=popink,left=15pt,right=15pt,top=12pt,bottom=11pt,before skip=3pt,after skip=18pt]
    \poppill{#2}{popink}{popcream}\par\vspace{9pt}
    {\raggedright\hyphenpenalty=10000\exhyphenpenalty=10000\popdisplay\fontsize{22}{24}\selectfont\color{popcream}\MakeUppercase{#1}\par}
  \end{tcolorbox}
}
% Section numérotée : pastille orange avec le numéro + titre Archivo
\newcounter{popsec}
\newcommand{\coursec}[1]{%
  \stepcounter{popsec}\setcounter{popsub}{0}%
  \par\vspace{16pt}\noindent
  \begin{minipage}{\linewidth}
  \tikz[baseline=-0.7ex]\node[draw=popink,line width=1.6pt,fill=poporange,rounded corners=4pt,minimum size=22pt,inner sep=0]{\popdisplay\fontsize{12}{12}\selectfont\color{popcream}\thepopsec};%
  \hspace{8pt}{\popdisplay\fontsize{15}{17}\selectfont\color{popink}\MakeUppercase{#1}}\par
  \vspace{4pt}\noindent{\color{poporange}\rule{36pt}{3.6pt}}
  \end{minipage}\par\nopagebreak\vspace{8pt}
}
\newcounter{popsub}
\newcommand{\courssub}[1]{%
  \stepcounter{popsub}%
  \par\vspace{9pt}\noindent{\popxbold\fontsize{11.5}{13}\selectfont\color{popink}{\color{poporange}\thepopsec.\thepopsub}\hspace{6pt}#1}\par\nopagebreak\vspace{4pt}
}

% ============================================================
% Boîtes pédagogiques — même recette : fond blanc, bordure épaisse colorée,
% ombre dure encre, badge plein. Titre optionnel [#1].
% ============================================================
\tcbset{popbase/.style={breakable,enhanced,colback=white,boxrule=2.3pt,arc=9pt,
  shadow={3pt}{-3pt}{0pt}{popink},left=14pt,right=14pt,top=11pt,bottom=11pt,before skip=11pt,after skip=15pt,
  fontupper=\popsemi\fontsize{10.6}{14.5}\selectfont\color{popink},
  fontlower=\popsemi\fontsize{10.6}{14.5}\selectfont\color{popink}}}
% Coche dessinée (le glyphe ✓ n'existe pas dans les polices du document)
\newcommand{\popcheck}[1]{\tikz[baseline=-0.5ex]\draw[#1,line width=1.6pt,line cap=round,line join=round](-0.13,0.0)--(-0.04,-0.09)--(0.13,0.1);}
\providecommand{\coche}{\popcheck{popgreen}}
\providecommand{\resultat}[1]{\par\smallskip\noindent\fcolorbox{popgreen}{popgreenL}{\parbox{\dimexpr\linewidth-2\fboxsep-2\fboxrule\relax}{\textbf{Résultat.} #1}}\par\smallskip}
\newcommand{\popboxhead}[3]{\poppill{#1}{#2}{white}\ifx\relax#3\relax\else\hspace{6pt}{\popxbold\fontsize{10.6}{12}\selectfont\color{popink}#3}\fi\par\vspace{7pt}}

\newtcolorbox{aretenir}[1][]{popbase,colframe=popgreen,before upper={\popboxhead{À retenir}{popgreen}{#1}}}
% Texte en chinois (document de lecture, dialogue, extrait) : cadre
% d'étude. Un titre optionnel (source, auteur) s'affiche dans l'en-tête.
\newtcolorbox{textebox}[1][]{popbase,colframe=popblue,before upper={\popboxhead{文本 · Texte}{popblue}{#1}},after upper={}}
% Marques ✓ / ✗ (forme correcte / fautive) souvent utilisées par les modèles
\providecommand{\cross}{\textcolor{poppink}{\ensuremath{\times}}}
\providecommand{\xmark}{\cross}\providecommand{\crossmark}{\cross}\providecommand{\cmark}{\ensuremath{\checkmark}}
% Ligne de réponse / justification à compléter (inventée par les modèles)
\providecommand{\justification}[1]{\par\noindent\textit{Justification~:} \dotfill\par}
% \textipa (package tipa non chargé) : la police ne couvre pas l'API -> simple passage
\providecommand{\textipa}[1]{#1}
% Soulignement (ulem non chargé) : \uline, \ul
\providecommand{\uline}[1]{\underline{#1}}\providecommand{\ul}[1]{\underline{#1}}
% \glossaire{\item ...} inventée par les modèles : liste de glossaire
\providecommand{\glossaire}[1]{\begin{itemize}#1\end{itemize}}
% \alert (beamer) inventée par les modèles : texte mis en évidence
\providecommand{\alert}[1]{\textbf{\textcolor{poppink}{#1}}}
% variantes ASCII d'ordinaux écrites en macro
\providecommand{\iere}{\textsuperscript{re}}\providecommand{\ieme}{\textsuperscript{e}}\providecommand{\ere}{\textsuperscript{re}}\providecommand{\eme}{\textsuperscript{e}}\providecommand{\er}{\textsuperscript{er}}\providecommand{\ie}{\textsuperscript{e}}
% Ordinaux français écrits en macro par les modèles : 1\ière, 3\ième
\newcommand{\ière}{\textsuperscript{re}}\newcommand{\ième}{\textsuperscript{e}}
% \exemple (inventée par les modèles) : étiquette « Exemple : »
\providecommand{\exemple}{\textbf{Exemple~:}\ }
% \square utilisable aussi hors mode math (cases à cocher)
\AtBeginDocument{\let\squareMath\square\renewcommand{\square}{\ensuremath{\squareMath}}}
% Mot ou phrase en chinois à l'intérieur d'un paragraphe français
\newcommand{\zh}[1]{\ifmmode\text{#1}\else{#1}\fi}
\let\eng\zh \let\esp\zh \let\all\zh % alias tolérés
% flèches d'intonation inventées par les modèles
\providecommand{\textsearrow}{}\renewcommand{\textsearrow}{\ensuremath{\searrow}}\providecommand{\textnearrow}{}\renewcommand{\textnearrow}{\ensuremath{\nearrow}}
% \fr{..} : mot français (inventé par les modèles)
\providecommand{\fr}[1]{\textit{#1}}
% zhbox : encadré de texte chinois inventé par les modèles
\newenvironment{zhbox}{\begin{quote}}{\end{quote}}
% \vs : « versus » inventé par les modèles
\providecommand{\vs}{\textit{vs}\ }
% \blank : case à compléter inventée par les modèles
\providecommand{\blank}[1][]{\underline{\hspace{1.6em}}}
\newtcolorbox{exemplebox}[1][]{popbase,colframe=poporange,before upper={\popboxhead{Exemple}{poporange}{#1}}}
\newtcolorbox{definition}[1][]{popbase,colframe=popblue,before upper={\popboxhead{Définition}{popblue}{#1}}}
\newtcolorbox{propriete}[1][]{popbase,colframe=poppurple,before upper={\popboxhead{Propriété}{poppurple}{#1}}}
\newtcolorbox{methode}[1][]{popbase,colframe=popgold,before upper={\popboxhead{Méthode}{popgold}{#1}}}
\newtcolorbox{attention}[1][]{popbase,colframe=poppink,before upper={\popboxhead{Attention}{poppink}{#1}}}
\newtcolorbox{experience}[1][]{popbase,colframe=popblue,colback=popblueL,before upper={\popboxhead{Expérience}{popblue}{#1}}}
% « Le savais-tu ? » : carte violette pleine (contexte, culture, Cameroun)
\newtcolorbox{savaistu}[1][]{popbase,colback=poppurple,colframe=popink,
  fontupper=\popsemi\fontsize{10.4}{14.2}\selectfont\color{white},
  before upper={\poppill{Le savais-tu\,?}{white}{poppurple}\ifx\relax#1\relax\else\hspace{6pt}{\popxbold\color{white}#1}\fi\par\vspace{7pt}}}
% Exercice résolu : énoncé en haut, \tcblower puis la solution sur fond crème
\newcounter{popexo}\newcounter{popexr}
\newtcolorbox{exoresolu}[1][]{popbase,colframe=poporange,colbacklower=popcream,
  segmentation style={popink,line width=1.2pt,dashed},
  before upper={\stepcounter{popexr}\popboxhead{Exercice résolu \thepopexr}{poporange}{#1}},
  before lower={\poppill{Solution}{popink}{popcream}\par\vspace{6pt}}}
% Exercice (non résolu en place) — la correction est dans la partie Corrigés
\newtcolorbox{exercice}[1][]{popbase,colframe=popink,
  before upper={\stepcounter{popexo}\popboxhead{Exercice \thepopexo}{popink}{#1}}}
% Corrigé
\newtcolorbox{corrige}[1][]{popbase,colframe=popgreen,colback=popgreenL,
  before upper={\popboxhead{Corrigé}{popgreen}{#1}}}
% Objectifs / prérequis / bilan
\newtcolorbox{objectifs}{popbase,colframe=popink,colback=poporangeL,
  before upper={\popboxhead{Objectifs de la leçon}{popink}{}}}
\newtcolorbox{prerequis}{popbase,colframe=popink,colback=popblueL,
  before upper={\popboxhead{Prérequis}{popblue}{}}}
\newtcolorbox{bilan}{popbase,colframe=popink,colback=white,boxrule=2.8pt,
  before upper={\popboxhead{Fiche bilan}{poporange}{L'essentiel à connaître pour l'examen}}}
% Figure encadrée avec légende
\newtcolorbox{popfigure}[1][]{enhanced,breakable=false,colback=white,colframe=popline,boxrule=1.4pt,arc=8pt,
  left=10pt,right=10pt,top=10pt,bottom=8pt,before skip=10pt,after skip=12pt,halign=center,
  lower separated=false,halign lower=center,
  fontlower=\popsemi\fontsize{9}{11.5}\selectfont\color{popmuted},
  #1}
\newcommand{\legende}[1]{\par\vspace{6pt}{\popsemi\fontsize{9}{11.5}\selectfont\color{popmuted}#1}}

% ============================================================
% Signature OnBuch+
% ============================================================
\newcommand{\poplogoinline}[1]{{\poplogo\fontsize{#1}{#1}\selectfont\color{popink}OnBuch\color{poporange}+}}
\newcommand{\signatureonbuch}{%
  \begin{tcolorbox}[enhanced,colback=popink,colframe=popink,boxrule=2.2pt,arc=10pt,
      drop shadow=poporange,left=14pt,right=14pt,top=10pt,bottom=10pt,before skip=0pt,after skip=0pt]
    \tikz[baseline=-0.7ex]\node[circle,fill=popgreen,draw=popcream,line width=1.3pt,minimum size=22pt,inner sep=0]{\popcheck{white}};%
    \hspace{9pt}\begin{minipage}[c]{0.62\linewidth}
      {\popxbold\fontsize{12}{13}\selectfont\color{popcream}Signé\ \textbullet\ {\poplogo OnBuch\color{poporange}+}}\par\vspace{2pt}
      {\raggedright\popsemi\fontsize{8}{10.5}\selectfont\color{popline}\MakeUppercase{Équipe pédagogique OnBuch+}\\\MakeUppercase{Conforme au programme officiel MINESEC}\par}
    \end{minipage}\hfill
    {\poplogo\fontsize{22}{22}\selectfont\color{popcream}OnBuch\color{poporange}+}
  \end{tcolorbox}%
}

% ============================================================
% Page de garde
% #1 titre · #2 matière · #3 niveau · #4 module · #5 n° leçon · #6 durée ·
% #7 description / objectifs (texte ou itemize)
% ============================================================
\newcommand{\pagegarde}[7]{%
  \thispagestyle{empty}
  \newgeometry{a4paper,margin=1.7cm,top=1.6cm,bottom=1.4cm}%
  \begin{tikzpicture}[remember picture,overlay]
    \fill[poporange!18] ([xshift=-3.2cm,yshift=-3.6cm]current page.north east) circle (4.6cm);
    \fill[poppurple!14] ([xshift=2.4cm,yshift=7.6cm]current page.south west) circle (3.8cm);
  \end{tikzpicture}%
  {\poplogo\fontsize{30}{30}\selectfont\color{popink}OnBuch\color{poporange}+}\hfill
  \raisebox{0.6ex}{\poppill{Cours signé \textbullet\ Premium}{popink}{popcream}}\par
  \vspace{1pt}\popsquiggle{poppurple}\par
  \vspace{8pt}
  \begin{tcolorbox}[enhanced,colback=poporange,colframe=popink,boxrule=2.8pt,arc=13pt,
      drop shadow=popink,left=18pt,right=18pt,top=12pt,bottom=13pt,after skip=10pt]
    \poppill{#2 \textbullet\ #3}{popink}{popcream}\hspace{5pt}\poppill{#4}{white}{popink}\par\vspace{8pt}
    {\popdisplay\fontsize{14}{14}\selectfont\color{popink}LEÇON #5}\par\vspace{4pt}
    {\raggedright\hyphenpenalty=10000\exhyphenpenalty=10000\popdisplay\fontsize{25}{27}\selectfont\color{popcream}\MakeUppercase{#1}\par}
    \vspace{8pt}
    {\popxbold\fontsize{9.5}{9.5}\selectfont\color{popink}\MakeUppercase{Durée conseillée : #6}}
  \end{tcolorbox}
  \begin{tcolorbox}[enhanced,colback=white,colframe=popink,boxrule=2.2pt,arc=10pt,
      drop shadow=popink,left=16pt,right=16pt,top=10pt,bottom=10pt,after skip=8pt]
    \poppill{Dans cette leçon}{poppurple}{white}\par\vspace{5pt}
    \popsemi\fontsize{9.3}{12.6}\selectfont\color{popink}#7
  \end{tcolorbox}
  \noindent\begin{minipage}[t]{0.487\linewidth}
    \begin{tcolorbox}[enhanced,colback=popgreen,colframe=popink,boxrule=2.2pt,arc=10pt,
        drop shadow=popink,left=11pt,right=11pt,top=10pt,bottom=10pt,height=3.1cm,valign=center]
      \tcbox[colback=white,colframe=popink,boxrule=1.3pt,arc=5pt,boxsep=0pt,nobeforeafter,
        left=3pt,right=3pt,top=3pt,bottom=3pt,valign=center]{\qrcode[height=1.9cm]{https://play.google.com/store/apps/details?id=com.luvvix.onbuch&hl=fr}}%
      \hspace{8pt}\begin{minipage}[c]{0.5\linewidth}
        {\popdisplay\fontsize{11}{12.5}\selectfont\color{white}\MakeUppercase{Télécharge OnBuch}}\par\vspace{4pt}
        {\raggedright\popsemi\fontsize{8.4}{11}\selectfont\color{white}Gratuit \textbullet\ Google Play\\Tuteur IA, annales, concours, résultats\par}
      \end{minipage}
    \end{tcolorbox}
  \end{minipage}\hfill
  \begin{minipage}[t]{0.487\linewidth}
    \begin{tcolorbox}[enhanced,colback=poppurple,colframe=popink,boxrule=2.2pt,arc=10pt,
        drop shadow=popink,left=11pt,right=11pt,top=10pt,bottom=10pt,height=3.1cm,valign=center]
      \tikz[baseline=-0.7ex]\node[circle,fill=white,draw=popink,line width=1.3pt,minimum size=1.6cm,inner sep=0]{\popdisplay\fontsize{24}{24}\selectfont\color{poppurple}+};%
      \hspace{8pt}\begin{minipage}[c]{0.6\linewidth}
        {\popdisplay\fontsize{11}{12.5}\selectfont\color{white}\MakeUppercase{Passe à OnBuch+}}\par\vspace{4pt}
        {\raggedright\popsemi\fontsize{8.4}{11}\selectfont\color{white}Tous les cours signés, Tuteur IA illimité, fiches PDF — onglet OnBuch+ de l'app\par}
      \end{minipage}
    \end{tcolorbox}
  \end{minipage}
  \vspace{8pt}
  \popcontenu
  \vfill
  \signatureonbuch
  \restoregeometry
  \newpage
}

% Rangée « Ce cours contient » (page de garde)
\newcommand{\poptile}[3]{% #1 couleur #2 gros texte #3 légende
  \begin{tcolorbox}[enhanced,colback=#1,colframe=popink,boxrule=2pt,arc=9pt,shadow={2.5pt}{-2.5pt}{0pt}{popink},
      left=6pt,right=6pt,top=9pt,bottom=9pt,halign=center,valign=center,height=1.75cm,width=\linewidth,nobeforeafter]
    {\popdisplay\fontsize{12.5}{13}\selectfont\color{white}\MakeUppercase{#2}}\par\vspace{3pt}
    {\centering\popsemi\fontsize{7.8}{9.5}\selectfont\color{white}#3\par}
  \end{tcolorbox}}
\newcommand{\popcontenu}{%
  \par\noindent{\popxboldls\fontsize{7.5}{7.5}\selectfont\color{popmuted}CE COURS SIGNÉ CONTIENT}\par\vspace{7pt}
  \noindent\begin{minipage}[t]{0.235\linewidth}\poptile{popblue}{Cours}{complet, illustré, pas à pas}\end{minipage}\hfill
  \begin{minipage}[t]{0.235\linewidth}\poptile{poporange}{Méthodes}{et exercices résolus}\end{minipage}\hfill
  \begin{minipage}[t]{0.235\linewidth}\poptile{poppink}{Exercices}{type examen + corrigés}\end{minipage}\hfill
  \begin{minipage}[t]{0.235\linewidth}\poptile{popgreen}{Fiche bilan}{l'essentiel à retenir}\end{minipage}\par}

% Sommaire
\newtcolorbox{sommaire}{enhanced,colback=white,colframe=popink,boxrule=2.2pt,arc=10pt,
  drop shadow=popink,left=18pt,right=18pt,top=13pt,bottom=13pt,before skip=6pt,after skip=20pt,
  before upper={\poppill{Sommaire}{poppurple}{white}\par\vspace{9pt}\popsemi\fontsize{10.6}{14.5}\selectfont\color{popink}}}

% Signature de fin de cours — poussée tout en bas de la dernière page
\newcommand{\findecours}{%
  \par\vfill\nopagebreak
  \begin{center}\popsquiggle{poporange}\end{center}\vspace{4pt}
  \signatureonbuch
}
\AtBeginDocument{\def\llbracket{\mathopen{[\mkern-3.5mu[}}\def\rrbracket{\mathclose{]\mkern-3.5mu]}}} % intervalles d'entiers (glyphe ⟦ absent de la police)
% Glyphes absents de Fira Math : redéfinis à partir de glyphes disponibles
\AtBeginDocument{%
  \def\setminus{\mathbin{\backslash}}\let\smallsetminus\setminus
  \def\vdots{\mathord{\vbox{\baselineskip3.2pt\lineskiplimit0pt\kern4pt\hbox{.}\hbox{.}\hbox{.}}}}%
  \def\ddots{\mathinner{\mkern1mu\raise6.5pt\hbox{.}\mkern2mu\raise3.6pt\hbox{.}\mkern2mu\raise0.7pt\hbox{.}\mkern1mu}}%
  \def\triangle{\mathord{\scalebox{0.9}{$\symup{\Delta}$}}}%
  \def\nabla{\mathord{\raisebox{\depth}{\scalebox{1}[-1]{$\symup{\Delta}$}}}}%
  \def\checkmark{\popcheck{popdark}}%
  \def\star{\mathbin{\ast}}\def\bigstar{\mathord{\ast}}%
  \def\diamond{\mathbin{\scalebox{0.6}{$\Diamond$}}}%
  \def\bigcup{\mathop{\mathchoice{\raisebox{-0.3em}{\scalebox{1.7}{$\cup$}}}{\scalebox{1.25}{$\cup$}}{\cup}{\cup}}\displaylimits}%
  \def\bigcap{\mathop{\mathchoice{\raisebox{-0.3em}{\scalebox{1.7}{$\cap$}}}{\scalebox{1.25}{$\cap$}}{\cap}{\cap}}\displaylimits}%
  \def\lesseqqgtr{\mathrel{\lessgtr}}\def\bigoplus{\mathop{\oplus}\displaylimits}%
}
\providecommand{\ssi}{si et seulement si} % abréviation fréquente
\AtBeginDocument{\providecommand{\wideparen}{\overparen}} % arc de cercle

\begin{document}

\pagegarde{Leçon 9 — Vocabulaire et compréhension de texte : importance de la pratique du sport}{Chinois}{1ère A}{Module 2 : Environnement et santé}{ZHA09}{2 h}{Cette leçon de vocabulaire et compréhension de texte porte sur l'importance de la pratique du sport. À partir d'un texte chinois original, les élèves acquièrent le lexique du sport et de la santé, étudient les radicaux et l'ordre des traits de caractères clés, puis répondent à des questions de compréhension.
\vspace{4pt}
\begin{multicols}{2}\small
\begin{itemize}
\item À la fin de cette leçon, je sais lire à voix haute un court texte chinois sur le sport avec une prononciation correcte (pinyin et tons).
\item À la fin de cette leçon, je sais reconnaître et traduire au moins 15 mots de vocabulaire liés au sport et à la santé.
\item À la fin de cette leçon, je sais identifier les radicaux (clés) des caractères du thème et les utiliser pour deviner le sens.
\item À la fin de cette leçon, je sais écrire correctement quelques caractères du thème en respectant l'ordre des traits.
\item À la fin de cette leçon, je sais répondre à des questions de compréhension portant sur le texte support.
\item À la fin de cette leçon, je sais employer les collocations usuelles du thème (运动, 锻炼, 健康...) dans des phrases simples.
\end{itemize}\end{multicols}}

\begin{sommaire}
\begin{multicols}{2}
\begin{enumerate}
\item Mise en route et découverte du texte
\item Glossaire du thème : sport et santé
\item Clés (radicaux) et ordre des traits
\item Compréhension détaillée du texte
\item Collocations et exemples d'emploi
\item Production guidée et synthèse
\item Activité d'intégration
\item Méthodes et exercices résolus
\item Exercices
\item Corrigés des exercices
\item Fiche bilan
\end{enumerate}
\end{multicols}
\end{sommaire}

\begin{objectifs}
\begin{itemize}
\item À la fin de cette leçon, je sais lire à voix haute un court texte chinois sur le sport avec une prononciation correcte (pinyin et tons).
\item À la fin de cette leçon, je sais reconnaître et traduire au moins 15 mots de vocabulaire liés au sport et à la santé.
\item À la fin de cette leçon, je sais identifier les radicaux (clés) des caractères du thème et les utiliser pour deviner le sens.
\item À la fin de cette leçon, je sais écrire correctement quelques caractères du thème en respectant l'ordre des traits.
\item À la fin de cette leçon, je sais répondre à des questions de compréhension portant sur le texte support.
\item À la fin de cette leçon, je sais employer les collocations usuelles du thème (运动, 锻炼, 健康...) dans des phrases simples.
\item À la fin de cette leçon, je sais produire un court paragraphe guidé sur l'importance du sport dans ma vie.
\end{itemize}
\end{objectifs}

\begin{prerequis}
\begin{itemize}
\item Vocabulaire de base de la famille, de l'école et des activités quotidiennes (Seconde)
\item Structures de phrase simples : sujet + verbe + complément ; 很 + adjectif
\item Verbes modaux déjà vus : 会, 能, 应该
\item Notions de pinyin et des quatre tons
\item Radicaux fréquents déjà rencontrés (人, 口, 心, 氵)
\end{itemize}
\end{prerequis}

\newpage

\coursec{Mise en route et découverte du texte}

\courssub{Brainstorming : quels sports pratiquez-vous ?}

C'est la saison des vacances à Yaoundé : les terrains de quartier se remplissent de jeunes qui jouent au football, tandis que d'autres passent leurs journées sur leur téléphone. Au lycée, le médecin scolaire rappelle que beaucoup d'élèves manquent d'exercice physique. En Chine, chaque matin, les élèves font de la gymnastique collective dans la cour de l'école. Comment dire en chinois que le sport est bon pour la santé ? Cette leçon nous donne les mots et le texte pour en parler.

\textbf{Problématique :} comment lire et comprendre un court texte chinois sur l'importance du sport, et comment employer le vocabulaire du sport et de la santé ?

Avant d'ouvrir le texte, activons ce que nous savons déjà. Répondez d'abord en français, puis cherchez l'équivalent chinois dans le tableau ci-dessous.

\begin{itemize}
\item Quels sports pratiquez-vous ou connaissez-vous ? (football, course, danse, basket-ball, natation...)
\item Pourquoi fait-on du sport ? (pour la santé, pour le plaisir, pour être fort...)
\item Quand et où fait-on du sport ? (le matin, le soir, au stade, dans la cour...)
\end{itemize}

\begin{center}
\begin{tabularx}{\linewidth}{|X|X|X|X|}
\hline
\rowcolor{popblueL} Caractères & Pinyin & Nature & Traduction \\
\hline
运动 & yùndòng & nom / verbe & sport ; faire du sport \\
\hline
踢足球 & tī zúqiú & locution verbale & jouer au football \\
\hline
跑步 & pǎobù & verbe & courir, faire du jogging \\
\hline
跳舞 & tiàowǔ & verbe & danser \\
\hline
打太极拳 & dǎ tàijíquán & locution verbale & faire du tai-chi \\
\hline
锻炼 & duànliàn & verbe & s'entraîner, faire de l'exercice \\
\hline
身体 & shēntǐ & nom & corps, santé \\
\hline
健康 & jiànkāng & adjectif / nom & sain ; santé \\
\hline
重要 & zhòngyào & adjectif & important \\
\hline
快乐 & kuàilè & adjectif & heureux, joyeux \\
\hline
\end{tabularx}
\end{center}

\begin{attention}[Faux amis et pièges des francophones]
En français, on dit « jouer \emph{au} football » avec un seul verbe ; en chinois, chaque sport a son verbe propre : \zh{踢} (tī, donner un coup de pied) pour le football, \zh{打} (dǎ, frapper) pour le tai-chi, le basket ou le ping-pong, \zh{跳} (tiào, sauter) pour la danse. On ne dit jamais \zh{玩足球}. De plus, \zh{运动} se prononce avec deux quatrièmes tons : \emph{yùndòng}, bien marqués, descendants.
\end{attention}

\courssub{Carte mentale du thème}

Pour organiser le vocabulaire avant la lecture, voici une carte mentale autour du mot central \zh{运动} (yùndòng, sport). Elle vous servira de plan mental pendant toute la leçon.

\begin{popfigure}
\begin{tikzpicture}[mindmap, grow cyclic, every node/.style={concept}, concept color=popblue!40, text=popdark, root concept/.append style={concept color=poporange!60, font=\large\bfseries}, level 1/.append style={level distance=3.2cm, sibling angle=90}]
\node[align=center]{运动\\ yùndòng\\ (sport)}
  child[concept color=popgreen!40] { node[align=center]{好处\\ hǎochu\\ bienfaits} }
  child[concept color=poppurple!40] { node[align=center]{运动项目\\ yùndòng xiàngmù\\ les sports} }
  child[concept color=poppink!40] { node[align=center]{时间\\ shíjiān\\ quand ?} }
  child[concept color=popgold!50] { node[align=center]{人\\ rén\\ qui ?} };
\end{tikzpicture}
\legende{Carte mentale du thème : le sport, ses bienfaits, les disciplines, le moment et les personnes.}
\end{popfigure}

Retenez cette organisation : un bon texte chinois simple répond presque toujours à quatre questions : \emph{de quoi parle-t-on ? qui ? quand et où ? pourquoi ?} Le texte que nous allons lire suit exactement cette logique.

\courssub{Texte support : 运动很重要}

L'enseignant lit le texte une première fois à voix haute, à vitesse normale. Les élèves écoutent sans lire, puis suivent le pinyin du doigt lors de la deuxième lecture. Ensuite, lecture chorale (toute la classe ensemble), puis lecture individuelle : chaque élève lit une phrase, l'enseignant corrige les tons.

\begin{textebox}[运动很重要 — Le sport est important]
运动很重要。每天运动对我们的身体很好。运动可以让我们更健康、更快乐。很多中国人早上在公园里锻炼：他们跑步、打太极拳、跳舞。学生也应该多运动。运动以后，我们学习得更好。所以，大家都喜欢运动。\\[4pt]
\emph{Yùndòng hěn zhòngyào. Měitiān yùndòng duì wǒmen de shēntǐ hěn hǎo. Yùndòng kěyǐ ràng wǒmen gèng jiànkāng, gèng kuàilè. Hěn duō Zhōngguórén zǎoshang zài gōngyuán li duànliàn : tāmen pǎobù, dǎ tàijíquán, tiàowǔ. Xuésheng yě yīnggāi duō yùndòng. Yùndòng yǐhòu, wǒmen xuéxí de gèng hǎo. Suǒyǐ, dàjiā dōu xǐhuan yùndòng.}\\[4pt]
« Le sport est très important. Faire du sport chaque jour est très bon pour notre corps. Le sport peut nous rendre plus sains et plus heureux. Beaucoup de Chinois s'entraînent le matin dans les parcs : ils courent, font du tai-chi, dansent. Les élèves aussi devraient faire plus de sport. Après le sport, nous apprenons mieux. C'est pourquoi tout le monde aime le sport. »
\end{textebox}

\begin{aretenir}[Mini-glossaire de première lecture]
\begin{itemize}
\item \zh{对……很好} (duì... hěn hǎo) : être très bon pour...
\item \zh{让} (ràng) : rendre, faire que (+ personne + adjectif)
\item \zh{更} (gèng) : encore plus
\item \zh{应该} (yīnggāi) : devoir (conseil)
\item \zh{以后} (yǐhòu) : après
\item \zh{所以} (suǒyǐ) : c'est pourquoi, donc
\item \zh{大家} (dàjiā) : tout le monde
\end{itemize}
\end{aretenir}

\courssub{Travail sur la prononciation : les tons du texte}

La lecture à voix haute est l'occasion de corriger les tons. Observez les groupes difficiles du texte :

\begin{center}
\begin{tabularx}{\linewidth}{|X|X|X|}
\hline
\rowcolor{popblueL} Mot & Tons & Conseil de prononciation \\
\hline
运动 yùndòng & 4 + 4 & deux tons descendants, secs et nets \\
\hline
重要 zhòngyào & 4 + 4 & même mélodie que 运动, ne pas adoucir \\
\hline
很好 hěn hǎo & 3 + 3 & le premier 3\textsuperscript{e} ton devient un 2\textsuperscript{e} ton : \emph{hén hǎo} \\
\hline
可以 kěyǐ & 3 + 3 & même règle : prononcer \emph{kéyǐ} \\
\hline
早上 zǎoshang & 3 + neutre & le deuxième caractère est atone, court et léger \\
\hline
喜欢 xǐhuan & 3 + neutre & même remarque : \emph{huan} est atone \\
\hline
\end{tabularx}
\end{center}

\begin{attention}[La règle des deux troisièmes tons]
Quand deux troisièmes tons se suivent, le premier se prononce comme un deuxième ton (montant). On écrit toujours les tons d'origine (\emph{hěn hǎo}) mais on prononce \emph{hén hǎo}. C'est une des erreurs les plus fréquentes des francophones, qui lisent les tons écrits sans appliquer cette transformation.
\end{attention}

\courssub{Repérage global : de quoi parle le texte ?}

Après deux ou trois lectures, répondez aux questions de repérage, sans encore chercher le détail :

\begin{enumerate}
\item De quoi parle le texte ? (le thème général)
\item Qui fait du sport dans ce texte ? (les personnes mentionnées)
\item Où et quand font-ils du sport ?
\item Quelle est l'idée principale défendue par l'auteur ?
\end{enumerate}

\begin{exoresolu}[Repérage global du texte]
Énoncé : répondez en français aux quatre questions de repérage ci-dessus, en citant le mot chinois qui justifie chaque réponse.
\tcblower
Solution détaillée :
\begin{enumerate}
\item Le texte parle du \cle{sport} et de son importance : le mot \zh{运动} (yùndòng) revient dans presque chaque phrase, et le titre est \zh{运动很重要} (le sport est très important).
\item Les personnes qui font du sport sont \zh{很多中国人} (beaucoup de Chinois), \zh{学生} (les élèves) et, à la fin, \zh{大家} (tout le monde), ce qui inclut \zh{我们} (nous).
\item Ils font du sport \zh{早上} (le matin) \zh{在公园里} (dans les parcs) ; la structure \zh{在 + lieu} indique le lieu de l'action.
\item L'idée principale : le sport est bon pour le corps et pour le moral (\zh{更健康、更快乐} : plus sain, plus heureux), et il aide même à mieux étudier (\zh{我们学习得更好} : nous apprenons mieux). La conclusion \zh{所以，大家都喜欢运动} (c'est pourquoi tout le monde aime le sport) résume le message.
\end{enumerate}
\end{exoresolu}

\courssub{Méthode : lire un texte chinois en trois étapes}

\begin{methode}[Les trois lectures d'un texte chinois]
\begin{enumerate}
\item \textbf{Lecture globale} : lisez le texte entier une fois, sans dictionnaire, pour saisir l'idée générale. Repérez les mots déjà connus et les mots qui reviennent souvent (ici \zh{运动}). Ne vous bloquez pas sur un caractère inconnu.
\item \textbf{Lecture détaillée} : relisez phrase par phrase avec le glossaire. Soulignez les structures (sujet + verbe + complément, \zh{在 + lieu}, \zh{以后}...). Traduisez mentalement chaque phrase.
\item \textbf{Relecture de vérification} : relisez une dernière fois pour répondre aux questions de compréhension et vérifier que vos réponses correspondent bien au texte, pas à vos idées personnelles.
\end{enumerate}
\end{methode}

\begin{experience}[Lecture en chaîne et mini-sondage]
\begin{enumerate}
\item Lecture en chaîne : chaque élève lit une phrase du texte à voix haute ; la classe corrige les tons ensemble avant de passer à la phrase suivante.
\item Mini-sondage : par groupes de quatre, posez-vous la question \zh{你喜欢什么运动？} (Nǐ xǐhuan shénme yùndòng ? — Quel sport aimes-tu ?) et répondez avec \zh{我喜欢……} (j'aime...) suivi d'un sport du tableau. Un rapporteur annonce le sport préféré du groupe.
\end{enumerate}
\end{experience}

\begin{savaistu}[La gymnastique matinale en Chine]
En Chine, la plupart des écoles organisent chaque matin une séance de gymnastique collective (\zh{早操}, zǎocāo) dans la cour : des centaines d'élèves font les mêmes mouvements en musique. Dans les parcs des villes chinoises, on voit aussi, dès six heures du matin, des groupes de retraités pratiquer le tai-chi (\zh{太极拳}) ou la danse en groupe. Au Cameroun, c'est plutôt le football du soir et les matchs de quartier qui rassemblent les jeunes : deux façons différentes, mais un même constat — \zh{运动很重要} !
\end{savaistu}

\begin{exercice}[À vous de jouer]
\begin{enumerate}
\item Relevez dans le texte trois mots contenant le caractère \zh{运} ou le mot \zh{运动} et traduisez-les.
\item Recopiez la phrase \zh{运动可以让我们更健康、更快乐} avec son pinyin et sa traduction.
\item Répondez en chinois : \zh{你早上锻炼吗？} (Fais-tu de l'exercice le matin ?) en utilisant \zh{我早上……} ou \zh{我不……}.
\end{enumerate}
\end{exercice}

\begin{corrige}[Exercice « À vous de jouer »]
\begin{enumerate}
\item \zh{运动} (yùndòng, sport), \zh{运动很重要} (le sport est important), \zh{多运动} (faire plus de sport). Le caractère \zh{运} porte l'idée de « mouvement, transporter ».
\item \zh{运动可以让我们更健康、更快乐。} — \emph{Yùndòng kěyǐ ràng wǒmen gèng jiànkāng, gèng kuàilè.} — « Le sport peut nous rendre plus sains et plus heureux. »
\item Modèles de réponse : \zh{我早上跑步。} (Wǒ zǎoshang pǎobù. — Le matin, je cours.) ou \zh{我不锻炼。} (Wǒ bù duànliàn. — Je ne fais pas d'exercice.) Attention : la négation \zh{不} se place avant le verbe, jamais après.
\end{enumerate}
\end{corrige}

Nous avons maintenant une vision globale du texte : son thème, ses personnages, son message. La prochaine étape consiste à entrer dans le détail du vocabulaire du sport et de la santé, mot par mot, avec leurs clés et leur écriture.

\coursec{Glossaire du thème : sport et santé}

Dans la section précédente, nous avons découvert le texte support sur l'importance de la pratique du sport : vous avez entendu et lu pour la première fois des mots comme \zh{运动}, \zh{锻炼} ou \zh{健康}. Certains vous ont peut-être semblé familiers, d'autres complètement nouveaux. C'est normal : un texte ne se comprend vraiment que lorsqu'on en maîtrise le vocabulaire. Cette deuxième section est donc consacrée à l'étude systématique du lexique du sport et de la santé. Nous allons observer chaque mot, le prononcer correctement, comprendre sa nature grammaticale et apprendre à l'employer dans des phrases authentiques. Prenez votre cahier de vocabulaire : à la fin de cette section, vous devrez être capable de reconnaître, prononcer et utiliser chacun de ces mots sans hésitation.

\courssub{Observer avant d'apprendre : les mots dans leur contexte}

Avant de présenter le tableau de vocabulaire, observons quelques phrases tirées du texte de la section 1. Lisez-les à voix haute et essayez de deviner le sens des mots en gras.

\begin{exemplebox}[Les mots du thème en contexte]
\zh{我\textbf{每天}\textbf{早上}在\textbf{公园}\textbf{跑步}。}\\
\emph{Wǒ měitiān zǎoshang zài gōngyuán pǎobù.}\\
« Chaque matin, je cours dans le parc. »

\zh{\textbf{运动}对\textbf{身体}很\textbf{好}，所以我们\textbf{应该}经常\textbf{锻炼}。}\\
\emph{Yùndòng duì shēntǐ hěn hǎo, suǒyǐ wǒmen yīnggāi jīngcháng duànliàn.}\\
« Le sport est très bon pour le corps, c'est pourquoi nous devons nous entraîner régulièrement. »

\zh{爷爷在公园里\textbf{打太极拳}，奶奶喜欢\textbf{跳舞}。}\\
\emph{Yéye zài gōngyuán lǐ dǎ tàijíquán, nǎinai xǐhuan tiàowǔ.}\\
« Grand-père fait du tai-chi dans le parc, grand-mère aime danser. »
\end{exemplebox}

Vous remarquez déjà une chose essentielle : en chinois, un même mot peut parfois fonctionner comme nom ou comme verbe. \zh{运动} signifie « le sport » (nom) mais aussi « faire du sport » (verbe). Cette souplesse n'existe pas en français, et c'est l'un des points que nous allons clarifier.

\courssub{Tableau de vocabulaire : les mots clés}

Voici le glossaire complet du thème. Pour chaque mot, mémorisez quatre informations : les caractères, le pinyin avec les tons, la nature grammaticale et la traduction.

\begin{center}
\begin{tabularx}{\linewidth}{|X|X|X|X|}
\hline
\rowcolor{popblueL} Caractères & Pinyin & Nature & Traduction \\
\hline
运动 & yùndòng & n. / v. & sport ; faire du sport \\
\hline
锻炼 & duànliàn & v. & s'entraîner, faire de l'exercice \\
\hline
身体 & shēntǐ & n. & corps, santé (physique) \\
\hline
健康 & jiànkāng & adj. / n. & en bonne santé ; santé \\
\hline
跑步 & pǎobù & v. & courir (loisir, sport) \\
\hline
散步 & sànbù & v. & se promener \\
\hline
打太极拳 & dǎ tàijíquán & loc. v. & faire du tai-chi \\
\hline
跳舞 & tiàowǔ & v. & danser \\
\hline
游泳 & yóuyǒng & v. & nager \\
\hline
踢足球 & tī zúqiú & loc. v. & jouer au football \\
\hline
打篮球 & dǎ lánqiú & loc. v. & jouer au basket-ball \\
\hline
\end{tabularx}
\end{center}

\begin{definition}[运动 et 锻炼 : les deux piliers du thème]
\zh{运动} \emph{yùndòng} est un mot à double fonction : comme \cle{nom}, il désigne « le sport » en général (\zh{我喜欢运动。} « J'aime le sport. ») ; comme \cle{verbe}, il signifie « faire du sport » (\zh{你常常运动吗？} « Fais-tu souvent du sport ? »).\\
\zh{锻炼} \emph{duànliàn} est un \cle{verbe} qui signifie « s'entraîner, faire de l'exercice physique ». Il s'emploie souvent avec \zh{身体} : \zh{锻炼身体} « entraîner son corps, se maintenir en forme ». À l'origine, \zh{锻} et \zh{炼} désignent le travail du métal à la forge : s'entraîner, c'est « forger » son corps !
\end{definition}

\begin{propriete}[Le choix du verbe selon le sport : 打 ou 踢 ?]
En français, « jouer » convient à presque tous les sports ; en chinois, chaque sport exige son verbe propre, qu'il faut mémoriser avec le nom du sport :
\begin{itemize}
\item \zh{打} \emph{dǎ} « frapper, jouer » pour les sports de balle avec les mains et les arts martiaux : \zh{打篮球} « jouer au basket », \zh{打太极拳} « faire du tai-chi » ;
\item \zh{踢} \emph{tī} « frapper du pied » pour le football : \zh{踢足球} « jouer au football » (le caractère contient la clé du pied \zh{足}) ;
\item les sports sans ballon sont des verbes autonomes : \zh{跑步、游泳、跳舞、散步}.
\end{itemize}
On ne dit donc jamais \zh{玩足球} ni \zh{踢篮球} : le verbe est fixe.
\end{propriete}

\courssub{Tableau de vocabulaire : les mots complémentaires}

Ces mots ne désignent pas des sports, mais ils sont indispensables pour parler de la pratique sportive : ils expriment le temps, le lieu, le jugement et la conséquence.

\begin{center}
\begin{tabularx}{\linewidth}{|X|X|X|X|}
\hline
\rowcolor{popblueL} Caractères & Pinyin & Nature & Traduction \\
\hline
重要 & zhòngyào & adj. & important \\
\hline
快乐 & kuàilè & adj. & joyeux, heureux \\
\hline
应该 & yīnggāi & v. mod. & devoir (il faut) \\
\hline
所以 & suǒyǐ & conj. & c'est pourquoi, donc \\
\hline
每天 & měitiān & expr. t. & chaque jour \\
\hline
早上 & zǎoshang & n. t. & (le) matin \\
\hline
公园 & gōngyuán & n. & parc \\
\hline
对……好 & duì... hǎo & struct. & être bon pour... \\
\hline
\end{tabularx}
\end{center}

\begin{propriete}[La structure 对……好 / 对……不好]
La structure \cle{Sujet + 对 + personne/chose + 好 (ou 不好)} exprime « être bon (mauvais) pour... ». Le mot \zh{对} \emph{duì} introduit la cible de l'action ou du jugement, exactement comme « pour » en français, mais il se place \textbf{avant} l'adjectif, jamais après.
\begin{itemize}
\item \zh{运动对身体很好。} \emph{Yùndòng duì shēntǐ hěn hǎo.} « Le sport est très bon pour le corps. »
\item \zh{游泳对健康很重要。} \emph{Yóuyǒng duì jiànkāng hěn zhòngyào.} « La natation est très importante pour la santé. »
\item \zh{抽烟对身体不好。} \emph{Chōuyān duì shēntǐ bù hǎo.} « Fumer n'est pas bon pour le corps. »
\end{itemize}
\end{propriete}

\begin{attention}[L'ordre des mots : ne dites jamais « 运动很好对身体 »]
Les francophones ont tendance à calquer l'ordre français « le sport est très bon \emph{pour le corps} » et à placer le complément après l'adjectif. En chinois, c'est impossible : le complément introduit par \zh{对} se place \textbf{obligatoirement avant} l'adjectif. Retenez le schéma : \cle{Sujet + 对 + complément + adverbe + adjectif}. On dit \zh{运动对身体很好}, jamais \zh{运动很好对身体}.
\end{attention}

\courssub{Verbes, noms, adjectifs : bien distinguer les natures}

En chinois, il n'y a ni conjugaison ni article : c'est la \textbf{position} du mot dans la phrase qui révèle sa nature. Classons notre glossaire.

\begin{center}
\begin{tabularx}{\linewidth}{|X|X|X|}
\hline
\rowcolor{popgreenL} Nature & Mots & Emploi typique \\
\hline
Verbes & 锻炼、跑步、跳舞、游泳、散步、踢足球、打篮球 & après le sujet : \zh{我每天锻炼。} \\
\hline
Noms & 身体、健康(n.)、公园、运动(n.) & sujet ou complément : \zh{公园很大。} \\
\hline
Adjectifs & 健康(adj.)、重要、快乐 & avec \zh{很} : \zh{他很快乐。} \\
\hline
Expressions de temps & 每天、早上 & après le sujet, avant le verbe : \zh{我早上跑步。} \\
\hline
\end{tabularx}
\end{center}

\begin{propriete}[Deux règles d'emploi essentielles]
\textbf{1. L'adjectif chinois exige un adverbe.} En français on dit « il est heureux » avec le verbe « être » ; en chinois, l'adjectif fait office de verbe et s'emploie presque toujours avec \zh{很} \emph{hěn} (qui ne signifie plus vraiment « très » mais sert de simple lien) : \zh{他很快乐。} « Il est heureux. » Dire \zh{他快乐。} seul sonne incomplet ou comparatif. À la négative, \zh{很} disparaît : \zh{他不快乐。} « Il n'est pas heureux. »\\
\textbf{2. Le complément de temps se place avant le verbe.} \zh{每天} et \zh{早上} se placent après le sujet et avant le verbe : \zh{我每天早上锻炼。} « Je m'entraîne chaque matin. » Ne placez jamais le temps à la fin de la phrase comme en français.
\end{propriete}

\begin{exemplebox}[Chaque nature en action]
\zh{我每天锻炼。} \emph{Wǒ měitiān duànliàn.} « Je m'entraîne chaque jour. » (verbe)\\
\zh{运动对身体很好。} \emph{Yùndòng duì shēntǐ hěn hǎo.} « Le sport est très bon pour le corps. » (nom + nom)\\
\zh{奶奶的身体很健康。} \emph{Nǎinai de shēntǐ hěn jiànkāng.} « La santé de grand-mère est très bonne. » (adjectif)\\
\zh{踢足球很快乐。} \emph{Tī zúqiú hěn kuàilè.} « Jouer au football, c'est joyeux. » (locution verbale en sujet)\\
\zh{早上，公园里有很多人打太极拳。} \emph{Zǎoshang, gōngyuán lǐ yǒu hěn duō rén dǎ tàijíquán.} « Le matin, il y a beaucoup de gens qui font du tai-chi dans le parc. »
\end{exemplebox}

\courssub{Carte mentale du lexique}

Pour mémoriser durablement, organisons le vocabulaire en réseau autour du mot central \zh{运动}.

\begin{popfigure}
\begin{tikzpicture}[mindmap, grow cyclic, every node/.style={concept, align=center, font=\small}, concept color=popblue!40, level 1/.append style={level distance=3.2cm, sibling angle=72}, level 2/.append style={level distance=2.2cm, sibling angle=50}]
\node[concept color=poporange!50, font=\small\bfseries, align=center]{运动\\yùndòng\\sport}
  child[concept color=popgreen!40] { node[align=center]{球类 sports\\de balle} child { node[align=center]{踢足球\\football} } child { node[align=center]{打篮球\\basket} } }
  child[concept color=popgreen!40] { node[align=center]{个人\\individuel} child { node[align=center]{跑步\\courir} } child { node[align=center]{游泳\\nager} } child { node[align=center]{跳舞\\danser} } }
  child[concept color=poppurple!40] { node[align=center]{传统\\tradition} child { node[align=center]{打太极拳\\tai-chi} } }
  child[concept color=popgold!40] { node[align=center]{好处\\bienfaits} child { node[align=center]{健康\\santé} } child { node[align=center]{快乐\\joie} } child { node[align=center]{身体好\\corps} } }
  child[concept color=poppink!40] { node[align=center]{时间地点\\temps, lieu} child { node[align=center]{每天\\chaque jour} } child { node[align=center]{早上\\matin} } child { node[align=center]{公园\\parc} } };
\end{tikzpicture}
\legende{Carte mentale du vocabulaire du sport : cinq familles de mots autour de 运动.}
\end{popfigure}

\courssub{Entraînement à la prononciation : les tons difficiles}

Trois mots de ce glossaire posent des difficultés de tons aux francophones. Travaillez-les en répétant après le professeur, d'abord syllabe par syllabe, puis le mot entier.

\begin{methode}[Comment travailler les tons d'un mot nouveau]
\begin{enumerate}
\item Identifiez le ton de chaque syllabe dans le pinyin (regardez les marques).
\item Prononcez chaque syllabe seule, lentement, en exagérant le ton.
\item Enchaînez les deux syllabes sans les séparer, comme un seul mot.
\item Placez le mot dans une phrase courte et lisez la phrase entière.
\end{enumerate}
\end{methode}

\begin{exemplebox}[Trois mots à travailler]
\textbf{锻炼 duànliàn (4\up{e} + 4\up{e} ton)} : deux tons descendants à la suite. Le deuxième 4\up{e} ton part un peu moins haut que le premier. Évitez de monter sur la deuxième syllabe : dites \emph{duànliàn}, pas \emph{duànlián}.\\
\textbf{健康 jiànkāng (4\up{e} + 1\up{er} ton)} : un ton descendant suivi d'un ton haut et plat, bien tenu. Le contraste est fort : descendez sur \emph{jiàn}, puis restez haut et stable sur \emph{kāng}.\\
\textbf{跑步 pǎobù (3\up{e} + 4\up{e} ton)} : le 3\up{e} ton descend puis remonte (ou reste bas devant un autre ton), le 4\up{e} ton chute nettement. Ne prononcez pas \emph{páobù} : la première syllabe n'est pas un ton montant !
\end{exemplebox}

\begin{attention}[Ne confondez pas 跑步 et 散步]
\zh{跑步} \emph{pǎobù} signifie « courir » (sport, effort), tandis que \zh{散步} \emph{sànbù} signifie « se promener, faire une promenade ». Le premier caractère fait toute la différence : \zh{跑} contient la clé du pied \zh{足} et signifie « courir » ; \zh{散} évoque la dispersion, la détente, d'où « se promener sans but précis ». Attention aussi aux tons : \emph{pǎo} (3\up{e} ton) contre \emph{sàn} (4\up{e} ton). Dire \zh{我早上散步} signifie « je me promène le matin », pas « je cours le matin » !
\end{attention}

\begin{savaistu}[Le sport du matin en Chine]
En Chine, la gymnastique collective du matin, \zh{早操} \emph{zǎocāo}, est pratiquée dans la plupart des écoles : avant les cours, tous les élèves se réunissent dans la cour et font ensemble des exercices en musique. C'est un moment de discipline mais aussi de convivialité. De même, le tai-chi \zh{太极拳} \emph{tàijíquán}, art martial aux mouvements lents et harmonieux, est pratiqué chaque matin par des millions de personnes, souvent âgées, dans les parcs des villes chinoises. Au Cameroun, c'est plutôt le football qui rassemble : le dimanche, les terrains de quartier de Yaoundé ou de Douala ne désemplissent pas. Dans les deux pays, le sport est une affaire de santé autant que de plaisir !
\end{savaistu}

\courssub{Exercice résolu et entraînement}

\begin{exoresolu}[Compléter avec le bon mot]
Énoncé — Complétez chaque phrase avec le mot qui convient : \zh{锻炼、跑步、健康、重要、公园、所以}.\\
1. \zh{我每天早上在\_\_\_\_里跑步。}\\
2. \zh{运动对身体很\_\_\_\_。}\\
3. \zh{天气很好，\_\_\_\_我们去游泳吧。}\\
4. \zh{爷爷的身体很\_\_\_\_。}\\
5. \zh{我们应该每天\_\_\_\_身体。}\\
6. \zh{他不喜欢打篮球，他喜欢\_\_\_\_。}
\tcblower
Solution détaillée —\\
1. \zh{公园} : le mot \zh{里} (« à l'intérieur de ») appelle un lieu ; on court « dans le parc ». \emph{Wǒ měitiān zǎoshang zài gōngyuán lǐ pǎobù.}\\
2. \zh{重要} : structure \zh{对……很重要} « très important pour ». \emph{Yùndòng duì shēntǐ hěn zhòngyào.}\\
3. \zh{所以} : la deuxième partie est une conséquence (« il fait beau, \emph{c'est pourquoi} allons nager »). \emph{Tiānqì hěn hǎo, suǒyǐ wǒmen qù yóuyǒng ba.}\\
4. \zh{健康} : adjectif après \zh{很}, « le corps de grand-père est en bonne santé ». \emph{Yéye de shēntǐ hěn jiànkāng.}\\
5. \zh{锻炼} : verbe suivi du complément \zh{身体}, « nous devons entraîner notre corps chaque jour ». \emph{Wǒmen yīnggāi měitiān duànliàn shēntǐ.}\\
6. \zh{跑步} : il faut un sport ; parmi les mots restants, seul \zh{跑步} convient après \zh{喜欢}. \emph{Tā bù xǐhuan dǎ lánqiú, tā xǐhuan pǎobù.}
\end{exoresolu}

\begin{exercice}[À vous de jouer]
1. Traduisez en chinois : a) « Je cours chaque matin dans le parc. » b) « La natation est très bonne pour la santé. » c) « Nous devons faire du sport, c'est important pour le corps. »\\
2. Classez ces mots selon leur nature : \zh{快乐、游泳、身体、应该、重要、公园}.\\
3. Lisez à voix haute en respectant les tons : \zh{锻炼、健康、跑步、太极拳、重要}.
\end{exercice}

\begin{corrige}[Exercice « À vous de jouer »]
1. a) \zh{我每天早上在公园里跑步。} \emph{Wǒ měitiān zǎoshang zài gōngyuán lǐ pǎobù.} b) \zh{游泳对健康很好。} \emph{Yóuyǒng duì jiànkāng hěn hǎo.} c) \zh{我们应该运动，这对身体很重要。} \emph{Wǒmen yīnggāi yùndòng, zhè duì shēntǐ hěn zhòngyào.}\\
2. Adjectifs : \zh{快乐、重要} ; verbe : \zh{游泳} ; noms : \zh{身体、公园} ; verbe modal : \zh{应该}.\\
3. Vérifiez : \emph{duànliàn} (4-4), \emph{jiànkāng} (4-1), \emph{pǎobù} (3-4), \emph{tàijíquán} (4-2-2), \emph{zhòngyào} (4-4).
\end{corrige}

\begin{aretenir}[L'essentiel du glossaire]
\begin{itemize}
\item \zh{运动} est à la fois nom (« le sport ») et verbe (« faire du sport ») ; \zh{锻炼} est le verbe « s'entraîner », souvent dans \zh{锻炼身体}.
\item Les sports avec ballon se construisent avec un verbe précis : \zh{踢足球} (on « frappe » le football du pied), \zh{打篮球} (on « joue » au basket), \zh{打太极拳} (on « pratique » le tai-chi).
\item La structure \cle{Sujet + 对 + complément + 很 + adjectif} exprime « être bon/important pour... » ; le complément se place avant l'adjectif.
\item Les expressions de temps (\zh{每天、早上}) se placent après le sujet et avant le verbe.
\item Surveillez les tons : \emph{duànliàn} (4-4), \emph{jiànkāng} (4-1), \emph{pǎobù} (3-4), et ne confondez pas \zh{跑步} « courir » avec \zh{散步} « se promener ».
\end{itemize}
\end{aretenir}

Maintenant que le vocabulaire est solide, nous pouvons passer à la section suivante : l'étude des clés (radicaux) et de l'ordre des traits de quelques caractères du thème, pour apprendre à les écrire correctement.

\coursec{Clés (radicaux) et ordre des traits}

Dans la section précédente, nous avons appris le vocabulaire du sport et de la santé : \zh{运动} \emph{yùndòng} « sport », \zh{健康} \emph{jiànkāng} « santé », \zh{跑步} \emph{pǎobù} « course à pied », \zh{游泳} \emph{yóuyǒng} « natation ». Mais comment ces caractères sont-ils construits ? Pourquoi \zh{跑} et \zh{跳} se ressemblent-ils à gauche ? Pourquoi \zh{泳} et \zh{洗} commencent-ils par les mêmes trois petits traits ? La réponse tient en un mot : les \cle{clés} (ou \cle{radicaux}). Comprendre les clés, c'est obtenir une véritable « clé » pour lire, deviner et mémoriser les caractères chinois.

\courssub{Qu'est-ce qu'une clé (radical) ?}

\begin{textebox}[Observation : des caractères qui se ressemblent]
跑步、跳高、踢球。\\
\emph{Pǎobù, tiàogāo, tī qiú.}\\
« Courir, sauter en hauteur, jouer au football. »\\[0.4em]
游泳、洗澡、流汗。\\
\emph{Yóuyǒng, xǐzǎo, liú hàn.}\\
« Nager, se laver, transpirer. »\\[0.4em]
他们很健康。\\
\emph{Tāmen hěn jiànkāng.}\\
« Ils sont en bonne santé. »
\end{textebox}

Observez bien : \zh{跑}, \zh{跳}, \zh{踢} contiennent tous la même partie à gauche : \zh{足}. Et \zh{泳}, \zh{洗}, \zh{汗} commencent tous par trois petites gouttes : \zh{氵}. Ce n'est pas un hasard !

\begin{definition}[La clé (radical)]
Une \cle{clé} (en chinois \zh{部首} \emph{bùshǒu}) est un élément graphique d'un caractère qui indique généralement son \textbf{champ de sens} (le domaine sémantique). Chaque caractère chinois est classé sous une clé dans les dictionnaires : la clé sert donc à deux choses :
\begin{enumerate}
\item \textbf{Comprendre} : elle donne un indice sur le sens du caractère (pied, eau, main, cœur, parole...).
\item \textbf{Chercher} : dans un dictionnaire chinois traditionnel, on retrouve un caractère grâce à sa clé, puis en comptant les traits restants.
\end{enumerate}
L'autre partie du caractère, appelée souvent \cle{partie phonétique}, donne un indice sur la \textbf{prononciation}.
\end{definition}

\begin{exemplebox}[Clé + phonétique : le caractère 泳]
\zh{泳} \emph{yǒng} « nager » se décompose ainsi :
\begin{itemize}
\item \zh{氵} (trois gouttes d'eau, variante de \zh{水} \emph{shuǐ} « eau ») = la \textbf{clé} : le sens est lié à l'eau ;
\item \zh{永} \emph{yǒng} « éternel » = la \textbf{phonétique} : elle donne la prononciation \emph{yǒng}.
\end{itemize}
Un caractère chinois est donc souvent un petit « rébus » : une partie pour le sens, une partie pour le son. Comparaison avec le français : c'est un peu comme si le mot « nager » contenait le radical « aqua- » qui annonce l'eau, et une terminaison qui indique la prononciation.
\end{exemplebox}

\begin{propriete}[La clé donne souvent le champ de sens]
La clé d'un caractère indique très souvent le domaine de sens auquel il appartient. Quelques familles utiles pour notre thème du sport :
\begin{itemize}
\item \zh{足} (\emph{zú}, « pied ») → les actions des jambes et des pieds : \zh{跑} \emph{pǎo} « courir », \zh{跳} \emph{tiào} « sauter », \zh{踢} \emph{tī} « donner un coup de pied » ;
\item \zh{氵} (\emph{shuǐ}, « eau ») → les actions liées à l'eau : \zh{泳} \emph{yǒng} « nager », \zh{洗} \emph{xǐ} « laver », \zh{汗} \emph{hàn} « sueur » ;
\item \zh{亻} (\emph{rén}, « homme ») → les personnes et leurs états : \zh{健} \emph{jiàn} « robuste », \zh{休} \emph{xiū} « se reposer », \zh{们} \emph{men} (marque du pluriel des personnes) ;
\item \zh{辶} (\emph{chuò}, « marche, mouvement ») → les déplacements : \zh{运} \emph{yùn} « transporter », \zh{还} \emph{hái} « retourner », \zh{进} \emph{jìn} « entrer » ;
\item \zh{力} (\emph{lì}, « force ») → la force et l'effort : \zh{动} \emph{dòng} « bouger », \zh{功} \emph{gōng} « mérite, travail ».
\end{itemize}
\end{propriete}

\begin{popfigure}
\begin{tikzpicture}[mindmap, grow cyclic, every node/.style={concept, align=center}, root concept/.append style={concept color=popblue, font=\large\bfseries}, level 1/.append style={level distance=3.4cm, sibling angle=72}, level 1 concept/.append style={concept color=poporangeL, font=\small}]
\node[root concept, align=center]{部首 bùshǒu\\les clés}
  child { node[align=center]{足 zú\\pied} }
  child { node[align=center]{氵 shuǐ\\eau} }
  child { node[align=center]{亻 rén\\homme} }
  child { node[align=center]{辶 chuò\\marche} }
  child { node[align=center]{力 lì\\force} };
\node[concept color=popgreenL, font=\small] at (4.6,1.6) {跑 pǎo 跳 tiào 踢 tī};
\node[concept color=popgreenL, font=\small] at (1.8,4.4) {泳 yǒng 洗 xǐ 汗 hàn};
\node[concept color=popgreenL, font=\small] at (-2.6,4.2) {健 jiàn 休 xiū 们 men};
\node[concept color=popgreenL, font=\small] at (-4.8,0.4) {运 yùn 还 hái 进 jìn};
\node[concept color=popgreenL, font=\small] at (-1.8,-4.0) {动 dòng 功 gōng};
\end{tikzpicture}
\legende{Carte mentale : cinq clés du thème et leurs familles de caractères.}
\end{popfigure}

\courssub{Les règles générales de l'ordre des traits}

Avant d'étudier chaque caractère, il faut connaître les lois d'écriture. En chinois, chaque trait s'écrit dans un ordre précis, fixe, non négociable : c'est ce qui rend l'écriture lisible, rapide et belle.

\begin{methode}[Les cinq règles d'or de l'ordre des traits]
\begin{enumerate}
\item \textbf{De haut en bas} : on écrit d'abord les traits du haut, puis ceux du bas. Exemple : \zh{三} \emph{sān} « trois » : la barre du haut, puis celle du milieu, puis celle du bas.
\item \textbf{De gauche à droite} : on écrit d'abord la partie gauche, puis la partie droite. Exemple : \zh{你} \emph{nǐ} : d'abord \zh{亻}, puis \zh{尔}.
\item \textbf{L'horizontal avant le vertical} : quand un trait horizontal croise un vertical, l'horizontal passe d'abord. Exemple : \zh{十} \emph{shí} « dix » : d'abord —, puis |.
\item \textbf{L'extérieur avant l'intérieur} : pour les caractères avec une enceinte, on trace d'abord le cadre extérieur, puis le contenu. Exemple : \zh{同} \emph{tóng}.
\item \textbf{On ferme en dernier} : le trait qui « ferme » une boîte (la barre du bas de \zh{口} quand elle est pleine) se trace à la fin. Exemple : \zh{国} \emph{guó} « pays » : le cadre, le contenu \zh{玉}, puis la fermeture du bas.
\end{enumerate}
\textbf{Astuce de mémorisation} : retenez la phrase « \emph{Haut → bas, gauche → droite, dehors → dedans, on ferme à la fin} ».
\end{methode}

\courssub{Étude détaillée des caractères du thème}

\textbf{1. Le caractère 运 \emph{yùn} (dans 运动 « sport ») — 7 traits}\par
\zh{运} se compose de la phonétique \zh{云} \emph{yún} « nuage » (qui donne le son) et de la clé \zh{辶} \emph{chuò} « marche, mouvement » (qui donne le sens : transporter, faire bouger). Ordre des traits :
\begin{enumerate}
\item \zh{云} d'abord : trait horizontal du haut (1), deuxième horizontal (2), puis le bas de \zh{云} : le \zh{厶} en deux traits (3, 4) ;
\item \zh{辶} ensuite : le point (5), le trait brisé (6), et la longue vague horizontale qui « porte » le caractère (7).
\end{enumerate}

\begin{attention}[La clé 辶 s'écrit en DERNIER]
Erreur très fréquente chez les débutants francophones : comme \zh{辶} se trouve visuellement à gauche et en bas, on croit qu'il faut l'écrire en premier (règle « gauche avant droite »). \textbf{Non !} Pour tous les caractères avec \zh{辶} — \zh{运} \emph{yùn}, \zh{还} \emph{hái}, \zh{进} \emph{jìn}, \zh{这} \emph{zhè}, \zh{道} \emph{dào} — on écrit \textbf{d'abord la partie intérieure (la phonétique), puis la clé 辶 en dernier}. C'est une exception à la règle « gauche avant droite » : la vague finale de \zh{辶} vient envelopper le caractère comme un petit bateau qui le transporte.
\end{attention}

\textbf{2. Le caractère 动 \emph{dòng} (dans 运动) — 6 traits}\par
Clé : \zh{力} \emph{lì} « force ». Logique : « bouger », c'est mettre de la force en mouvement. Ordre des traits : d'abord \zh{云} à gauche (traits 1 à 4 : deux horizontaux, puis \zh{厶}), puis \zh{力} à droite : le trait brisé (trait brisé) en crochet (5), puis le grand trait descendant vers la gauche (6). Règle appliquée : gauche avant droite.

\textbf{3. Le caractère 健 \emph{jiàn} (dans 健康) — 10 traits}\par
Clé : \zh{亻} \emph{rén}, variante de \zh{人} « homme » : la santé concerne la personne. Ordre des traits : d'abord \zh{亻} à gauche (2 traits : le petit trait oblique, puis le vertical), puis la partie droite \zh{建} (8 traits) : on écrit \zh{聿} d'abord (de haut en bas), et la clé \zh{廴} en dernier — même logique que \zh{辶} !

\textbf{4. Le caractère 康 \emph{kāng} (dans 健康) — 11 traits}\par
Clé : \zh{广} \emph{guǎng} « abri, large toit » (à l'origine, un bâtiment : la santé était liée à la demeure paisible). Ordre des traits : d'abord le point du haut (1), puis l'horizontal (2), puis le grand trait oblique vers la gauche (3) — c'est l'enceinte \zh{广} — ensuite l'intérieur \zh{隶} (traits 4 à 11), de haut en bas. Règle appliquée : l'extérieur avant l'intérieur.

\begin{exemplebox}[健康 jiànkāng : « en bonne santé »]
\zh{他很健康。}\\
\emph{Tā hěn jiànkāng.}\\
« Il est en bonne santé. »\\[0.3em]
\zh{运动对健康有好处。}\\
\emph{Yùndòng duì jiànkāng yǒu hǎochù.}\\
« Le sport est bon pour la santé. »
\end{exemplebox}

\textbf{5. Le caractère 跑 \emph{pǎo} « courir » — 12 traits}\par
Clé : \zh{足} \emph{zú} « pied ». Lien logique évident : \textbf{on court avec les pieds} ! La partie droite \zh{包} \emph{bāo} donne la prononciation. Ordre des traits : d'abord \zh{足} à gauche (7 traits : le petit \zh{口} en haut — vertical, horizontal-crochet, horizontal de fermeture — puis le vertical, l'horizontal, le vertical court et le trait final), puis \zh{包} à droite (5 traits : le petit trait oblique, le crochet enveloppant, puis \zh{己} à l'intérieur, fermé en dernier).

\begin{popfigure}
\begin{tikzpicture}[font=\large]
\node[draw=popblue, line width=1pt, minimum size=1.9cm, font=\Huge, fill=popblueL] (zu) at (0,0) {足};
\node[below=2pt of zu, font=\small, align=center] {clé : pied\\7 traits (1--7)};
\node[font=\Huge] at (1.6,0) {+};
\node[draw=poporange, line width=1pt, minimum size=1.9cm, font=\Huge, fill=poporangeL] (bao) at (3.2,0) {包};
\node[below=2pt of bao, font=\small, align=center] {phonétique : bāo\\5 traits (8--12)};
\node[font=\Huge] at (4.8,0) {=};
\node[draw=popgreen, line width=1.5pt, minimum size=1.9cm, font=\Huge, fill=popgreenL] (pao) at (6.4,0) {跑};
\node[below=2pt of pao, font=\small, align=center] {跑 pǎo\\« courir »};
\draw[-{Stealth}, very thick, popdark] (zu.east) -- (bao.west);
\draw[-{Stealth}, very thick, popdark] (bao.east) -- (pao.west);
\draw[-{Stealth}, thick, poppurple, bend left=25] (zu.north) to node[above, font=\small\itshape]{1. gauche d'abord} (pao.north west);
\draw[-{Stealth}, thick, poppink, bend right=25] (bao.south) to node[below, font=\small\itshape]{2. droite ensuite} (pao.south west);
\end{tikzpicture}
\legende{Décomposition de 跑 : on écrit d'abord la clé 足 (pied, traits 1 à 7), puis la phonétique 包 (traits 8 à 12).}
\end{popfigure}

\textbf{6. Le caractère 泳 \emph{yǒng} « nager » — 8 traits}\par
Clé : \zh{氵} \emph{shuǐ}, les « trois gouttes d'eau ». Lien logique : \textbf{on nage dans l'eau} ! Ordre des traits : les trois gouttes à gauche (traits 1, 2, 3 — de haut en bas), puis \zh{永} à droite (5 traits : le point, l'horizontale avec crochet, le trait vertical central, le trait oblique gauche, le trait oblique droit).

\begin{center}
\begin{tabularx}{\linewidth}{|X|X|X|X|}
\hline
\rowcolor{popblueL} Caractère & Clé (sens) & Nombre de traits & Ordre : point de vigilance \\
\hline
运 yùn & 辶 marche & 7 & 辶 en dernier, après 云 \\
\hline
动 dòng & 力 force & 6 & 云 à gauche, puis 力 \\
\hline
健 jiàn & 亻 homme & 10 & 亻, puis 聿, puis 廴 en dernier \\
\hline
康 kāng & 广 abri & 11 & l'enceinte 广, puis l'intérieur \\
\hline
跑 pǎo & 足 pied & 12 & 足 à gauche, puis 包 \\
\hline
泳 yǒng & 氵 eau & 8 & 氵 (3 gouttes), puis 永 \\
\hline
\end{tabularx}
\end{center}

\begin{exoresolu}[Deviner le sens grâce à la clé]
Énoncé : sans dictionnaire, devinez le domaine de sens des caractères suivants grâce à leur clé : \zh{跳}, \zh{踢}, \zh{休}, \zh{河}.
\tcblower
Solution détaillée :
\begin{enumerate}
\item \zh{跳} \emph{tiào} : clé \zh{足} (pied) → action du pied ou de la jambe. Sens : « sauter ». Exemple : \zh{跳高} \emph{tiàogāo} « saut en hauteur ».
\item \zh{踢} \emph{tī} : clé \zh{足} (pied) → action du pied. Sens : « donner un coup de pied ». Exemple : \zh{踢球} \emph{tī qiú} « jouer au ballon (au pied) ».
\item \zh{休} \emph{xiū} : clé \zh{亻} (homme) + \zh{木} (arbre) → un homme appuyé contre un arbre : « se reposer ». Exemple : \zh{休息} \emph{xiūxi} « se reposer ».
\item \zh{河} \emph{hé} : clé \zh{氵} (eau) → lié à l'eau. Sens : « rivière, fleuve ».
\end{enumerate}
Conclusion : la clé ne donne pas toujours le sens exact, mais elle donne presque toujours le \textbf{domaine} : c'est un formidable outil de devinette en lecture.
\end{exoresolu}

\begin{attention}[Pièges fréquents des francophones]
\begin{itemize}
\item Ne confondez pas \zh{氵} (clé de l'eau, trois traits) avec \zh{冫} (clé de la glace, deux traits, comme dans \zh{冷} \emph{lěng} « froid »). Trois gouttes = eau ; deux gouttes = froid/glace.
\item La clé \zh{亻} (homme debout, à gauche) et le caractère \zh{人} (homme, seul) ont le même sens mais des formes différentes selon leur place.
\item Attention aux homophones du thème : \zh{云} \emph{yún} « nuage » et \zh{运} \emph{yùn} « transporter » n'ont pas le même ton (2\textsuperscript{e} ton contre 4\textsuperscript{e} ton) : la clé \zh{辶} change aussi le sens.
\item Ne « dessinez » pas les caractères dans le désordre : un mauvais ordre des traits rend l'écriture lente et illisible, et vous sera pénalisé à l'examen.
\end{itemize}
\end{attention}

\begin{exercice}[À vous de jouer]
\begin{enumerate}
\item Indiquez la clé de chaque caractère et devinez son domaine de sens : \zh{喝}, \zh{吃}, \zh{打}, \zh{拍}.
\item Recopiez cinq fois \zh{跑} et \zh{泳} en respectant l'ordre des traits vu en classe.
\item Pourquoi la clé de \zh{运} s'écrit-elle en dernier ? Expliquez avec vos mots.
\end{enumerate}
\end{exercice}

\begin{corrige}[Exercice « À vous de jouer »]
\begin{enumerate}
\item \zh{喝} \emph{hē} « boire » et \zh{吃} \emph{chī} « manger » : clé \zh{口} (bouche) → actions de la bouche. \zh{打} \emph{dǎ} « frapper, jouer (au ballon avec les mains) » et \zh{拍} \emph{pāi} « frapper, applaudir » : clé \zh{扌} (main) → actions de la main.
\item Vérification en classe : \zh{跑} = 足 (7 traits) puis 包 (5 traits) ; \zh{泳} = 氵 (3 traits) puis 永 (5 traits).
\item Parce que \zh{辶} est une clé « enveloppante » de déplacement : on écrit d'abord la partie intérieure (\zh{云}), puis la vague de \zh{辶} vient la porter, comme un chemin sous les pas.
\end{enumerate}
\end{corrige}

\begin{aretenir}[L'essentiel de la section]
\begin{itemize}
\item Une \cle{clé} (部首) indique le champ de sens d'un caractère et permet de le chercher au dictionnaire.
\item Familles à connaître : 足 (pied) → 跑, 跳, 踢 ; 氵 (eau) → 泳, 洗 ; 亻 (homme) → 健, 休 ; 辶 (marche) → 运, 还, 进 ; 力 (force) → 动.
\item Ordre des traits : haut → bas, gauche → droite, horizontal avant vertical, extérieur avant intérieur, on ferme en dernier.
\item Exception capitale : la clé 辶 (et 廴) s'écrit \textbf{en dernier}, après la partie phonétique.
\end{itemize}
\end{aretenir}

\coursec{Compréhension détaillée du texte}

Dans la section précédente, nous avons appris à écrire les caractères clés du thème du sport, en respectant les clés (radicaux) et l'ordre des traits. Savoir écrire un caractère, c'est bien ; comprendre ce qu'il dit dans un texte, c'est encore mieux. Nous allons maintenant relire ensemble le texte support, phrase par phrase, vérifier notre compréhension par des questions, repérer les petits mots qui relient les idées, puis apprendre à résumer le texte en français et en chinois.

\courssub{Relecture du texte support, phrase par phrase}

\begin{textebox}[Le sport est important — 运动很重要]
运动很重要。Yùndòng hěn zhòngyào. — Le sport est très important.\\
它对我们的身体好。Tā duì wǒmen de shēntǐ hǎo. — Il est bon pour notre corps.\\
运动让我们更健康，也让我们更快乐。Yùndòng ràng wǒmen gèng jiànkāng, yě ràng wǒmen gèng kuàilè. — Le sport nous rend plus en forme, et il nous rend aussi plus heureux.\\
很多中国人早上在公园里锻炼。Hěn duō Zhōngguó rén zǎoshang zài gōngyuán li duànliàn. — Beaucoup de Chinois font de l'exercice le matin dans les parcs.\\
他们跑步、打太极拳，也跳舞。Tāmen pǎobù, dǎ tàijíquán, yě tiàowǔ. — Ils courent, font du tai-chi, et dansent aussi.\\
运动以后，我们学习得更好。Yùndòng yǐhòu, wǒmen xuéxí de gèng hǎo. — Après le sport, nous apprenons mieux.\\
所以，学生应该每天运动。Suǒyǐ, xuésheng yīnggāi měitiān yùndòng. — Donc, les élèves devraient faire du sport tous les jours.\\
大家都喜欢运动！Dàjiā dōu xǐhuan yùndòng! — Tout le monde aime le sport !
\end{textebox}

Relisons phrase par phrase, en observant comment chaque idée s'enchaîne :

\begin{enumerate}
\item \zh{运动很重要。} — Phrase d'ouverture : elle annonce le \cle{thème} et l'\cle{idée principale}. En chinois comme en français, un court texte argumentatif commence souvent par sa thèse.
\item \zh{它对我们的身体好。} — Premier argument : la \cle{santé}. Remarquez la structure \zh{对...好} (être bon pour...), détaillée plus bas. Le pronom \zh{它} (tā, « il/elle » pour les choses) reprend \zh{运动}.
\item \zh{运动让我们更健康，也让我们更快乐。} — Deuxième argument : le \cle{bonheur}. Le mot \zh{也} (aussi) relie les deux effets : plus en forme \emph{et} plus heureux. L'adverbe \zh{更} (gèng) signifie « plus, encore plus » devant un adjectif.
\item \zh{很多中国人早上在公园里锻炼。} — Un exemple concret, tiré de la vie chinoise : la gymnastique matinale dans les parcs. Observez l'ordre : sujet + temps (\zh{早上}) + lieu (\zh{在公园里}) + verbe.
\item \zh{他们跑步、打太极拳，也跳舞。} — Précision de l'exemple : la virgule énumérative chinoise \zh{、} sépare les activités.
\item \zh{运动以后，我们学习得更好。} — Troisième argument : le sport améliore les \cle{études}. Structure \zh{verbe + 以后} = « après avoir fait ».
\item \zh{所以，学生应该每天运动。} — La \cle{conclusion}, introduite par \zh{所以} (donc) : la recommandation aux élèves. Le verbe modal \zh{应该} (yīnggāi) exprime le devoir, la recommandation : « devrait ».
\item \zh{大家都喜欢运动！} — Phrase finale enthousiaste, qui boucle le texte sur une note positive. \zh{都} (dōu, « tous ») se place après le sujet et avant le verbe.
\end{enumerate}

\begin{aretenir}[La logique du texte en trois temps]
Le texte suit un plan classique : \textbf{thèse} (运动很重要) → \textbf{arguments} (santé, bonheur, études) → \textbf{conclusion} (所以，学生应该每天运动). Retrouver ce plan, c'est déjà comprendre le texte.
\end{aretenir}

\begin{popfigure}
\begin{tikzpicture}[
  node distance=0.9cm and 1.2cm,
  main/.style={rectangle, rounded corners, draw=popblue, fill=popblueL, thick, align=center, minimum width=4.2cm, minimum height=0.9cm, font=\small},
  arg/.style={rectangle, rounded corners, draw=popgreen, fill=popgreenL, thick, align=center, minimum width=3.2cm, minimum height=0.9cm, font=\small},
  conc/.style={rectangle, rounded corners, draw=poporange, fill=poporangeL, thick, align=center, minimum width=4.6cm, minimum height=0.9cm, font=\small},
  fleche/.style={-{Stealth[length=2.5mm]}, thick, popdark}
]
\node[main, align=center] (idee){Idée principale\\运动很重要};
\node[arg, below left=1.0cm and -0.6cm of idee, align=center] (a1){健康\\la santé};
\node[arg, below=1.0cm of idee, align=center] (a2){快乐\\le bonheur};
\node[arg, below right=1.0cm and -0.6cm of idee, align=center] (a3){学习更好\\mieux étudier};
\node[conc, below=3.6cm of idee, align=center] (ccl){Conclusion\\学生应该每天运动};
\draw[fleche] (idee) -- (a1);
\draw[fleche] (idee) -- (a2);
\draw[fleche] (idee) -- (a3);
\draw[fleche] (a1.south) |- ($(ccl.north)+(0,0.25)$) -- (ccl.north);
\draw[fleche] (a2) -- (ccl);
\draw[fleche] (a3.south) |- ($(ccl.north)+(0,0.25)$) -- (ccl.north);
\end{tikzpicture}
\legende{Organigramme de la structure du texte : de l'idée principale vers les arguments, puis la conclusion.}
\end{popfigure}

\courssub{Questions de compréhension en français}

Répondre à une question de compréhension est une compétence méthodique : on ne répond pas « de mémoire », on répond \emph{avec le texte}.

\begin{methode}[Répondre à une question de compréhension en trois étapes]
\begin{enumerate}
\item \textbf{Souligner le mot-clé de la question.} Exemple : dans « Que font les Chinois le matin dans les parcs ? », les mots-clés sont « Chinois / matin / parcs ».
\item \textbf{Retrouver ce mot-clé dans le texte.} Ici : \zh{很多中国人早上在公园里锻炼} — on repère \zh{中国人}, \zh{早上}, \zh{公园}.
\item \textbf{Recopier ou adapter la phrase complète du texte} qui contient la réponse, puis la traduire proprement en français. Une réponse en un seul mot est toujours insuffisante.
\end{enumerate}
\end{methode}

\begin{exoresolu}[Questions de compréhension sur le texte]
\textbf{Énoncé.} Répondez en français, par des phrases complètes :\\
1. Pourquoi le sport est-il important ?\\
2. Que font les Chinois le matin dans les parcs ?\\
3. Que devraient faire les élèves ?
\tcblower
\textbf{Solution détaillée.}\\
1. Mot-clé : « important » → \zh{运动很重要}. Le texte donne trois raisons : le sport est bon pour le corps (\zh{它对我们的身体好}), il nous rend plus en forme et plus heureux (\zh{更健康、更快乐}), et il nous aide à mieux apprendre (\zh{学习得更好}).\\
\emph{Réponse : Le sport est important parce qu'il est bon pour notre corps, qu'il nous rend plus en forme et plus heureux, et qu'il nous aide à mieux étudier.}\\
2. Mots-clés : « Chinois / matin / parcs » → \zh{很多中国人早上在公园里锻炼。他们跑步、打太极拳，也跳舞。}\\
\emph{Réponse : Le matin, beaucoup de Chinois font de l'exercice dans les parcs : ils courent, font du tai-chi et dansent.}\\
3. Mot-clé : « élèves » → \zh{所以，学生应该每天运动。}\\
\emph{Réponse : Les élèves devraient faire du sport tous les jours.}
\end{exoresolu}

\courssub{Questions en chinois simple}

Passons maintenant aux questions posées directement en chinois : c'est l'exercice le plus formateur, car il oblige à lire la question \emph{en chinois} et à y répondre \emph{en chinois}, en s'appuyant sur le texte.

\begin{exemplebox}[Deux questions, deux réponses à trouver]
\textbf{Question 1 :} 运动对身体好不好？Yùndòng duì shēntǐ hǎo bu hǎo? — Le sport est-il bon ou non pour le corps ?\\
\textbf{Réponse :} 运动对身体很好。Yùndòng duì shēntǐ hěn hǎo. — Le sport est très bon pour le corps.\\[0.3em]
\textbf{Question 2 :} 很多中国人早上在哪儿锻炼？Hěn duō Zhōngguó rén zǎoshang zài nǎr duànliàn? — Où beaucoup de Chinois font-ils de l'exercice le matin ?\\
\textbf{Réponse :} 他们早上在公园里锻炼。Tāmen zǎoshang zài gōngyuán li duànliàn. — Le matin, ils font de l'exercice dans les parcs.
\end{exemplebox}

Remarquez la forme de la question 1 : \zh{好不好} est une \cle{question affirmative-négative} (adjectif + 不 + adjectif). Pour répondre « oui », on reprend l'adjectif seul : \zh{很好} ; pour répondre « non », on dirait \zh{不好} (bù hǎo). Pour la question 2, le mot interrogatif \zh{哪儿} (où) se place à l'endroit exact où se trouvera la réponse dans la phrase : \zh{在公园里}. L'ordre des mots ne change pas entre question et réponse — c'est une grande différence avec le français, qui inverse souvent sujet et verbe.

\courssub{Repérage des connecteurs du texte}

Les \cle{connecteurs} sont les petits mots qui relient les idées. Les repérer, c'est voir le « squelette » du texte.

\begin{center}
\begin{tabularx}{\linewidth}{|X|X|X|X|}
\hline
\rowcolor{popblueL} Connecteur & Pinyin & Sens & Emploi dans le texte \\
\hline
所以 & suǒyǐ & donc, c'est pourquoi & 所以，学生应该每天运动。 \\
\hline
以后 & yǐhòu & après, après avoir fait & 运动以后，我们学习得更好。 \\
\hline
也 & yě & aussi, également & 也让我们更快乐。 \\
\hline
更 & gèng & plus, encore plus & 更健康、更快乐、更好 \\
\hline
\end{tabularx}
\end{center}

\begin{propriete}[La structure 对...好 : être bon pour]
\textbf{Structure :} A + 对 + B + 好 (ou 有好处).\\
\zh{运动对身体好。} Yùndòng duì shēntǐ hǎo. — Le sport est bon pour le corps.\\
La préposition \zh{对} (duì, « envers, pour ») introduit le complément et se place \textbf{avant} celui-ci, entre le sujet et l'adjectif. On peut renforcer : \zh{很好} (très bon pour), ou nier : \zh{对身体不好} (mauvais pour le corps).
\end{propriete}

\begin{exemplebox}[Verbe + 以后 : après avoir fait]
\zh{运动以后，我们学习得更好。} Yùndòng yǐhòu, wǒmen xuéxí de gèng hǎo. — Après le sport, nous apprenons mieux.\\
Autres exemples :\\
\zh{吃饭以后，我去散步。} Chīfàn yǐhòu, wǒ qù sànbù. — Après le repas, je vais me promener.\\
\zh{下课以后，我们踢足球。} Xiàkè yǐhòu, wǒmen tī zúqiú. — Après les cours, nous jouons au football.
\end{exemplebox}

\begin{attention}[对 se place AVANT le complément]
Les francophones traduisent mot à mot « le corps est bon grâce au sport » et produisent la phrase fautive \zh{身体好运动}. C'est impossible en chinois. Dites toujours : \zh{运动对身体好} (sujet + 对 + complément + 好). De même, \zh{也} se place \textbf{avant le verbe} (\zh{也喜欢}), jamais en fin de phrase comme le « aussi » français. Enfin, ne confondez pas \zh{以后} (après) et \zh{以前} (yǐqián, avant) : \zh{运动以前} = « avant le sport ».
\end{attention}

\begin{savaistu}[Sport et réussite scolaire en Chine]
En Chine, l'éducation physique fait partie du quotidien des élèves : beaucoup d'écoles commencent la journée par des exercices collectifs, et le tai-chi (\zh{太极拳}) est pratiqué par des millions de personnes, jeunes et âgées, dans les parcs dès le lever du soleil. On observe que les élèves qui pratiquent un sport régulièrement sont souvent plus concentrés en classe : c'est exactement l'idée de notre texte, \zh{运动以后，我们学习得更好}. Au Cameroun aussi, beaucoup de professeurs remarquent que les élèves sportifs sont plus attentifs. Le sport et les études ne sont pas des ennemis, mais des alliés !
\end{savaistu}

\courssub{Vrai ou faux : vérification rapide}

\begin{exercice}[Vrai ou faux ? Justifiez avec le texte]
Cochez V (vrai) ou F (faux), puis citez la phrase du texte qui justifie votre choix.\\
1. 运动让我们更快乐。Yùndòng ràng wǒmen gèng kuàilè. \hspace{1cm} V / F\\
2. 学生不应该运动。Xuésheng bù yīnggāi yùndòng. \hspace{1cm} V / F\\
3. 很多中国人晚上在公园里锻炼。Hěn duō Zhōngguó rén wǎnshang zài gōngyuán li duànliàn. \hspace{1cm} V / F
\end{exercice}

\begin{corrige}[Exercice Vrai ou faux]
1. \textbf{V} — Le texte dit : \zh{也让我们更快乐} (il nous rend aussi plus heureux).\\
2. \textbf{F} — Le texte dit le contraire : \zh{学生应该每天运动} (les élèves devraient faire du sport tous les jours). Attention au mot négatif \zh{不} qui inverse le sens !\\
3. \textbf{F} — Le texte dit \zh{早上} (le matin), pas \zh{晚上} (le soir). Piège classique : une seule information changée rend la phrase fausse.
\end{corrige}

\courssub{Reformulation : résumer le texte}

Résumer, c'est garder l'essentiel avec ses propres mots. Procédons en deux temps : d'abord en français, puis en chinois guidé.

\begin{methode}[Résumer un court texte argumentatif]
\begin{enumerate}
\item Repérer l'idée principale (souvent la première phrase).
\item Choisir les arguments essentiels (ici : santé, bonheur, études) et supprimer les exemples de détail.
\item Garder la conclusion. En chinois, réutiliser les structures du texte (\zh{对...好}, \zh{所以}) en raccourcissant les phrases.
\end{enumerate}
\end{methode}

\begin{exemplebox}[Résumé en deux phrases]
\textbf{En français :} Le sport est très important parce qu'il est bon pour la santé, qu'il rend heureux et qu'il aide à mieux étudier. C'est pourquoi les élèves devraient faire du sport tous les jours.\\[0.4em]
\textbf{En chinois guidé :}\\
\zh{运动对身体好，也让我们更快乐。} Yùndòng duì shēntǐ hǎo, yě ràng wǒmen gèng kuàilè. — Le sport est bon pour le corps, il nous rend aussi plus heureux.\\
\zh{所以，学生应该每天运动。} Suǒyǐ, xuésheng yīnggāi měitiān yùndòng. — Donc, les élèves devraient faire du sport tous les jours.
\end{exemplebox}

\begin{exoresolu}[À vous de résumer]
\textbf{Énoncé.} Complétez ce résumé en chinois avec les mots suivants : 所以 / 对 / 也 / 以后.\\
运动\_\_\_\_身体好，\_\_\_\_让我们更快乐。运动\_\_\_\_，我们学习得更好。\_\_\_\_，学生应该每天运动。
\tcblower
\textbf{Solution.} 运动\textbf{对}身体好，\textbf{也}让我们更快乐。运动\textbf{以后}，我们学习得更好。\textbf{所以}，学生应该每天运动。\\
Vérifiez le sens : le sport est bon \emph{pour} (对) le corps, il nous rend \emph{aussi} (也) plus heureux ; \emph{après} (以后) le sport, nous apprenons mieux ; \emph{donc} (所以), les élèves devraient faire du sport chaque jour.
\end{exoresolu}

\begin{aretenir}[Ce qu'il faut retenir de cette section]
\begin{itemize}
\item Un texte argumentatif simple se lit en trois temps : thèse → arguments → conclusion.
\item Pour répondre à une question : souligner le mot-clé, le retrouver dans le texte, recopier la phrase complète.
\item Les connecteurs \zh{所以} (donc), \zh{以后} (après), \zh{也} (aussi) et \zh{更} (plus) guident la lecture.
\item Structure clé : A + 对 + B + 好 = « A est bon pour B », avec \zh{对} toujours avant le complément.
\item Dans une question avec mot interrogatif (\zh{哪儿}, \zh{谁}, \zh{什么}), l'ordre des mots ne change pas : le mot interrogatif occupe la place de la réponse.
\end{itemize}
\end{aretenir}

\coursec{Collocations et exemples d'emploi}

Dans la section précédente, nous avons lu et compris en détail le texte sur l'importance de la pratique du sport. Nous y avons rencontré des expressions comme \zh{打篮球} et \zh{踢足球}. Une question se pose alors naturellement : pourquoi deux verbes différents, \zh{打} et \zh{踢}, pour traduire le même mot français « jouer » ? C'est toute la logique des \cle{collocations} chinoises : en chinois, un verbe ne se combine pas avec n'importe quel nom. Cette section vous apprend à associer correctement verbes et sports, puis à construire vos propres phrases pour parler de vos activités sportives.

\courssub{Observer avant de comprendre : le verbe change selon le sport}

Lisons d'abord ces quatre phrases, toutes traduites en français par « jouer » ou « faire » :

\begin{exemplebox}[Observation : quel verbe avec quel sport ?]
\zh{我喜欢打篮球。} \emph{Wǒ xǐhuan dǎ lánqiú.} « J'aime jouer au basket. »\\
\zh{他周末踢足球。} \emph{Tā zhōumò tī zúqiú.} « Le week-end, il joue au foot. »\\
\zh{爷爷每天早上打太极拳。} \emph{Yéye měi tiān zǎoshang dǎ tàijíquán.} « Grand-père fait du tai-chi tous les matins. »\\
\zh{我们下午去游泳。} \emph{Wǒmen xiàwǔ qù yóuyǒng.} « Nous allons nager cet après-midi. »
\end{exemplebox}

Observez bien : le basket et le tai-chi prennent \zh{打}, le football prend \zh{踢}, et la natation n'a pas de verbe séparé du tout : \zh{游泳} est déjà un verbe à elle seule. Ce n'est pas un hasard. Regardons les caractères eux-mêmes : \zh{打} contient la clé de la \cle{main} (\zh{扌}), et \zh{踢} contient la clé du \cle{pied} (\zh{⻊}). Le chinois choisit le verbe selon la \textbf{partie du corps} utilisée dans le sport !

\begin{propriete}[La règle d'or : 打 avec les mains, 踢 avec les pieds]
\begin{itemize}
\item \zh{打} \emph{dǎ} (frapper, avec la main) s'emploie avec les sports de ballon joués \textbf{avec les mains} ou avec une raquette : \zh{打篮球} (jouer au basket), \zh{打排球} (jouer au volley), \zh{打乒乓球} (jouer au ping-pong), \zh{打网球} (jouer au tennis), \zh{打羽毛球} (jouer au badminton). Il s'emploie aussi avec le tai-chi : \zh{打太极拳}.
\item \zh{踢} \emph{tī} (donner un coup de pied) s'emploie avec les sports joués \textbf{avec les pieds} : \zh{踢足球} (jouer au football).
\item Avec les activités pratiquées « seules », le verbe est \textbf{incorporé} dans le mot : on dit directement \zh{跑步} \emph{pǎobù} (courir), \zh{游泳} \emph{yóuyǒng} (nager), \zh{跳舞} \emph{tiàowǔ} (danser), sans ajouter de verbe devant.
\end{itemize}
Structures générales : \textbf{Sujet + 打/踢 + nom du sport} ou \textbf{Sujet + verbe-sport (跑步, 游泳, 跳舞)}.
\end{propriete}

\begin{attention}[Le piège numéro un des francophones : 打足球]
En français, « jouer » s'utilise partout : jouer au foot, jouer au basket, jouer au tennis. Le réflexe du francophone est donc de traduire « jouer » partout par \zh{打}. C'est une \textbf{erreur très fréquente à l'examen} : on ne dit JAMAIS \zh{打足球}. Le football se joue avec les pieds, donc obligatoirement \zh{踢足球}. De même, ne dites pas \zh{打游泳} : la natation n'a pas besoin de verbe, on dit simplement \zh{我去游泳} (je vais nager). Retenez : mains $\rightarrow$ \zh{打}, pieds $\rightarrow$ \zh{踢}, activité seule $\rightarrow$ verbe incorporé.
\end{attention}

\begin{popfigure}
\begin{tikzpicture}[
  box/.style={draw=popblue, thick, rounded corners, fill=popblueL, align=center, minimum width=3.2cm, minimum height=1cm},
  boxb/.style={draw=poporange, thick, rounded corners, fill=poporangeL, align=center, minimum width=3.2cm, minimum height=1cm},
  head/.style={font=\bfseries, align=center}
]
\node[head, text=popblue] at (-3,2.6) {打 dǎ (avec les mains)};
\node[box, align=center] at (-3,1.5){打篮球\\ dǎ lánqiú (basket)};
\node[box, align=center] at (-3,0.3){打排球\\ dǎ páiqiú (volley)};
\node[box, align=center] at (-3,-0.9){打太极拳\\ dǎ tàijíquán (tai-chi)};
\node[head, text=poporange] at (3,2.6) {踢 tī (avec les pieds)};
\node[boxb, align=center] at (3,1.5){踢足球\\ tī zúqiú (football)};
\node[head, text=popgreen] at (3,0.3) {Verbe incorporé};
\node[draw=popgreen, thick, rounded corners, fill=popgreenL, align=center, minimum width=3.2cm] at (3,-0.9) {跑步 / 游泳 / 跳舞\\ courir / nager / danser};
\end{tikzpicture}
\legende{Carte mentale : le choix du verbe dépend de la partie du corps utilisée.}
\end{popfigure}

\courssub{Tableau récapitulatif des collocations}

\begin{center}
\begin{tabularx}{\linewidth}{|X|X|X|X|}
\hline
\rowcolor{popblueL} Verbe & Sport & Exemple (caractères + pinyin) & Traduction \\
\hline
打 (mains) & 篮球 & 我打篮球。Wǒ dǎ lánqiú. & Je joue au basket. \\
\hline
打 (mains) & 排球 & 她们打排球。Tāmen dǎ páiqiú. & Elles jouent au volley. \\
\hline
打 (mains) & 太极拳 & 老人打太极拳。Lǎorén dǎ tàijíquán. & Les personnes âgées font du tai-chi. \\
\hline
踢 (pieds) & 足球 & 我们踢足球。Wǒmen tī zúqiú. & Nous jouons au foot. \\
\hline
(incorporé) & 跑步 & 他每天早上跑步。Tā měi tiān zǎoshang pǎobù. & Il court tous les matins. \\
\hline
(incorporé) & 游泳 & 她夏天游泳。Tā xiàtiān yóuyǒng. & Elle nage en été. \\
\hline
(incorporé) & 跳舞 & 学生们在跳舞。Xuéshengmen zài tiàowǔ. & Les élèves sont en train de danser. \\
\hline
\end{tabularx}
\end{center}

Comparaison avec le français : là où le français utilise un seul verbe (« jouer », « faire »), le chinois en utilise plusieurs, choisis selon le sens concret de l'action. C'est plus précis, et c'est une richesse de la langue : le verbe chinois « dessine » le geste sportif. Remarquez aussi la dernière ligne du tableau : \zh{在} devant le verbe marque l'action en cours, comme notre « être en train de ».

\courssub{Expressions utiles pour parler du sport et de la santé}

Au-delà des collocations verbe + sport, quatre expressions du texte sont indispensables pour construire vos propres phrases :

\begin{definition}[Quatre expressions à maîtriser]
\begin{itemize}
\item \zh{做运动} \emph{zuò yùndòng} : faire du sport (expression générale, tous sports confondus). Exemple : \zh{我们应该每天做运动。} \emph{Wǒmen yīnggāi měi tiān zuò yùndòng.} « Nous devrions faire du sport tous les jours. »
\item \zh{多运动} \emph{duō yùndòng} : faire plus de sport. \zh{多} (beaucoup, plus) placé devant un verbe exprime une recommandation : \zh{你要多运动。} \emph{Nǐ yào duō yùndòng.} « Tu dois faire plus de sport. »
\item \zh{对……很好} \emph{duì\ldots{} hěn hǎo} : être très bon pour\ldots{} Structure : \textbf{A + 对 + B + 很好}. Exemple : \zh{运动对身体很好。} \emph{Yùndòng duì shēntǐ hěn hǎo.} « Le sport est très bon pour le corps. »
\item \zh{越来越健康} \emph{yuè lái yuè jiànkāng} : de plus en plus en forme. Structure : \textbf{越来越 + adjectif}. Exemple : \zh{他越来越健康。} \emph{Tā yuè lái yuè jiànkāng.} « Il est de plus en plus en forme. »
\end{itemize}
\end{definition}

\begin{exemplebox}[Exemples traduits et commentés]
\zh{我周末和朋友踢足球。} \emph{Wǒ zhōumò hé péngyou tī zúqiú.} « Le week-end, je joue au foot avec mes amis. » — Remarquez l'ordre chinois : sujet + temps (\zh{周末}) + \zh{和} + personne + verbe. Le complément de temps se place AVANT le verbe, contrairement au français.\\
\zh{游泳对身体很好。} \emph{Yóuyǒng duì shēntǐ hěn hǎo.} « La natation est très bonne pour le corps. » — Ici \zh{游泳} est le sujet de la phrase : un verbe-sport peut occuper la place du sujet, sans article, contrairement au français « la natation ».\\
\zh{我每天早上去跑步。} \emph{Wǒ měi tiān zǎoshang qù pǎobù.} « Je vais courir tous les matins. » — \zh{去} + verbe exprime le déplacement vers l'activité.\\
\zh{多做运动，你会越来越健康。} \emph{Duō zuò yùndòng, nǐ huì yuè lái yuè jiànkāng.} « Fais plus de sport, tu seras de plus en plus en forme. »
\end{exemplebox}

\courssub{Construire ses propres phrases et dialoguer}

Pour parler de vos goûts sportifs, deux modèles suffisent :

\begin{itemize}
\item \textbf{Sujet + 喜欢 + verbe + sport} : \zh{我喜欢踢足球。} \emph{Wǒ xǐhuan tī zúqiú.} « J'aime jouer au foot. »
\item \textbf{Sujet + moment + verbe-sport} : \zh{我每天早上跑步。} \emph{Wǒ měi tiān zǎoshang pǎobù.} « Je cours tous les matins. »
\end{itemize}

Ces deux modèles permettent déjà un échange authentique. Voici le mini-dialogue modèle à mémoriser et à réutiliser à l'oral :

\begin{textebox}[Mini-dialogue modèle : parler de son sport préféré]
A：你喜欢什么运动？\emph{Nǐ xǐhuan shénme yùndòng?} — Quel sport aimes-tu ?\\
B：我喜欢打篮球。你呢？\emph{Wǒ xǐhuan dǎ lánqiú. Nǐ ne?} — J'aime jouer au basket. Et toi ?\\
A：我喜欢踢足球。足球对身体很好！\emph{Wǒ xǐhuan tī zúqiú. Zúqiú duì shēntǐ hěn hǎo!} — J'aime jouer au foot. Le foot est très bon pour la santé !\\
B：对！我们都应该多运动。\emph{Duì! Wǒmen dōu yīnggāi duō yùndòng.} — C'est vrai ! Nous devrions tous faire plus de sport.
\end{textebox}

Notez l'outil de conversation \zh{你呢？} \emph{Nǐ ne?} « Et toi ? » : il suffit d'ajouter \zh{呢} après le pronom pour retourner la question. C'est un réflexe très utile à l'oral comme à l'écrit.

\begin{exoresolu}[Choisir le bon verbe]
Énoncé : complétez chaque phrase avec \zh{打}, \zh{踢}, \zh{对} ou rien (verbe incorporé) : 1) 我们每天……篮球。2) 喀麦隆的学生喜欢……足球。3) 我妈妈早上……太极拳。4) 他下午去……（nager）。5) 游泳……身体很好。
\tcblower
Solution détaillée : 1) \zh{打} : le basket se joue avec les mains $\rightarrow$ \zh{我们每天打篮球。} \emph{Wǒmen měi tiān dǎ lánqiú.} 2) \zh{踢} : le football se joue avec les pieds $\rightarrow$ \zh{喀麦隆的学生喜欢踢足球。} \emph{Kāmàilóng de xuésheng xǐhuan tī zúqiú.} 3) \zh{打} : le tai-chi prend toujours \zh{打} $\rightarrow$ \zh{我妈妈早上打太极拳。} \emph{Wǒ māma zǎoshang dǎ tàijíquán.} 4) rien : \zh{游泳} est un verbe incorporé $\rightarrow$ \zh{他下午去游泳。} \emph{Tā xiàwǔ qù yóuyǒng.} 5) \zh{对} : structure \zh{对……很好} $\rightarrow$ \zh{游泳对身体很好。} \emph{Yóuyǒng duì shēntǐ hěn hǎo.} « La natation est très bonne pour le corps. »
\end{exoresolu}

\begin{exercice}[À vous de jouer]
Traduisez en chinois (caractères + pinyin) : 1) J'aime jouer au football. 2) Mon grand-père fait du tai-chi tous les matins. 3) Le sport est très bon pour le corps. 4) Elle court le week-end avec ses amis. 5) Tu dois faire plus de sport, tu seras de plus en plus en forme.
\end{exercice}

\begin{corrige}[Exercice « À vous de jouer »]
1) \zh{我喜欢踢足球。} \emph{Wǒ xǐhuan tī zúqiú.} 2) \zh{我爷爷每天早上打太极拳。} \emph{Wǒ yéye měi tiān zǎoshang dǎ tàijíquán.} 3) \zh{运动对身体很好。} \emph{Yùndòng duì shēntǐ hěn hǎo.} 4) \zh{她周末和朋友跑步。} \emph{Tā zhōumò hé péngyou pǎobù.} 5) \zh{你要多运动，你会越来越健康。} \emph{Nǐ yào duō yùndòng, nǐ huì yuè lái yuè jiànkāng.}
\end{corrige}

\begin{savaistu}[Le football, roi des deux côtés du monde]
Au Cameroun comme en Chine, le football est le sport le plus populaire : des stades de Yaoundé et de Douala aux terrains d'école de Pékin, on joue au foot partout. Pour briller à l'examen et dans vos rédactions, retenez cette phrase complète : \zh{足球是世界上最受欢迎的运动。} \emph{Zúqiú shì shìjiè shàng zuì shòu huānyíng de yùndòng.} « Le football est le sport le plus populaire du monde. » Elle combine le superlatif \zh{最} (le plus) et la structure \zh{受} + verbe (recevoir, « être l'objet de ») : une phrase riche qui montre votre niveau.
\end{savaistu}

\begin{aretenir}[L'essentiel de la section]
Mains $\rightarrow$ \zh{打} (\zh{打篮球}, \zh{打排球}, \zh{打太极拳}) ; pieds $\rightarrow$ \zh{踢} (\zh{踢足球}) ; activités seules $\rightarrow$ verbe incorporé (\zh{跑步}, \zh{游泳}, \zh{跳舞}). Expressions clés : \zh{做运动}, \zh{多运动}, \zh{对……很好}, \zh{越来越健康}. Le complément de temps se place avant le verbe : \zh{我周末踢足球。} Et surtout : jamais \zh{打足球} !
\end{aretenir}

\coursec{Production guidée et synthèse}

Dans la section précédente, nous avons appris à combiner le vocabulaire du sport en collocations naturelles : \zh{踢足球} (tī zúqiú, jouer au football), \zh{锻炼身体} (duànliàn shēntǐ, entretenir son corps), \zh{对健康有好处} (duì jiànkāng yǒu hǎochù, être bon pour la santé). Il est maintenant temps de passer de la \emph{réutilisation} à la \emph{production} : vous allez rédiger vos propres phrases en chinois, puis nous ferons la synthèse complète de la leçon.

\courssub{Production écrite guidée : « le sport et moi »}

\textbf{Observer d'abord : un modèle de production}\par

Avant d'écrire, observons un texte modèle rédigé par un élève de Première. Repérez comment chaque phrase suit un schéma simple et régulier.

\begin{exemplebox}[Modèle de production : « le sport et moi »]
\zh{我喜欢运动。我每个周末踢足球。运动对我的身体很好，也让我更快乐。所以我经常锻炼。}\\
\emph{Wǒ xǐhuan yùndòng. Wǒ měi ge zhōumò tī zúqiú. Yùndòng duì wǒ de shēntǐ hěn hǎo, yě ràng wǒ gèng kuàilè. Suǒyǐ wǒ jīngcháng duànliàn.}\\
« J'aime le sport. Chaque week-end, je joue au football. Le sport est très bon pour mon corps et me rend aussi plus heureux. C'est pourquoi je m'entraîne souvent. »
\end{exemplebox}

Observez la construction : quatre phrases courtes, chacune bâtie sur la trame \cle{Sujet + temps + verbe + complément}. Aucune phrase ne dépasse une dizaine de caractères : en chinois débutant, la clarté prime toujours sur la longueur. Notez aussi l'emploi de \zh{让} (ràng, faire que, rendre) dans \zh{让我更快乐} : « me rend plus heureux » ; c'est une structure fréquente et facile à réutiliser.

\begin{methode}[Rédiger en chinois : la méthode en cinq étapes]
\begin{enumerate}
\item \textbf{Choisir la trame.} Pour ce sujet, utilisez les quatre amorces : \zh{我喜欢…} (j'aime…), \zh{我每天/周末…} (chaque jour / le week-end…), \zh{运动对身体…} (le sport pour le corps…), \zh{所以…} (c'est pourquoi…).
\item \textbf{Respecter l'ordre des mots.} Sujet + complément de temps + verbe + complément d'objet. Le temps se place \emph{avant} le verbe.
\item \textbf{Réutiliser le texte support comme modèle.} Reprenez les phrases de la leçon et ne changez qu'un élément à la fois (le sport, le moment, la personne).
\item \textbf{Vérifier les tons et les caractères.} Relisez en prononçant à voix haute : un ton faux peut changer le sens.
\item \textbf{Conclure par} \zh{所以}\textbf{.} La dernière phrase tire la conséquence logique des précédentes.
\end{enumerate}
\end{methode}

\begin{attention}[Le complément de temps se place AVANT le verbe]
En français on dit « je cours \emph{chaque jour} » : le complément de temps est à la fin. En chinois, il se place \textbf{avant le verbe}, juste après le sujet : \zh{我每天跑步。} (Wǒ měitiān pǎobù. — Je cours chaque jour.) Ne dites jamais \zh{我跑步每天}, qui est incorrect. Même règle pour \zh{周末} : \zh{我周末打篮球。} (Wǒ zhōumò dǎ lánqiú. — Le week-end, je joue au basket.)
\end{attention}

\begin{center}
\begin{tabularx}{\linewidth}{|X|X|X|}\hline
\rowcolor{popblueL} Trame & Exemple (caractères + pinyin) & Traduction \\ \hline
\zh{我喜欢} + sport & \zh{我喜欢游泳。} Wǒ xǐhuan yóuyǒng. & J'aime la natation. \\ \hline
\zh{我每天} + verbe & \zh{我每天跑步。} Wǒ měitiān pǎobù. & Je cours chaque jour. \\ \hline
\zh{我周末} + verbe & \zh{我周末踢足球。} Wǒ zhōumò tī zúqiú. & Le week-end, je joue au foot. \\ \hline
\zh{运动对身体} + adjectif & \zh{运动对身体很好。} Yùndòng duì shēntǐ hěn hǎo. & Le sport est très bon pour le corps. \\ \hline
\zh{所以} + conclusion & \zh{所以我经常锻炼。} Suǒyǐ wǒ jīngcháng duànliàn. & C'est pourquoi je m'entraîne souvent. \\ \hline
\end{tabularx}
\end{center}

\begin{attention}[Quel verbe pour quel sport ?]
Le chinois choisit le verbe selon la partie du corps dominante : \zh{踢} (tī, frapper du pied) pour le football, \zh{打} (dǎ, frapper de la main) pour le basket-ball, le volley-ball et le tennis de table, \zh{游} (yóu, nager) pour la natation (\zh{游泳}), \zh{跑} (pǎo, courir) pour la course (\zh{跑步}). On ne dit donc jamais \zh{打足球} ni \zh{踢篮球}.
\end{attention}

\begin{exoresolu}[Production guidée : compléter la trame]
Énoncé : complétez les quatre phrases avec le vocabulaire de la leçon. 1) 我喜欢\underline{\hspace{2cm}}。 2) 我每天\underline{\hspace{2cm}}。 3) 运动对身体\underline{\hspace{2cm}}。 4) 所以我\underline{\hspace{2cm}}。
\tcblower
Solution détaillée (une réponse possible ; d'autres sont correctes) :
1) \zh{我喜欢打篮球。} Wǒ xǐhuan dǎ lánqiú. — J'aime jouer au basket. Le verbe \zh{打} s'emploie avec les sports de ballon joués avec les mains.
2) \zh{我每天跑步。} Wǒ měitiān pǎobù. — Je cours chaque jour. Attention : \zh{每天} est placé avant le verbe \zh{跑步}.
3) \zh{运动对身体很好。} Yùndòng duì shēntǐ hěn hǎo. — Le sport est très bon pour le corps. Structure : \zh{对} + nom + \zh{很} + adjectif.
4) \zh{所以我经常锻炼。} Suǒyǐ wǒ jīngcháng duànliàn. — C'est pourquoi je m'entraîne souvent. \zh{所以} introduit la conséquence, comme en français.
\end{exoresolu}

\courssub{Correction collective d'un exemple au tableau}

Voici la production d'un élève, avec ses erreurs. Cherchons-les ensemble avant de lire la correction.

Phrase de l'élève : \zh{我跑步每天，运动很好对身体。}

\begin{corrige}[Correction collective]
Deux erreurs d'ordre des mots, typiques des francophones :
\begin{itemize}
\item \zh{我跑步每天} : le complément de temps \zh{每天} doit précéder le verbe. Correction : \zh{我每天跑步。} (Wǒ měitiān pǎobù.)
\item \zh{运动很好对身体} : la préposition \zh{对} et son nom forment un bloc qui se place \emph{avant} l'adjectif. Correction : \zh{运动对身体很好。} (Yùndòng duì shēntǐ hěn hǎo.)
\end{itemize}
Phrase corrigée complète : \zh{我每天跑步，运动对身体很好。} — « Je cours chaque jour ; le sport est très bon pour le corps. » Vérification finale : tons de \zh{跑} (pǎo, 3\ieme{} ton) et de \zh{好} (hǎo, 3\ieme{} ton) bien marqués à l'oral ; à la suite, deux 3\ieme{} tons se prononcent avec le premier transformé en 2\ieme{} ton (règle d'euphonie), mais le pinyin écrit garde le ton d'origine.
\end{corrige}

\courssub{Synthèse de la leçon}

\begin{aretenir}[Ce qu'il faut retenir de la leçon 9]
\begin{itemize}
\item \textbf{Vocabulaire} : \zh{运动} (yùndòng, sport), \zh{锻炼} (duànliàn, s'entraîner), \zh{身体} (shēntǐ, corps), \zh{健康} (jiànkāng, santé), \zh{重要} (zhòngyào, important), \zh{踢足球} (tī zúqiú), \zh{打篮球} (dǎ lánqiú), \zh{跑步} (pǎobù), \zh{游泳} (yóuyǒng).
\item \textbf{Radicaux} : \zh{身} est la clé du corps (dans \zh{身体}) ; \zh{足} (pied) est la clé de \zh{跑} et \zh{踢} ; \zh{氵} (eau) est la clé de \zh{游} et \zh{泳}.
\item \textbf{Idées du texte} : le sport renforce le corps, aide à mieux travailler à l'école et rend plus heureux ; il faut s'entraîner régulièrement.
\item \textbf{Structures clés} : Sujet + temps + verbe + complément ; \zh{对} + nom + \zh{很} + adjectif ; \zh{所以} pour conclure.
\end{itemize}
\end{aretenir}

\begin{textebox}[Phrase-clé à retenir pour l'examen]
\zh{运动很重要，我们应该每天锻炼。}\\
\emph{Yùndòng hěn zhòngyào, wǒmen yīnggāi měitiān duànliàn.}\\
« Le sport est très important, nous devrions nous entraîner chaque jour. »
\end{textebox}

\begin{popfigure}
\begin{tikzpicture}[
  etape/.style={rectangle, rounded corners=4pt, draw=popblue, fill=popblueL, align=center, minimum width=2.5cm, minimum height=1cm, font=\small},
  fleche/.style={-{Stealth[length=3mm]}, thick, poporange}
]
\node[etape, align=center] (t) at (0,0){\zh{文本}\\Texte};
\node[etape, right=0.6cm of t, align=center] (v){\zh{词汇}\\Vocabulaire};
\node[etape, right=0.6cm of v, align=center] (r){\zh{部首}\\Radicaux};
\node[etape, right=0.6cm of r, align=center] (c){\zh{理解}\\Compréhension};
\node[etape, right=0.6cm of c, fill=poporangeL, draw=poporange, align=center] (p){\zh{写作}\\Production};
\draw[fleche] (t) -- (v);
\draw[fleche] (v) -- (r);
\draw[fleche] (r) -- (c);
\draw[fleche] (c) -- (p);
\end{tikzpicture}
\legende{Frise de synthèse : le parcours de la leçon, du texte support jusqu'à votre propre production écrite.}
\end{popfigure}

\courssub{Évaluation formative : quiz de cinq questions}

\begin{exercice}[Quiz rapide (10 minutes, sans document)]
\begin{enumerate}
\item Traduisez en français : \zh{我每天锻炼。} (Wǒ měitiān duànliàn.)
\item Traduisez en chinois : « Le sport est bon pour la santé. »
\item Quelle est la clé (radical) du caractère \zh{跑} ? Que signifie cette clé ?
\item D'après le texte de la leçon, citez deux bienfaits du sport.
\item Corrigez la phrase : \zh{我踢足球周末。}
\end{enumerate}
\end{exercice}

\begin{corrige}[Quiz]
\begin{enumerate}
\item « Je m'entraîne chaque jour. »
\item \zh{运动对健康很好。} (Yùndòng duì jiànkāng hěn hǎo.) — structure \zh{对} + nom + \zh{很} + adjectif. On accepte aussi \zh{运动对健康有好处。} (Yùndòng duì jiànkāng yǒu hǎochù.)
\item La clé est \zh{足} (zú), le pied : elle indique que le caractère a un rapport avec le pied ou la course.
\item Deux réponses possibles parmi : le sport rend le corps fort (\zh{身体好}), il aide à mieux étudier, il rend plus heureux (\zh{更快乐}).
\item Le complément de temps \zh{周末} doit précéder le verbe : \zh{我周末踢足球。} (Wǒ zhōumò tī zúqiú.)
\end{enumerate}
\end{corrige}

\begin{experience}[Pour aller plus loin : mini-exposé oral]
Par binôme : chaque élève lit sa production « le sport et moi » à son camarade, qui doit poser une question de compréhension (\zh{你喜欢什么运动？} Nǐ xǐhuan shénme yùndòng ? — Quel sport aimes-tu ?). Puis deux volontaires présentent leur texte à la classe. Critères d'évaluation : exactitude des tons, ordre des mots, utilisation d'au moins trois mots du glossaire de la leçon.
\end{experience}

\newpage

\coursec{Activité d'intégration}

\coursec{Leçon 9 — Vocabulaire et compréhension de texte : importance de la pratique du sport}

\courssub{Mise en route et découverte du texte}

\begin{textebox}[Situation de communication : la Semaine du sport au lycée]
Tu es élève en Première A au Lycée bilingue de Yaoundé. Le club « Sport et Santé » organise la \emph{Semaine du sport}. Le proviseur a constaté que beaucoup d'élèves restent assis pendant la récréation. Il demande à chaque classe de chinois de réaliser une \cle{affiche en chinois} pour encourager la pratique sportive. Les affiches seront exposées près du terrain de football ; les élèves de l'école primaire voisine viendront les admirer.

Voici le modèle d'affiche affiché au tableau d'information, qui servira de texte support à la leçon :
\end{textebox}

\begin{textebox}[Affiche du club « Sport et Santé » — Texte support]
运动很重要！\\
Yùndòng hěn zhòngyào !\\
« Le sport est très important ! »\\[4pt]
每天运动三十分钟，身体好，学习也好。\\
Měitiān yùndòng sānshí fēnzhōng, shēntǐ hǎo, xuéxí yě hǎo.\\
« Trente minutes de sport par jour : bonne santé et bonnes études. »\\[4pt]
跑步、游泳、打球、踢足球……你选一个吧！\\
Pǎobù, yóuyǒng, dǎ qiú, tī zúqiú…… nǐ xuǎn yí ge ba !\\
« Courir, nager, jouer au ballon, jouer au football… choisis-en un ! »\\[4pt]
为了健康，我们一起运动吧！\\
Wèile jiànkāng, wǒmen yìqǐ yùndòng ba !\\
« Pour la santé, faisons du sport ensemble ! »
\end{textebox}

\begin{exercice}[Compréhension globale du texte support]
Répondez en français aux questions suivantes :
\begin{enumerate}
\item Quel est le message principal de cette affiche ?
\item Quelle durée d'exercice quotidien est recommandée ?
\item Citez les quatre activités sportives mentionnées.
\item Quelle structure grammaticale exprime le but dans la dernière phrase ?
\item Quelle particule finale marque la suggestion ou l'invitation ?
\end{enumerate}
\end{exercice}

\begin{corrige}[Exercice : Compréhension globale]
\begin{enumerate}
\item L'importance du sport pour la santé et la réussite scolaire ; invitation à pratiquer une activité.
\item Trente minutes par jour (\zh{三十分钟}).
\item \zh{跑步} (courir), \zh{游泳} (nager), \zh{打球} (jouer au ballon), \zh{踢足球} (jouer au football).
\item La structure \zh{为了} + nom + ，+ phrase (\zh{为了健康，我们一起运动吧！}).
\item La particule \zh{吧} (ba), qui adoucit l'impératif et invite à l'action.
\end{enumerate}
\end{corrige}

\courssub{Glossaire du thème : sport et santé}

\begin{definition}[Vocabulaire de la leçon (HSK 2--3)]
Le tableau ci-dessous regroupe le lexique essentiel du texte support et de l'activité d'intégration. La nature grammaticale est indiquée selon l'usage dans la leçon (n = nom, v = verbe, adj = adjectif, adv = adverbe, prep = préposition, part = particule, num = numéral, mw = classificateur).
\end{definition}

\begin{center}
\begin{tabularx}{\linewidth}{|X|X|X|X|}\hline
\rowcolor{popblueL} \textbf{Caractères} & \textbf{Pinyin} & \textbf{Nature} & \textbf{Traduction} \\ \hline
运动 & yùndòng & n / v & sport ; faire du sport, s'exercer \\ \hline
健康 & jiànkāng & adj / n & en bonne santé ; santé \\ \hline
身体 & shēntǐ & n & corps, condition physique \\ \hline
跑步 & pǎobù & v & courir, faire du jogging \\ \hline
游泳 & yóuyǒng & v & nager \\ \hline
踢足球 & tī zúqiú & v + n & jouer au football (litt. « donner un coup de pied au ballon de pied ») \\ \hline
打篮球 & dǎ lánqiú & v + n & jouer au basket-ball (litt. « frapper le ballon panier ») \\ \hline
打球 & dǎ qiú & v + n & jouer au ballon (général) \\ \hline
每天 & měitiān & adv & chaque jour, tous les jours \\ \hline
三十分钟 & sānshí fēnzhōng & num + mw & trente minutes \\ \hline
重要 & zhòngyào & adj & important \\ \hline
一起 & yìqǐ & adv & ensemble \\ \hline
为了 & wèile & prep & pour, dans le but de (structure de but) \\ \hline
选择 & xuǎnzé & v & choisir, sélectionner \\ \hline
学校 & xuéxiào & n & école \\ \hline
体育场 & tǐyùchǎng & n & stade \\ \hline
\end{tabularx}
\end{center}

\courssub{Clés (radicaux) et ordre des traits}

\begin{propriete}[Radicaux clés du vocabulaire sportif]
La connaissance des radicaux (clés) permet de mémoriser l'écriture, de deviner le sens et de classer les caractères dans le dictionnaire.
\end{propriete}

\begin{center}
\begin{tabularx}{\linewidth}{|X|X|X|X|}\hline
\rowcolor{popgreenL} \textbf{Radical (clé)} & \textbf{Nom / Sens} & \textbf{Caractères de la leçon} & \textbf{Explication} \\ \hline
\zh{辶} & \emph{chuò} / clé du mouvement & \zh{运} (dans \zh{运动}), \zh{过}, \zh{这} & Indique le déplacement, l'action de marcher/aller. S'écrit \textbf{en dernier}. \\ \hline
\zh{氵} & \emph{shuǐ} / clé de l'eau (3 traits) & \zh{游} (dans \zh{游泳}), \zh{泳}, \zh{汗} & Indique un lien avec l'eau, le liquide. S'écrit \textbf{en premier} (haut → bas). \\ \hline
\zh{⻊} (\zh{足}) & \emph{zú} / clé du pied & \zh{跑}, \zh{踢}, \zh{跳} & Indique une action du pied : courir, donner un coup de pied, sauter. \\ \hline
\zh{亻} & \emph{rén} / clé de l'homme (2 traits) & \zh{体} (dans \zh{身体}), \zh{健} (dans \zh{健康}), \zh{他} & Indique un lien avec l'être humain, le corps, l'action humaine. \\ \hline
\end{tabularx}
\end{center}

\begin{popfigure}
\begin{tikzpicture}[scale=0.8, every node/.style={font=\small}]
% Caractère 运
\node[align=center] at (0,2) {\textbf{Ordre des traits : \zh{运} (yùn) -- 7 traits}};
\node[draw, thick, minimum width=2cm, minimum height=2cm] (yun) at (0,0) {};
\node[align=center] at (0,0) {\zh{云}\\\textcolor{popblue}{1 à 4}\\\textcolor{poporange}{\zh{辶} 5 à 7}};
% Annotations
\draw[popblue, thick, ->] (-1.5,1) -- (-2.5,1.5) node[left, align=left] {1. \zh{一} horiz.\\(云)};
\draw[popblue, thick, ->] (-1.5,0.5) -- (-2.5,0.5) node[left, align=left] {2. \zh{乛} pli\\(云)};
\draw[popblue, thick, ->] (-1.5,-0.5) -- (-2.5,-0.5) node[left, align=left] {3. \zh{乚} crochet\\(云)};
\draw[popblue, thick, ->] (-1.5,-1) -- (-2.5,-1) node[left, align=left] {4. \zh{丶} point\\(cloud)};
\draw[poporange, thick, ->] (1.5,-0.5) -- (2.5,-0.5) node[right, align=left] {5. \zh{丶} point\\(辶)};
\draw[poporange, thick, ->] (1.5,-1) -- (2.5,-1) node[right, align=left] {6. \zh{乛} pli\\(辶)};
\draw[poporange, thick, ->] (1.5,-1.5) -- (2.5,-1.5) node[right, align=left] {7. \zh{乙} crochet long\\(辶)};

% Caractère 游
\node[align=center] at (6,2) {\textbf{Ordre des traits : \zh{游} (yóu) -- 12 traits}};
\node[draw, thick, minimum width=2cm, minimum height=2cm] (you) at (6,0) {};
\node[align=center] at (6,0) {\zh{氵}\\\textcolor{popblue}{1--3}\\\zh{方}\\\textcolor{popgreen}{4--7}\\\zh{子}\\\textcolor{poppurple}{8--12}};
\draw[popblue, thick, ->] (4.5,1) -- (3.5,1.5) node[left, align=left] {1--3. \zh{氵} (eau)\\points, crochet, point};
\draw[popgreen, thick, ->] (4.5,0) -- (3.5,0) node[left, align=left] {4--7. \zh{方} (carré)\\horiz, vert, pli, horiz};
\draw[poppurple, thick, ->] (4.5,-1) -- (3.5,-1) node[left, align=left] {8--12. \zh{子} (enfant)\\horiz, crochet, horiz, vert, horiz};
\end{tikzpicture}
\legende{Analyse de la composition et ordre des traits pour \zh{运} (clé \zh{辶} en dernier) et \zh{游} (clé \zh{氵} en premier, puis phonétique).}
\end{popfigure}

\begin{attention}[Pièges d'écriture fréquents]
\begin{itemize}
\item \zh{运} : N'écrivez \textbf{pas} la clé \zh{辶} en premier ! L'intérieur \zh{云} se fait complètement (traits 1 à 4), puis \zh{辶} (traits 5 à 7).
\item \zh{游} : La clé \zh{氵} (trois points d'eau) s'écrit en premier (haut → bas : point, point, crochet montant). Le reste (\zh{方} + \zh{子}) suit.
\item \zh{跑} et \zh{踢} : La clé du pied \zh{⻊} (simplifiée de \zh{足}) se place en bas à gauche. Elle compte 7 traits. Attention à ne pas confondre avec \zh{身} (corps) dans \zh{身体}.
\item \zh{健} : Clé \zh{亻} (2 traits : trait horizontal court, trait vertical long) à gauche. La partie droite \zh{建} (construire) sert de phonétique.
\end{itemize}
\end{attention}

\courssub{Compréhension détaillée et structures grammaticales}

\begin{propriete}[Structures grammaticales observées dans le texte support]
Chaque structure est illustrée par une phrase du texte. Maîtrisez le schéma pour l'activité de production.
\end{propriete}

\begin{center}
\begin{tabularx}{\linewidth}{|X|X|X|}\hline
\rowcolor{popblueL} \textbf{Structure / Règle} & \textbf{Exemple du texte (caractères + pinyin)} & \textbf{Traduction / Note} \\ \hline
\textbf{Prédicat adjectival} : Sujet + \zh{很} + Adjectif & \zh{运动很重要。} \emph{Yùndòng hěn zhòngyào.} & « Le sport est très important. » \zh{很} lie le sujet à l'adjectif (pas de verbe \zh{是}). \\ \hline
\textbf{Complément de durée} : Sujet + Temps (quand) + Verbe + Durée & \zh{每天运动三十分钟。} \emph{Měitiān yùndòng sānshí fēnzhōng.} & « Faire du sport 30 min chaque jour. » La durée suit le verbe. \zh{每天} (chaque jour) est un adverbe de temps placé avant le verbe. \\ \hline
\textbf{Structure parallèle / énumération} : A 好，B 也好 & \zh{身体好，学习也好。} \emph{Shēntǐ hǎo, xuéxí yě hǎo.} & « Le corps va bien, les études aussi. » \zh{也} (aussi) marque la similarité. Structure rythmée, facile à mémoriser. \\ \hline
\textbf{Complément de but} : \zh{为了} + But (nom) + ，+ Sujet + Verbe & \zh{为了健康，我们一起运动吧！} \emph{Wèile jiànkāng, wǒmen yìqǐ yùndòng ba !} & « Pour la santé, faisons du sport ensemble ! » \zh{为了} introduit la finalité. La virgule sépare le but de l'action principale. \\ \hline
\textbf{Complément de lieu avant le verbe} : \zh{在} + Lieu + Verbe & (Adaptation) \zh{在学校的体育场跑步。} \emph{Zài xuéxiào de tǐyùchǎng pǎobù.} & « Courir au stade de l'école. » En chinois, le lieu (\zh{在} + lieu) se place \textbf{avant} le verbe. \\ \hline
\textbf{Particule de suggestion} : Verbe / Phrase + \zh{吧} & \zh{你选一个吧！} \emph{Nǐ xuǎn yí ge ba !} & « Choisis-en un ! » \zh{吧} adoucit l'ordre, exprime une proposition, un accord ou une estimation. \\ \hline
\textbf{Virgule d'énumération} : \zh{、} & \zh{跑步、游泳、打球、踢足球} & Remplace la virgule française pour lister des noms ou verbes de même nature. Pas d'espace après. \\ \hline
\end{tabularx}
\end{center}

\begin{exemplebox}[Collocations usuelles (搭配) du thème]
Ces associations « mot + mot » sont fixes et naturelles. Apprenez-les par blocs.
\begin{itemize}
\item \zh{做运动} (zuò yùndòng) — faire du sport (verbe léger \zh{做} + nom).
\item \zh{锻炼身体} (duànliàn shēntǐ) — exercer son corps, se muscler (verbe composé + objet).
\item \zh{保持健康} (bǎochí jiànkāng) — maintenir la santé.
\item \zh{去跑步 / 去游泳} (qù pǎobù / qù yóuyǒng) — aller courir / nager (\zh{去} + activité).
\item \zh{踢足球 / 打篮球 / 打乒乓球} (tī zúqiú / dǎ lánqiú / dǎ pīngpāngqiú) — sports à ballon : \zh{踢} (pied) pour foot, \zh{打} (main/raquette) pour basket, ping-pong, badminton (\zh{打羽毛球}).
\item \zh{每天坚持} (měitiān jiānchí) — persévérer chaque jour.
\end{itemize}
\end{exemplebox}

\courssub{Production guidée et synthèse : Création d'une affiche « Sport et Santé »}

\courssub{Objectifs de l'activité}

À l'issue de cette activité d'intégration (50 min), tu seras capable de :
\begin{itemize}
\item \cle{réutiliser} le vocabulaire du sport et de la santé (\zh{运动}, \zh{健康}, \zh{身体}, \zh{跑步}, \zh{游泳}, \zh{重要}…) dans une production écrite réelle ;
\item \cle{écrire correctement} les caractères du glossaire en respectant l'ordre des traits et les radicaux (\zh{辶}, \zh{氵}, \zh{⻊}, \zh{亻}) ;
\item \cle{comprendre et adapter} les structures du texte support (prédicat adjectival, but \zh{为了}, lieu \zh{在}...V, suggestion \zh{吧}) pour créer des phrases simples et correctes ;
\item \cle{présenter à l'oral} un document collectif en lisant les phrases avec une prononciation et des tons corrects ;
\item \cle{travailler en équipe} : répartir les tâches, vérifier le travail des camarades, défendre un choix devant la classe.
\end{itemize}

\begin{methode}[Déroulement de l'activité — 50 minutes]
\begin{enumerate}
\item \textbf{Relecture et analyse (5 min).} En groupe de 4, relisez le texte support (\emph{Affiche du club}). Soulignez 3 phrases réutilisables. Vérifiez la compréhension de chacun (sens, tons, structure).
\item \textbf{Choix du slogan (5 min).} Créez ou choisissez un slogan court et percutant (ex. \zh{运动很重要！} ou \zh{为了健康，动起来！}). Écrivez-le en grand en haut de l'affiche (format A3 ou A2), \textbf{pinyin au-dessus des caractères}.
\item \textbf{Sélection du vocabulaire (10 min).} Choisissez \textbf{au moins 5 mots} du glossaire. Chaque élève écrit 1--2 mots sur l'affiche : \textbf{Caractères + Pinyin + Traduction française} (petit). Vérification croisée de l'ordre des traits.
\item \textbf{Rédaction de 3 phrases (10 min).} Adaptez le texte support. Contraintes : (1) Une phrase avec \zh{很} + adj. (2) Une phrase avec lieu \zh{在}... avant le verbe. (3) Une phrase énumérant des sports avec \zh{、} et finissant par \zh{吧}. Relisez : ordre SUJET -- TEMPS -- LIEU -- VERBE -- DURÉE/OBJET, tons, ponctuation chinoise (\zh{。}，\zh{！}，\zh{、}).
\item \textbf{Illustration (5 min).} Dessin simple en lien avec le slogan (élève qui court, ballons, terrain du lycée, gourde d'eau). Couleurs vives.
\item \textbf{Présentation orale et vote (15 min).} Affichage au tableau. 1--2 élèves/groupe lisent le slogan et les 3 phrases \textbf{à voix haute, tons soignés}. Classe : écoute active, notation sur grille (Caractères / Prononciation / Vocabulaire). Vote. Correction collective par le prof.
\end{enumerate}
\end{methode}

\begin{attention}[Pièges à éviter dans votre affiche — RAPPEL]
\begin{itemize}
\item \textbf{Confusion \zh{运动} vs \zh{活动}} : \zh{运动} = sport/exercice physique ; \zh{活动} = activité (général). Ne les mélangez pas.
\item \textbf{Tons critiques} : \zh{跑步} = \emph{pǎobù} (3\textsuperscript{e} + 4\textsuperscript{e}), pas *pàobū. \zh{健康} = \emph{jiànkāng} (4\textsuperscript{e} + 1\textsuperscript{er}). \zh{重要} = \emph{zhòngyào} (4\textsuperscript{e} + 4\textsuperscript{e}).
\item \textbf{Ordre des traits \zh{运}} : \zh{云} (1--4) PUIS \zh{辶} (5--7). Jamais l'inverse.
\item \textbf{Ponctuation} : Pleine chasse \zh{。！，、} collée au caractère précédent. Pas d'espace. Pas de virgule française `,` ni point `.`.
\item \textbf{Radical \zh{⻊} (pied)} : Dans \zh{跑} et \zh{踢}, c'est la clé du pied. Dans \zh{身体}, \zh{体} a la clé \zh{亻} (homme), pas le pied !
\item \textbf{Classificateur} : \zh{一个} (yí ge) pour « un » (objet/choix). \zh{三十分钟} : \zh{分钟} est le classificateur de durée, pas de \zh{个} entre nombre et minute.
\end{itemize}
\end{attention}

\courssub{Ressources mobilisables pour l'activité}

\begin{definition}[Glossaire de référence rapide]
\zh{运动} \emph{yùndòng} (sport) | \zh{健康} \emph{jiànkāng} (santé) | \zh{身体} \emph{shēntǐ} (corps) | \zh{跑步} \emph{pǎobù} (courir) | \zh{游泳} \emph{yóuyǒng} (nager) | \zh{踢足球} \emph{tī zúqiú} (foot) | \zh{打篮球} \emph{dǎ lánqiú} (basket) | \zh{打球} \emph{dǎ qiú} (ballon) | \zh{每天} \emph{měitiān} (chaque jour) | \zh{重要} \emph{zhòngyào} (important) | \zh{一起} \emph{yìqǐ} (ensemble) | \zh{为了} \emph{wèile} (pour) | \zh{选择} \emph{xuǎnzé} (choisir) | \zh{学校} \emph{xuéxiào} (école) | \zh{体育场} \emph{tǐyùchǎng} (stade).
\end{definition}

\begin{propriete}[Tableau récapitulatif des structures modèles]
\begin{center}
\begin{tabularx}{\linewidth}{|X|X|X|}\hline
\rowcolor{popblueL} \textbf{Structure} & \textbf{Modèle (car. + pinyin)} & \textbf{Traduction} \\ \hline
Adj + \zh{很} + Adj & \zh{运动很重要。} \emph{Yùndòng hěn zhòngyào.} & Le sport est très important. \\ \hline
Temps + V + Durée & \zh{每天运动三十分钟。} \emph{Měitiān yùndòng sānshí fēnzhōng.} & 30 min de sport par jour. \\ \hline
\zh{在} + Lieu + V & \zh{在体育场跑步。} \emph{Zài tǐyùchǎng pǎobù.} & Courir au stade. \\ \hline
\zh{为了} + But + ，+ Phrase & \zh{为了健康，我们运动吧！} \emph{Wèile jiànkāng, wǒmen yùndòng ba !} & Pour la santé, bougeons ! \\ \hline
Énumération + \zh{吧} & \zh{跑步、游泳……选一个吧！} \emph{Pǎobù, yóuyǒng…… xuǎn yí ge ba !} & Course, natation… choisis-en un ! \\ \hline
\end{tabularx}
\end{center}
\end{propriete}

\begin{exoresolu}[Modèle d'affiche « Bougeons pour notre santé ! »]
Énoncé : Rédigez l'affiche complète de votre groupe (slogan, 5 mots de vocabulaire, 3 phrases), prête à être présentée.

\tcblower

\textbf{Solution détaillée — Affiche modèle :}

\textbf{Slogan (en grand, centré en haut) :}\\
\zh{为了健康，我们一起动起来！}\\
\emph{Wèile jiànkāng, wǒmen yìqǐ dòng qǐlái !}\\
« Pour la santé, mettons-nous en mouvement ensemble ! »\\
\emph{(Variante plus dynamique que le texte support : \zh{动起来} « se mettre en mouvement / bouger » = verbe composé directionnel).}

\textbf{Vocabulaire (au centre, 2 colonnes, écriture soignée) :}\\
\begin{tabular}{lll}
\zh{运动} \emph{yùndòng} — le sport & \zh{健康} \emph{jiànkāng} — la santé \\
\zh{跑步} \emph{pǎobù} — courir & \zh{游泳} \emph{yóuyǒng} — nager \\
\zh{踢足球} \emph{tī zúqiú} — jouer au football &
\end{tabular}

\textbf{Trois phrases (en bas, espacées) :}\\
1. \zh{运动很重要，身体好，学习也好。}\\
\emph{Yùndòng hěn zhòngyào, shēntǐ hǎo, xuéxí yě hǎo.}\\
« Le sport est très important : bonne santé et bonnes études. »\\[4pt]
2. \zh{每天在学校的体育场跑步三十分钟。}\\
\emph{Měitiān zài xuéxiào de tǐyùchǎng pǎobù sānshí fēnzhōng.}\\
« Chaque jour, courir trente minutes au stade de l'école. »\\[4pt]
3. \zh{游泳、踢足球、打篮球……你选一个吧！}\\
\emph{Yóuyǒng, tī zúqiú, dǎ lánqiú…… nǐ xuǎn yí ge ba !}\\
« Natation, football, basket… choisis-en un ! »

\textbf{Commentaire linguistique des choix :}
\begin{itemize}
\item \textbf{Slogan} : Reprend la structure de but \zh{为了} + Nom + Phrase. \zh{动起来} (dòng qǐlái) est un verbe composé de direction (HSK 3) signifiant « commencer à bouger / se mettre en mouvement », plus incitatif que \zh{运动} seul. La particule \zh{吧} (ou ici implicite dans l'impératif doux) invite à l'action collective.
\item \textbf{Vocabulaire} : Les 5 mots couvrent les 4 radicaux étudiés : \zh{辶} (\zh{运}), \zh{氵} (\zh{游}), \zh{⻊} (\zh{跑}, \zh{踢}), \zh{亻} (\zh{健}, \zh{体}). L'affiche devient un outil de révision orthographique.
\item \textbf{Phrase 1} : Adaptation directe du texte support. Structure parallèle \zh{A好，B也好} (A bien, B aussi bien) : rythme binaire, facile à retenir, valeur proverbe.
\item \textbf{Phrase 2} : Insertion du complément de lieu \zh{在学校的体育场} \textbf{avant} le verbe \zh{跑步}. Rappel de la règle d'or : \textbf{Lieu avant Verbe}. \zh{三十分钟} (durée) reste \textbf{après} le verbe.
\item \textbf{Phrase 3} : Énumération avec la virgule chinoise \zh{、}. Remplace \zh{打球} (générique) par \zh{打篮球} (spécifique, connu au Cameroun). Terminaison par \zh{吧} pour l'interpellation directe du lecteur.
\end{itemize}
\end{exoresolu}

\courssub{Bilan de l'activité}

Grâce à cette tâche authentique, tu as mobilisé l'ensemble des savoir-faire de la leçon :

\begin{itemize}
\item \textbf{Compréhension de l'écrit} : Relire un texte modèle, en extraire le sens global et les structures réutilisables.
\item \textbf{Lexique} : Sélectionner, écrire (caractères + pinyin + tons) et prononcer au moins 5 mots du thème « sport/santé ».
\item \textbf{Écriture des caractères} : Appliquer l'ordre des traits standard et identifier les radicaux \zh{辶}, \zh{氵}, \zh{⻊}, \zh{亻} comme indices de sens et de classification.
\item \textbf{Grammaire} : Réemployer correctement : prédicat adjectival (\zh{很}), complément de durée, structure de but (\zh{为了}), place du lieu (\zh{在}...V), particule de suggestion (\zh{吧}), virgule d'énumération (\zh{、}).
\item \textbf{Expression orale} : Lecture à voix haute en groupe, attention aux tons (3\textsuperscript{e} ton \zh{跑}, 4\textsuperscript{e} ton \zh{步}, \zh{重}, \zh{要}...), écoute évaluative des pairs.
\item \textbf{Compétence socioculturelle} : Contexte camerounais réel (lycée bilingue, terrain de foot, semaine du sport), valeur éducative partagée Cameroun/Chine (sport = santé + réussite).
\end{itemize}

\begin{aretenir}[L'essentiel de la Leçon 9]
Une communication efficace en chinois sur le sport repose sur trois piliers :
\begin{enumerate}
\item Un \cle{slogan court} structuré par la finalité : \zh{为了} + But + \zh{，} + Action + \zh{吧！}
\item Un \cle{vocabulaire précis} écrit avec des caractères corrects : radicaux (\zh{辶} mouvement, \zh{氵} eau, \zh{⻊} pied, \zh{亻} homme) + ordre des traits rigoureux.
\item Des \cle{phrases types} adaptées d'un modèle : \zh{很} + adj., \zh{在} Lieu + V, Durée post-verbale, \zh{、} pour lister, \zh{吧} pour suggérer.
\end{enumerate}
Dire et écrire \zh{运动很重要！} (\emph{Yùndòng hěn zhòngyào !}) avec des caractères justes et des tons exacts, c'est déjà savoir convaincre en chinois.
\end{aretenir}

\begin{savaistu}[Regard croisé Cameroun -- Chine : Le sport à l'école]
\begin{itemize}
\item \textbf{Cameroun} : L'EPS (Éducation Physique et Sportive) est obligatoire au secondaire. Le football est roi (\zh{足球} \emph{zúqiú} = « ballon de pied »). Les interclasses et la FENASSCO (Fédération Nationale des Sports Scolaires) rythment l'année. Le stade Ahmadou Ahidjo (Yaoundé) et Japoma (Douala) sont les temples du sport.
\item \textbf{Chine} : L'éducation physique est une matière majeure (\zh{体育课} \emph{tǐyùkè}). Le test national \zh{体测} \emph{tǐcè} (évaluation de la condition physique) compte pour le \zh{高考} \emph{gāokǎo} (bac). Sports traditionnels : \zh{太极拳} \emph{tàijíquán} (tai-chi), \zh{武术} \emph{wǔshù} (arts martiaux), \zh{羽毛球} \emph{yǔmáoqiú} (badminton), \zh{乒乓球} \emph{pīngpāngqiú} (tennis de table = sport national). Le dicton \zh{文明其精神，野蛮其体魄} \emph{(Wénmíng qí jīngshén, yěmán qí tǐpò)} — « Civiliser l'esprit, fortifier le corps » (Mao Zedong, 1917) — guide encore la politique éducative.
\item \textbf{Point commun} : Dans les deux pays, le sport scolaire est vu comme levier de santé publique, de cohésion sociale et de réussite académique (\zh{身体好，学习也好}).
\end{itemize}
\end{savaistu}

\newpage

\coursec{Méthodes et exercices résolus}

\courssub{Mise en route et découverte du texte}

Au Cameroun, le football est roi : dans les quartiers de Yaoundé et de Douala, on joue dès la fin des classes. En Chine, les élèves commencent souvent la journée par des exercices collectifs dans la cour de l'école. Dans les deux pays, le sport fait partie de la vie des lycéens. Cette leçon te donne le vocabulaire et les outils pour lire, comprendre et produire un texte chinois sur l'importance de la pratique du sport.

\begin{textebox}[运动很重要 — Le sport est très important]
李明是雅温得一所高中的学生。他每天早上六点起床，然后去跑步。他说：« 运动对身体很好。每天跑步让我更健康，也让我上课更有精神。» 周末，他常常和朋友们一起打篮球或者游泳。他觉得运动不但对身体好，而且对学习也有帮助。\\
Lǐ Míng shì Yǎwēndé yì suǒ gāozhōng de xuésheng. Tā měi tiān zǎoshang liù diǎn qǐchuáng, ránhòu qù pǎobù. Tā shuō : « Yùndòng duì shēntǐ hěn hǎo. Měi tiān pǎobù ràng wǒ gèng jiànkāng, yě ràng wǒ shàngkè gèng yǒu jīngshen. » Zhōumò, tā chángcháng hé péngyoumen yìqǐ dǎ lánqiú huòzhě yóuyǒng. Tā juéde yùndòng bùdàn duì shēntǐ hǎo, érqiě duì xuéxí yě yǒu bāngzhù.\\
Li Ming est élève dans un lycée de Yaoundé. Chaque matin, il se lève à six heures, puis va courir. Il dit : « Le sport est très bon pour le corps. Courir chaque jour me rend plus en forme et me donne plus d'énergie en classe. » Le week-end, il joue souvent au basket ou nage avec ses amis. Il trouve que le sport est non seulement bon pour le corps, mais aussi utile pour les études.
\end{textebox}

\courssub{Glossaire du thème : sport et santé}

\begin{definition}[Vocabulaire de la leçon]
Voici le lexique à connaître pour cette leçon. Apprends chaque mot avec son pinyin, ses tons et sa nature grammaticale.
\end{definition}

\begin{center}
\begin{tabularx}{\linewidth}{|X|X|X|X|}
\hline
\rowcolor{popblueL} Caractères & Pinyin & Nature & Traduction \\
\hline
运动 & yùndòng & nom / verbe & sport ; faire du sport \\
\hline
跑步 & pǎobù & verbe & courir (footing) \\
\hline
游泳 & yóuyǒng & verbe & nager \\
\hline
打篮球 & dǎ lánqiú & verbe + objet & jouer au basket-ball \\
\hline
踢足球 & tī zúqiú & verbe + objet & jouer au football \\
\hline
身体 & shēntǐ & nom & corps, santé physique \\
\hline
健康 & jiànkāng & adjectif / nom & en bonne santé ; santé \\
\hline
锻炼 & duànliàn & verbe & s'entraîner, faire de l'exercice \\
\hline
休息 & xiūxi & verbe & se reposer \\
\hline
精神 & jīngshen & nom / adjectif & énergie, forme ; en forme \\
\hline
起床 & qǐchuáng & verbe & se lever (du lit) \\
\hline
周末 & zhōumò & nom & week-end \\
\hline
常常 & chángcháng & adverbe & souvent \\
\hline
有时候 & yǒu shíhou & locution adverbiale & parfois \\
\hline
或者 & huòzhě & conjonction & ou (dans une affirmation) \\
\hline
帮助 & bāngzhù & nom / verbe & aide ; aider \\
\hline
觉得 & juéde & verbe & trouver, penser (opinion) \\
\hline
应该 & yīnggāi & verbe modal & devoir (il faut) \\
\hline
体育场 & tǐyùchǎng & nom & stade \\
\hline
\end{tabularx}
\end{center}

\begin{attention}[Faux amis et pièges de vocabulaire]
\begin{itemize}
\item \zh{或者} (huòzhě) s'emploie dans les phrases \emph{affirmatives} : \zh{打篮球或者游泳}. Dans une question offrant un choix, on utilise \zh{还是} (háishi) : \zh{你打篮球还是游泳？} « Tu joues au basket ou tu nages ? »
\item \zh{觉得} exprime une opinion personnelle (« je trouve que »), pas un fait : \zh{我觉得运动很重要。}
\item \zh{身体} désigne le corps et, par extension, la santé : \zh{身体好} = « être en bonne santé ». Ne pas traduire mot à mot « le corps est bon ».
\end{itemize}
\end{attention}

\courssub{Méthode 1 — Lire et comprendre un court texte chinois sur le sport}

\begin{methode}[Comprendre un texte sur la pratique du sport]
Pour lire efficacement un texte chinois sur le sport et la santé, suis ces étapes :
\begin{enumerate}
\item \textbf{Lecture globale silencieuse} : lis tout le texte une première fois sans dictionnaire. Repère le thème général (ici : pourquoi faire du sport ?) et les mots que tu connais déjà (\zh{运动}, \zh{身体}, \zh{健康}).
\item \textbf{Repérage des mots-clés du thème} : souligne le vocabulaire du sport (\zh{跑步}, \zh{游泳}, \zh{打篮球}) et de la santé (\zh{锻炼}, \zh{休息}, \zh{精神}). Ces mots portent l'essentiel du sens.
\item \textbf{Découpage des phrases} : en chinois, chaque phrase suit le schéma Sujet + (complément de temps/lieu) + Verbe + Complément. Identifie d'abord le sujet, puis le verbe.
\item \textbf{Lecture phrase par phrase avec le pinyin} : prononce à voix haute en respectant les tons ; la prononciation aide à mémoriser et à vérifier le sens.
\item \textbf{Vérification avec la traduction} : compare ta compréhension avec la version française, puis relis le chinois une dernière fois sans aide.
\item \textbf{Résumé oral} : sois capable de dire en deux phrases chinoises de quoi parle le texte (\zh{这篇文章说……} Zhè piān wénzhāng shuō… « Cet article dit que… »).
\end{enumerate}
Réflexe d'examen : les questions de compréhension suivent presque toujours l'ordre du texte. La réponse à la question 1 se trouve au début, la dernière vers la fin.
\end{methode}

\begin{exoresolu}[Compréhension de texte : le sport au lycée]
\textbf{Énoncé.} Relis le texte \zh{运动很重要} (ci-dessus), puis réponds aux questions en chinois, par des phrases complètes.
\begin{enumerate}
\item 李明几点起床？(Lǐ Míng jǐ diǎn qǐchuáng ? — Li Ming se lève à quelle heure ?)
\item 他早上做什么运动？(Tā zǎoshang zuò shénme yùndòng ? — Quel sport fait-il le matin ?)
\item 周末他常常和谁一起运动？(Zhōumò tā chángcháng hé shéi yìqǐ yùndòng ? — Avec qui fait-il souvent du sport le week-end ?)
\end{enumerate}
\tcblower
\textbf{Solution détaillée.}
\begin{enumerate}
\item \textbf{Question 1.} Le mot interrogatif \zh{几点} demande l'heure. On cherche dans le texte un nombre suivi de \zh{点} : on trouve \zh{每天早上六点起床}. Réponse rédigée : \zh{他每天早上六点起床。} (Tā měi tiān zǎoshang liù diǎn qǐchuáng.) « Il se lève chaque matin à six heures. »
\item \textbf{Question 2.} \zh{什么运动} demande le type de sport. Le texte dit \zh{然后去跑步}. Réponse : \zh{他早上去跑步。} (Tā zǎoshang qù pǎobù.) « Le matin, il va courir. »
\item \textbf{Question 3.} \zh{和谁} demande la personne accompagnante. Le texte dit \zh{和朋友们一起打篮球或者游泳}. Réponse : \zh{周末他常常和朋友们一起运动。} (Zhōumò tā chángcháng hé péngyoumen yìqǐ yùndòng.) « Le week-end, il fait souvent du sport avec ses amis. »
\end{enumerate}
\textbf{Justification grammaticale} : dans les réponses, on garde l'ordre Sujet + complément de temps (\zh{周末}, \zh{早上}) + verbe. Le complément de temps se place avant le verbe, jamais à la fin comme en français.
\textbf{Conclusion} : chaque réponse reprend la formulation du texte en la complétant en phrase entière avec sujet explicite (\zh{他}) — c'est ce qu'attend l'examinateur.
\end{exoresolu}

\courssub{Méthode 2 — Reconnaître et employer le vocabulaire du sport et de la santé}

\begin{methode}[Mémoriser et utiliser le lexique du sport]
\begin{enumerate}
\item \textbf{Apprends les mots par familles} : les sports (\zh{跑步}, \zh{游泳}, \zh{打篮球}, \zh{踢足球}), le corps et la santé (\zh{身体}, \zh{健康}, \zh{锻炼}, \zh{休息}), les fréquences (\zh{每天}, \zh{常常}, \zh{有时候}).
\item \textbf{Mémorise le verbe qui accompagne chaque sport} : on dit \zh{打篮球} (jouer au basket, avec les mains) et \zh{踢足球} (jouer au football, avec le pied), jamais l'inverse. Pour \zh{跑步} et \zh{游泳}, on utilise \zh{去} + sport : \zh{去跑步}, \zh{去游泳}.
\item \textbf{Prononce avec les tons} : répète chaque mot trois fois à voix haute. Attention : \zh{yùndòng} (4\textsuperscript{e} + 4\textsuperscript{e} ton), \zh{jiànkāng} (4\textsuperscript{e} + 1\textsuperscript{er} ton).
\item \textbf{Fabrique une phrase personnelle} avec chaque mot nouveau : \zh{我每天锻炼半个小时。} (Wǒ měi tiān duànliàn bàn ge xiǎoshí.) « Je m'entraîne une demi-heure chaque jour. »
\item \textbf{Structure à retenir} : \zh{对 + nom/pronom + 好/有帮助} = « être bon pour / utile à » : \zh{运动对健康有好处。} (Yùndòng duì jiànkāng yǒu hǎochu.) « Le sport est bénéfique pour la santé. »
\end{enumerate}
\end{methode}

\begin{propriete}[La structure \zh{对……好} et \zh{对……有帮助}]
Schéma : \textbf{A + 对 + B + 好 / 有帮助 / 有好处} = « A est bon / utile / bénéfique pour B ».
\begin{itemize}
\item \zh{运动对身体很好。} (Yùndòng duì shēntǐ hěn hǎo.) « Le sport est très bon pour le corps. »
\item \zh{运动对学习也有帮助。} (Yùndòng duì xuéxí yě yǒu bāngzhù.) « Le sport est aussi utile pour les études. »
\end{itemize}
Le groupe \zh{对 + B} se place toujours \emph{avant} l'adjectif ou le verbe, jamais après.
\end{propriete}

\begin{exoresolu}[Vocabulaire : complète les phrases]
\textbf{Énoncé.} Complète avec le mot qui convient : \zh{锻炼} — \zh{踢} — \zh{打} — \zh{健康} — \zh{休息}.
\begin{enumerate}
\item 我们下午 \ldots\ldots 篮球。
\item 他每天跑步，身体很 \ldots\ldots 。
\item 运动以后，你应该 \ldots\ldots 一会儿。
\item 喀麦隆人很喜欢 \ldots\ldots 足球。
\item 早上 \ldots\ldots 对身体很好。
\end{enumerate}
\tcblower
\textbf{Solution détaillée.}
\begin{enumerate}
\item \zh{打} : le basket se joue avec les mains, donc \zh{打篮球}. Phrase : \zh{我们下午打篮球。} (Wǒmen xiàwǔ dǎ lánqiú.) « Nous jouons au basket l'après-midi. »
\item \zh{健康} : \zh{身体很健康} = « être en bonne santé ». Phrase : \zh{他每天跑步，身体很健康。} (Tā měi tiān pǎobù, shēntǐ hěn jiànkāng.) « Il court chaque jour, il est en très bonne santé. »
\item \zh{休息} : après le sport, on se repose. \zh{休息一会儿} = « se reposer un moment » (\zh{一会儿} exprime une courte durée). Phrase : \zh{运动以后，你应该休息一会儿。} (Yùndòng yǐhòu, nǐ yīnggāi xiūxi yíhuìr.) « Après le sport, tu devrais te reposer un moment. »
\item \zh{踢} : le football se joue avec le pied, donc \zh{踢足球}. Phrase : \zh{喀麦隆人很喜欢踢足球。} (Kāmàilóngrén hěn xǐhuan tī zúqiú.) « Les Camerounais aiment beaucoup jouer au football. »
\item \zh{锻炼} : \zh{锻炼} = « s'entraîner, faire de l'exercice », ici sujet de la phrase. Phrase : \zh{早上锻炼对身体很好。} (Zǎoshang duànliàn duì shēntǐ hěn hǎo.) « S'entraîner le matin est très bon pour le corps. »
\end{enumerate}
\textbf{Conclusion} : le choix du verbe dépend de l'objet : \zh{打} + ballon de main, \zh{踢} + ballon de pied, \zh{锻炼} seul ou avec une durée. C'est la collocation qui donne la réponse, pas la traduction française.
\end{exoresolu}

\courssub{Clés (radicaux) et ordre des traits}

\begin{propriete}[Les clés des caractères du thème]
La clé (radical) donne une indication de sens ; l'autre partie donne souvent une indication de prononciation (phonétique).
\begin{itemize}
\item \zh{跑} pǎo « courir » : clé \zh{足} « pied » à gauche ; la partie droite \zh{包} (bāo) indique le son.
\item \zh{泳} yǒng « nager » : clé \zh{氵} « eau » (trois gouttes) à gauche ; \zh{永} (yǒng) donne le son.
\item \zh{休} xiū « se reposer » : clé \zh{亻} « personne » + \zh{木} « arbre » : une personne contre un arbre se repose.
\item \zh{健} jiàn « en bonne santé » : clé \zh{亻} « personne » ; \zh{建} (jiàn) donne le son.
\item \zh{体} tǐ « corps » : clé \zh{亻} « personne » + \zh{本} (běn).
\end{itemize}
\end{propriete}

\begin{methode}[Analyser et écrire les caractères du sport]
\begin{enumerate}
\item \textbf{Identifie la clé (radical)} : elle donne le sens général. Dans \zh{跑} (courir), la clé est \zh{足} « pied » ; dans \zh{泳} (nager), la clé est \zh{氵} « eau » ; dans \zh{健} et \zh{休}, la clé est \zh{亻} « personne ».
\item \textbf{Respecte l'ordre des traits} : de haut en bas, de gauche à droite, l'horizontal avant le vertical, l'extérieur avant l'intérieur, on ferme la boîte en dernier. Pour les caractères gauche-droite comme \zh{跑}, \zh{泳}, \zh{健}, on écrit d'abord la clé à gauche, puis la partie droite.
\item \textbf{Décompose avant d'écrire} : \zh{休} = \zh{亻} + \zh{木} ; \zh{跑} = \zh{足} + \zh{包}.
\item \textbf{Écris chaque caractère cinq fois} en prononçant le pinyin à voix haute à chaque caractère terminé.
\item \textbf{Repère les pièges} : \zh{篮} (dans \zh{篮球}) a la clé bambou \zh{⺮} en haut ; ne pas confondre \zh{休} (xiū, se reposer) et \zh{体} (tǐ, corps).
\end{enumerate}
\end{methode}

\begin{exoresolu}[Caractères : clés et sens]
\textbf{Énoncé.} Pour chaque caractère, donne la clé (radical), le pinyin et le sens, puis explique le lien entre la clé et le sens : \zh{跑} — \zh{泳} — \zh{休} — \zh{健}.
\tcblower
\textbf{Solution détaillée.}
\begin{enumerate}
\item \zh{跑} pǎo « courir » : clé \zh{足} « pied » (à gauche). On court avec les pieds : la clé indique l'action du pied. La partie droite \zh{包} (bāo) donne une indication phonétique.
\item \zh{泳} yǒng « nager » : clé \zh{氵} « eau » (trois gouttes à gauche). Nager se fait dans l'eau. La partie droite \zh{永} (yǒng) donne le son.
\item \zh{休} xiū « se reposer » : clé \zh{亻} « personne ». Composition : une personne (\zh{亻}) appuyée contre un arbre (\zh{木}) se repose — caractère idéographique facile à retenir.
\item \zh{健} jiàn « en bonne santé » : clé \zh{亻} « personne ». La santé concerne la personne ; la partie droite \zh{建} (jiàn) donne le son.
\end{enumerate}
\textbf{Conclusion} : la clé sémantique (à gauche dans ces quatre caractères) oriente vers le sens, la partie phonétique vers la prononciation. Connaître les clés permet de deviner le sens de caractères inconnus et de chercher dans un dictionnaire.
\end{exoresolu}

\courssub{Méthode 3 — Traduire et construire des phrases sur le sport (ordre des mots)}

\begin{methode}[Traduire du français vers le chinois sans calquer]
\begin{enumerate}
\item \textbf{Repère le squelette} : Sujet + Verbe + Complément d'objet. Exemple : « Il fait du sport » → \zh{他运动。}
\item \textbf{Place les compléments de temps et de lieu AVANT le verbe} : « Il court le matin au stade » → \zh{他早上在体育场跑步。} (Tā zǎoshang zài tǐyùchǎng pǎobù.) Jamais après le verbe comme en français.
\item \textbf{Utilise les mots-outils correctement} : \zh{很} devant un adjectif attribut (\zh{身体很好}), \zh{得} après un verbe pour évaluer la manière (\zh{他跑得很快。} Tā pǎo de hěn kuài. « Il court vite »).
\item \textbf{Pour « non seulement… mais aussi »}, utilise la structure : Sujet + \zh{不但} + A + \zh{，而且} + B : \zh{运动不但对身体好，而且对学习也有帮助。}
\item \textbf{Vérifie les tons et la ponctuation chinoise} (，。) avant de rendre ta copie.
\end{enumerate}
\end{methode}

\begin{exoresolu}[Thème : traduis en chinois]
\textbf{Énoncé.} Traduis les phrases suivantes en chinois (caractères + pinyin).
\begin{enumerate}
\item Le sport est très bon pour la santé.
\item Le week-end, je joue souvent au football avec mes camarades à Douala.
\item Elle court très vite.
\end{enumerate}
\tcblower
\textbf{Solution détaillée.}
\begin{enumerate}
\item Structure « être bon pour » = \zh{对……很好}. \zh{运动对健康很好。} (Yùndòng duì jiànkāng hěn hǎo.) Le sujet \zh{运动} ouvre la phrase ; \zh{对健康} est un complément placé avant l'adjectif \zh{好}.
\item Le complément de temps \zh{周末} et le complément de lieu \zh{在杜阿拉} se placent avant le verbe ; \zh{和同学们一起} (avec les camarades) aussi. \zh{周末我常常和同学们一起在杜阿拉踢足球。} (Zhōumò wǒ chángcháng hé tóngxuémen yìqǐ zài Dù'ālā tī zúqiú.) Attention : \zh{踢足球}, pas \zh{打足球}.
\item Évaluation de la manière : verbe + \zh{得} + adverbe. \zh{她跑得很快。} (Tā pǎo de hěn kuài.) Le \zh{得} (avec le radical « pas » \zh{彳}) introduit le complément de manière ; ne pas le confondre avec \zh{的}.
\end{enumerate}
\textbf{Conclusion} : la traduction correcte exige de déplacer les compléments de temps et de lieu avant le verbe et de choisir le bon mot-outil (\zh{得} pour la manière, \zh{对} pour « pour »).
\end{exoresolu}

\courssub{Méthode 4 — Produire un court paragraphe sur sa pratique sportive}

\begin{methode}[Rédiger 4 à 6 phrases sur « mon sport »]
\begin{enumerate}
\item \textbf{Phrase 1 — présentation} : \zh{我叫……，我是……的学生。} (Je m'appelle…, je suis élève de…)
\item \textbf{Phrase 2 — le sport pratiqué + fréquence} : \zh{我每天/常常 + verbe de sport。} Exemple : \zh{我常常打篮球。}
\item \textbf{Phrase 3 — quand et avec qui} : compléments de temps et d'accompagnement avant le verbe : \zh{周末我和朋友们一起去游泳。}
\item \textbf{Phrase 4 — l'importance (le cœur du sujet)} : \zh{运动对身体很好，也让我更有精神。}
\item \textbf{Phrase 5 — conclusion} : \zh{我很喜欢运动。} ou \zh{运动很重要。}
\item \textbf{Relis} : vérifie l'ordre des mots, les tons du pinyin et la ponctuation chinoise.
\end{enumerate}
\end{methode}

\begin{exoresolu}[Production guidée : « Le sport et moi »]
\textbf{Énoncé.} En t'inspirant du texte de Li Ming, rédige un paragraphe de 4 à 5 phrases en chinois sur ta pratique sportive et l'importance du sport. Donne le pinyin et la traduction.
\tcblower
\textbf{Solution détaillée (modèle de réponse).}

\zh{我叫玛丽，我是杜阿拉一所高中的学生。我非常喜欢运动。每天下午，我和同学们一起踢足球。周末，我有时候去游泳。运动不但让我的身体更健康，而且让我上课更有精神。}

Wǒ jiào Mǎlì, wǒ shì Dù'ālā yì suǒ gāozhōng de xuésheng. Wǒ fēicháng xǐhuan yùndòng. Měi tiān xiàwǔ, wǒ hé tóngxuémen yìqǐ tī zúqiú. Zhōumò, wǒ yǒu shíhou qù yóuyǒng. Yùndòng bùdàn ràng wǒ de shēntǐ gèng jiànkāng, érqiě ràng wǒ shàngkè gèng yǒu jīngshen.

« Je m'appelle Marie, je suis élève dans un lycée de Douala. J'aime beaucoup le sport. Chaque après-midi, je joue au football avec mes camarades. Le week-end, je vais parfois nager. Le sport rend non seulement mon corps plus sain, mais il me donne aussi plus d'énergie en classe. »

\textbf{Justifications grammaticales} :
\begin{itemize}
\item \zh{每天下午} et \zh{周末} (temps) placés avant le verbe, conformément à l'ordre chinois.
\item \zh{和同学们一起} (accompagnement) placé avant le verbe \zh{踢}.
\item Structure \zh{不但……而且……} correctement appariée dans la dernière phrase.
\item \zh{让 + pronom/nom + adjectif} : \zh{让我的身体更健康} = « rend mon corps plus sain ».
\end{itemize}
\textbf{Conclusion} : un paragraphe réussi reprend les structures du texte support en changeant les informations personnelles — c'est la stratégie la plus sûre à l'examen.
\end{exoresolu}

\begin{experience}[Jeu de rôle : l'interview sportive]
Par groupes de deux : l'un joue un journaliste de la CRTV, l'autre un élève sportif. Le journaliste pose trois questions en chinois : \zh{你做什么运动？} \zh{你什么时候运动？} \zh{你觉得运动重要吗？} (Quel sport fais-tu ? Quand fais-tu du sport ? Trouves-tu le sport important ?). L'élève répond par des phrases complètes en utilisant \zh{常常}, \zh{和……一起} et \zh{对……好}. Puis on échange les rôles.
\end{experience}

\begin{savaistu}[Le sport en classe : Cameroun et Chine]
En Chine, la plupart des écoles organisent chaque matin des exercices collectifs (\zh{早操} zǎocāo, « gymnastique du matin ») dans la cour, en musique, avant le début des cours. Au Cameroun, c'est surtout le football qui rythme la vie des lycéens, avec les tournois inter-classes et les vacances sportives. Dans les deux pays, les enseignants répètent aux élèves que le sport aide à mieux étudier : \zh{运动对学习有帮助。}
\end{savaistu}

\begin{aretenir}[L'essentiel de la leçon]
\begin{itemize}
\item Vocabulaire clé : \zh{运动} (sport), \zh{跑步} (courir), \zh{游泳} (nager), \zh{打篮球} (jouer au basket), \zh{踢足球} (jouer au football), \zh{身体} (corps), \zh{健康} (santé), \zh{锻炼} (s'entraîner), \zh{休息} (se reposer), \zh{精神} (énergie).
\item Collocations : \zh{打} + basket, \zh{踢} + football, \zh{去} + \zh{跑步}/\zh{游泳}.
\item Structures : \zh{对……好/有帮助} (être bon/utile pour) ; \zh{不但……而且……} (non seulement… mais aussi) ; \zh{让 + pronom + adjectif} (rendre quelqu'un + adjectif) ; verbe + \zh{得} + adverbe (manière).
\item Ordre des mots : Sujet + temps + lieu + accompagnement + verbe.
\item Clés : \zh{足} (pied) dans \zh{跑}, \zh{氵} (eau) dans \zh{泳}, \zh{亻} (personne) dans \zh{休}, \zh{健}, \zh{体}.
\end{itemize}
\end{aretenir}

\begin{attention}[Les 5 erreurs qui coûtent des points à l'examen]
\begin{enumerate}
\item \textbf{Confondre \zh{打篮球} et \zh{踢足球}} : le basket se « frappe » (\zh{打}), le football se « botte » (\zh{踢}). Écrire \zh{打足球} est une faute de collocation systématiquement pénalisée.
\item \textbf{Placer le complément de temps après le verbe} (calque du français) : « Il court le matin » ne se traduit jamais par \zh{他跑步早上} mais par \zh{他早上跑步。} Temps et lieu AVANT le verbe.
\item \textbf{Confondre \zh{的}, \zh{得} et \zh{地}} : \zh{得} sert pour la manière après un verbe (\zh{他跑得很快。}), \zh{的} pour la possession ou la qualification (\zh{我的身体}). À l'écrit, ces trois caractères homophones (de) doivent être distingués.
\item \textbf{Négliger les tons} : \zh{yùndòng} (运动, sport) avec de mauvais tons devient incompréhensible ; \zh{jiànkāng} (健康) prononcé à plat perd son sens. À l'oral comme au pinyin écrit, chaque ton compte.
\item \textbf{Confondre des caractères proches} : \zh{休} (xiū, se reposer) et \zh{体} (tǐ, corps) partagent la clé \zh{亻} ; \zh{蓝} (lán, bleu) et \zh{篮} (lán, panier, dans \zh{篮球}) se distinguent par la clé bambou \zh{⺮}. Vérifie chaque trait avant de rendre ta copie.
\end{enumerate}
\end{attention}

\begin{exercice}[Pour t'entraîner]
\begin{enumerate}
\item Réponds en chinois par des phrases complètes : 你常常做什么运动？你觉得运动对身体好吗？
\item Traduis en chinois : « Le matin, je cours au stade avec mes amis. Le sport est bon pour la santé. »
\item Donne la clé et le sens des caractères : \zh{泳}, \zh{休}, \zh{跑}.
\end{enumerate}
\end{exercice}

\begin{corrige}[Exercice « Pour t'entraîner »]
\begin{enumerate}
\item Modèle de réponse : \zh{我常常打篮球。我觉得运动对身体很好。} (Wǒ chángcháng dǎ lánqiú. Wǒ juéde yùndòng duì shēntǐ hěn hǎo.) « Je joue souvent au basket. Je trouve que le sport est très bon pour le corps. » Toute réponse avec un sport correctement collé à son verbe et la structure \zh{对……好} est acceptée.
\item \zh{早上我和朋友们一起在体育场跑步。运动对健康很好。} (Zǎoshang wǒ hé péngyoumen yìqǐ zài tǐyùchǎng pǎobù. Yùndòng duì jiànkāng hěn hǎo.) Points de vigilance : temps \zh{早上} et lieu \zh{在体育场} avant le verbe ; \zh{和朋友们一起} avant le verbe.
\item \zh{泳} : clé \zh{氵} (eau), « nager » ; \zh{休} : clé \zh{亻} (personne), « se reposer » ; \zh{跑} : clé \zh{足} (pied), « courir ».
\end{enumerate}
\end{corrige}

\newpage

\coursec{Exercices}

\courssub{Je vérifie mes connaissances}

\begin{exercice}[QCM : le vocabulaire du sport]
Pour chaque question, entoure la bonne réponse.
\begin{enumerate}
\item \zh{运动} (yùndòng) signifie :
\begin{enumerate}
\item la santé \item le sport \item la danse
\end{enumerate}
\item \zh{健康} (jiànkāng) signifie :
\begin{enumerate}
\item malade \item fatigué \item en bonne santé
\end{enumerate}
\item \zh{跑步} (pǎobù) signifie :
\begin{enumerate}
\item courir \item nager \item sauter
\end{enumerate}
\item \zh{每天} (měitiān) signifie :
\begin{enumerate}
\item chaque semaine \item chaque jour \item parfois
\end{enumerate}
\item \zh{身体} (shēntǐ) signifie :
\begin{enumerate}
\item l'esprit \item le corps \item le cœur
\end{enumerate}
\end{enumerate}
\end{exercice}

\begin{exercice}[Vrai ou faux ?]
D'après le texte support de la leçon, indique si chaque phrase est \textbf{vraie} ou \textbf{fausse}, puis justifie ta réponse en recopiant une phrase du texte.
\begin{enumerate}
\item Le sport est bon pour le corps et pour l'esprit.
\item Xiao Ming n'aime pas du tout le sport.
\item Xiao Ming court tous les matins avant les cours.
\item Faire du sport fatigue trop les élèves, ils ne doivent pas en faire.
\item Les élèves devraient faire du sport chaque jour.
\end{enumerate}
\end{exercice}

\begin{exercice}[Complète le tableau de vocabulaire]
Complète le tableau : pour chaque caractère donné, écris le pinyin, la nature et la traduction ; pour chaque traduction donnée, retrouve les caractères.
\begin{center}
\begin{tabularx}{\linewidth}{|X|X|X|X|}\hline
\rowcolor{popblueL} Caractères & Pinyin & Nature & Traduction \\ \hline
\zh{游泳} & \ldots\ldots & \ldots\ldots & \ldots\ldots \\ \hline
\zh{踢足球} & \ldots\ldots & \ldots\ldots & \ldots\ldots \\ \hline
\ldots\ldots & \ldots\ldots & \ldots\ldots & la santé \\ \hline
\ldots\ldots & \ldots\ldots & \ldots\ldots & faire de l'exercice, s'entraîner \\ \hline
\zh{篮球} & \ldots\ldots & \ldots\ldots & \ldots\ldots \\ \hline
\ldots\ldots & \ldots\ldots & \ldots\ldots & important \\ \hline
\end{tabularx}
\end{center}
\end{exercice}

\begin{exercice}[Associe le pinyin aux caractères]
Relie chaque mot en caractères au bon pinyin.
\begin{center}
\begin{tabularx}{\linewidth}{|X|X|}\hline
\rowcolor{popblueL} Caractères & Pinyin \\ \hline
1. \zh{跳舞} & a. duànliàn \\ \hline
2. \zh{锻炼} & b. tiàowǔ \\ \hline
3. \zh{比赛} & c. xǐhuan \\ \hline
4. \zh{喜欢} & d. bǐsài \\ \hline
5. \zh{应该} & e. yīnggāi \\ \hline
\end{tabularx}
\end{center}
\end{exercice}

\courssub{Je m'entraîne}

\begin{exercice}[Traduis en chinois]
Traduis les phrases suivantes en chinois (caractères et pinyin).
\begin{enumerate}
\item Le sport est bon pour la santé.
\item Je cours chaque matin.
\item Les élèves devraient faire plus de sport.
\item Mon ami aime jouer au basket-ball.
\end{enumerate}
\end{exercice}

\begin{exercice}[Relie chaque sport au bon verbe]
Relie chaque verbe au nom de sport qui convient, puis écris les quatre expressions complètes avec leur traduction.
\begin{center}
\begin{tabularx}{\linewidth}{|X|X|}\hline
\rowcolor{popblueL} Verbe & Sport \\ \hline
\zh{打} & \zh{足球} \\ \hline
\zh{踢} & \zh{泳} \\ \hline
\zh{跳} & \zh{篮球} \\ \hline
\zh{游} & \zh{舞} \\ \hline
\end{tabularx}
\end{center}
\end{exercice}

\begin{textebox}[Texte à compléter (support de l'exercice suivant)]
\zh{我（一）\ldots\ldots 运动，因为运动对（二）\ldots\ldots 很好。我（三）\ldots\ldots 早上跑步，下午和朋友们一起踢足球。（四）\ldots\ldots 让我更（五）\ldots\ldots ，也让我更快乐。}\\[4pt]
\emph{Wǒ (yī) \ldots\ yùndòng, yīnwèi yùndòng duì (èr) \ldots\ hěn hǎo. Wǒ (sān) \ldots\ zǎoshang pǎobù, xiàwǔ hé péngyoumen yìqǐ tī zúqiú. (Sì) \ldots\ ràng wǒ gèng (wǔ) \ldots, yě ràng wǒ gèng kuàilè.}\\[4pt]
« Je \ldots\ (1) le sport, parce que le sport est très bon pour \ldots\ (2). \ldots\ (3) matin, je cours ; l'après-midi, je joue au football avec mes amis. \ldots\ (4) me rend plus \ldots\ (5) et me rend aussi plus heureux. »
\end{textebox}

\begin{exercice}[Texte à trous]
Complète le texte ci-dessus avec les mots de la liste : \zh{健康}、\zh{每天}、\zh{运动}、\zh{身体}、\zh{喜欢}。Attention : chaque mot ne peut être utilisé qu'une seule fois.
\end{exercice}

\begin{exercice}[Remets les mots dans l'ordre]
Remets les mots dans le bon ordre pour former des phrases correctes, puis traduis-les en français.
\begin{enumerate}
\item \zh{对 / 运动 / 健康 / 很 / 有好处}
\item \zh{早上 / 我 / 跑步 / 每天 / 在 / 学校}
\item \zh{应该 / 多 / 学生们 / 锻炼 / 身体}
\item \zh{喜欢 / 打 / 他 / 最 / 篮球}
\end{enumerate}
\end{exercice}

\courssub{Je me prépare à l'évaluation}

\begin{textebox}[Le sport à l'école (support de l'exercice suivant)]
\zh{运动对我们的健康非常重要。在雅温得的一所中学里，学生们每天下午四点以后做运动。有的学生踢足球，有的学生打篮球，还有的学生跑步。老师说：« 运动不但对身体好，而且对学习也有帮助。 » 学生们都觉得运动让他们更快乐、更健康。学校每年还有一次运动会，大家都很喜欢。}\\[4pt]
\emph{Yùndòng duì wǒmen de jiànkāng fēicháng zhòngyào. Zài Yǎwēndé de yí suǒ zhōngxué li, xuéshengmen měitiān xiàwǔ sì diǎn yǐhòu zuò yùndòng. Yǒu de xuésheng tī zúqiú, yǒu de xuésheng dǎ lánqiú, hái yǒu de xuésheng pǎobù. Lǎoshī shuō : « Yùndòng búdàn duì shēntǐ hǎo, érqiě duì xuéxí yě yǒu bāngzhù. » Xuéshengmen dōu juéde yùndòng ràng tāmen gèng kuàilè, gèng jiànkāng. Xuéxiào měinián hái yǒu yí cì yùndònghuì, dàjiā dōu hěn xǐhuan.}\\[4pt]
« Le sport est très important pour notre santé. Dans un lycée de Yaoundé, les élèves font du sport chaque jour après seize heures. Certains élèves jouent au football, d'autres jouent au basket-ball, d'autres encore courent. Le professeur dit : "Le sport est non seulement bon pour le corps, mais il aide aussi les études." Les élèves trouvent tous que le sport les rend plus heureux et en meilleure santé. L'école organise aussi une rencontre sportive chaque année, et tout le monde l'apprécie beaucoup. »
\end{textebox}

\begin{exercice}[Type évaluation : compréhension de texte]
Lis le texte ci-dessus, puis réponds aux questions.
\textbf{Questions de compréhension (en français) :}
\begin{enumerate}
\item Pourquoi le sport est-il important selon le texte ?
\item À quelle heure les élèves font-ils du sport ?
\item Cite les trois sports pratiqués par les élèves.
\item Que pense le professeur du sport ?
\item Qu'est-ce que l'école organise chaque année ?
\end{enumerate}
\textbf{Questions de langue (en chinois, réponds en chinois) :}
\begin{enumerate}
\item \zh{学生们什么时候做运动？} \emph{(Xuéshengmen shénme shíhou zuò yùndòng ?)}
\item \zh{运动对学生有什么好处？} \emph{(Yùndòng duì xuésheng yǒu shénme hǎochu ?)}
\end{enumerate}
\end{exercice}

\begin{exercice}[Type évaluation : thème et version]
\textbf{A. Thème.} Traduis en chinois (caractères et pinyin) :
\begin{enumerate}
\item Le sport est très important pour la santé.
\item Chaque matin, je cours avec mes amis.
\item Les élèves devraient faire du sport tous les jours.
\end{enumerate}
\textbf{B. Version.} Traduis en français :
\begin{enumerate}
\item \zh{运动让我的身体更健康。} \emph{(Yùndòng ràng wǒ de shēntǐ gèng jiànkāng.)}
\item \zh{他最喜欢踢足球，也喜欢游泳。} \emph{(Tā zuì xǐhuan tī zúqiú, yě xǐhuan yóuyǒng.)}
\item \zh{我们应该每天锻炼身体。} \emph{(Wǒmen yīnggāi měitiān duànliàn shēntǐ.)}
\end{enumerate}
\end{exercice}

\begin{exercice}[Type évaluation : expression écrite]
\textbf{Sujet :} \zh{我最喜欢的运动} \emph{(Wǒ zuì xǐhuan de yùndòng)} — « Mon sport préféré ».

Rédige un court texte en chinois de \textbf{5 à 8 phrases} (au moins 40 caractères) en suivant la trame vue en classe :
\begin{enumerate}
\item Présente ton sport préféré : \zh{我最喜欢的运动是\ldots}
\item Dis quand et avec qui tu le pratiques : \zh{我每天 / 每个星期\ldots 和\ldots 一起\ldots}
\item Explique pourquoi tu l'aimes (santé, plaisir, amis) : \zh{因为\ldots}
\item Conclus sur l'importance du sport : \zh{运动对\ldots 很重要。}
\end{enumerate}
Tu utiliseras au moins cinq mots du vocabulaire de la leçon (\zh{运动、健康、身体、锻炼、喜欢、应该}\ldots) et tu écriras le pinyin sous chaque phrase.
\end{exercice}

\courssub{Corrigés}



\newpage

\coursec{Corrigés des exercices}

\begin{corrige}[Exercice 1 — QCM : le vocabulaire du sport]
1. \textbf{b} — \zh{运动} (yùndòng) signifie « le sport ». 2. \textbf{c} — \zh{健康} (jiànkāng) signifie « en bonne santé ». 3. \textbf{a} — \zh{跑步} (pǎobù) signifie « courir ». 4. \textbf{b} — \zh{每天} (měitiān) signifie « chaque jour ». 5. \textbf{b} — \zh{身体} (shēntǐ) signifie « le corps ».

\textbf{Points de vigilance :} ne pas confondre \zh{健康} (jiànkāng, la santé) et \zh{身体} (shēntǐ, le corps) ; attention au troisième ton de \zh{跑} (pǎo) et au ton descendant de \zh{运动} (yùndòng, 4\textsuperscript{e} + 4\textsuperscript{e} ton).
\end{corrige}

\begin{corrige}[Exercice 2 — Vrai ou faux ?]
\begin{enumerate}
\item \textbf{Vrai.} Le texte dit : \zh{运动对身体好，也对心情好。} \emph{(Yùndòng duì shēntǐ hǎo, yě duì xīnqíng hǎo.)} « Le sport est bon pour le corps et aussi pour le moral. »
\item \textbf{Faux.} Le texte dit : \zh{小明很喜欢运动。} \emph{(Xiǎo Míng hěn xǐhuan yùndòng.)} « Xiao Ming aime beaucoup le sport. »
\item \textbf{Vrai.} Le texte dit : \zh{他每天早上上课以前跑步。} \emph{(Tā měitiān zǎoshang shàngkè yǐqián pǎobù.)} « Chaque matin, avant les cours, il court. »
\item \textbf{Faux.} Le texte affirme le contraire : le sport est bon pour la santé et les élèves devraient en faire : \zh{我们应该每天做运动。}
\item \textbf{Vrai.} Le texte dit : \zh{我们应该每天做运动。} \emph{(Wǒmen yīnggāi měitiān zuò yùndòng.)} « Nous devrions faire du sport chaque jour. »
\end{enumerate}
\textbf{Méthode :} pour justifier, on recopie toujours une phrase complète du texte, jamais une simple impression personnelle.
\end{corrige}

\begin{corrige}[Exercice 3 — Complète le tableau de vocabulaire]
\begin{center}
\begin{tabularx}{\linewidth}{|X|X|X|X|}\hline
\rowcolor{popblueL} Caractères & Pinyin & Nature & Traduction \\ \hline
\zh{游泳} & yóuyǒng & verbe & nager \\ \hline
\zh{踢足球} & tī zúqiú & locution verbale & jouer au football \\ \hline
\zh{健康} & jiànkāng & nom / adjectif & la santé \\ \hline
\zh{锻炼} & duànliàn & verbe & faire de l'exercice, s'entraîner \\ \hline
\zh{篮球} & lánqiú & nom & basket-ball \\ \hline
\zh{重要} & zhòngyào & adjectif & important \\ \hline
\end{tabularx}
\end{center}
\textbf{Points de vigilance :} \zh{锻炼} (duànliàn) porte deux quatrièmes tons ; \zh{游泳} s'écrit avec la clé de l'eau \zh{氵} à gauche dans les deux caractères ; ne pas confondre \zh{重要} (zhòngyào, important) et \zh{主要} (zhǔyào, principal).
\end{corrige}

\begin{corrige}[Exercice 4 — Associe le pinyin aux caractères]
1 -- b (\zh{跳舞} tiàowǔ, danser) ; 2 -- a (\zh{锻炼} duànliàn, s'entraîner) ; 3 -- d (\zh{比赛} bǐsài, match, compétition) ; 4 -- c (\zh{喜欢} xǐhuan, aimer) ; 5 -- e (\zh{应该} yīnggāi, devoir).

\textbf{Points de vigilance :} dans \zh{喜欢}, le second caractère \zh{欢} se lit au ton neutre (huan) ; dans \zh{应该}, \zh{应} se lit ici au premier ton (yīng) et non au quatrième.
\end{corrige}

\begin{corrige}[Exercice 5 — Traduis en chinois]
\begin{enumerate}
\item \zh{运动对健康很好。} \emph{(Yùndòng duì jiànkāng hěn hǎo.)} Structure : \cle{A + 对 + B + 很 + adjectif} « A est bon pour B ».
\item \zh{我每天早上跑步。} \emph{(Wǒ měitiān zǎoshang pǎobù.)} Le complément de temps \zh{每天早上} se place juste après le sujet, avant le verbe.
\item \zh{学生们应该多做运动。} \emph{(Xuéshengmen yīnggāi duō zuò yùndòng.)} \zh{应该} (devoir) précède le verbe ; \zh{多} (davantage) se place avant le verbe \zh{做}.
\item \zh{我的朋友喜欢打篮球。} \emph{(Wǒ de péngyou xǐhuan dǎ lánqiú.)} Pour le basket, on dit \zh{打篮球} et non \zh{踢篮球}.
\end{enumerate}
\textbf{Erreur fréquente :} placer le complément de temps à la fin de la phrase comme en français (« je cours chaque matin ») ; en chinois, il se met avant le verbe.
\end{corrige}

\begin{corrige}[Exercice 6 — Relie chaque sport au bon verbe]
\begin{itemize}
\item \zh{打篮球} \emph{(dǎ lánqiú)} : jouer au basket-ball.
\item \zh{踢足球} \emph{(tī zúqiú)} : jouer au football.
\item \zh{跳舞} \emph{(tiàowǔ)} : danser.
\item \zh{游泳} \emph{(yóuyǒng)} : nager.
\end{itemize}
\textbf{Règle à retenir :} on utilise \zh{打} (frapper avec la main) pour les sports de ballon joués avec les mains (\zh{打篮球、打排球、打乒乓球}) ; \zh{踢} (donner un coup de pied) pour le football ; \zh{游} se combine avec \zh{泳} dans le mot unique \zh{游泳}. Dire \zh{踢篮球} est une erreur classique de francophone.
\end{corrige}

\begin{corrige}[Exercice 7 — Texte à trous]
(一) \zh{喜欢} ; (二) \zh{身体} ; (三) \zh{每天} ; (四) \zh{运动} ; (五) \zh{健康}.

Texte complet : \zh{我喜欢运动，因为运动对身体很好。我每天早上跑步，下午和朋友们一起踢足球。运动让我更健康，也让我更快乐。}

\emph{Wǒ xǐhuan yùndòng, yīnwèi yùndòng duì shēntǐ hěn hǎo. Wǒ měitiān zǎoshang pǎobù, xiàwǔ hé péngyoumen yìqǐ tī zúqiú. Yùndòng ràng wǒ gèng jiànkāng, yě ràng wǒ gèng kuàilè.}

« J'aime le sport, parce que le sport est très bon pour le corps. Chaque matin, je cours ; l'après-midi, je joue au football avec mes amis. Le sport me rend plus sain et me rend aussi plus heureux. »

\textbf{Justification :} en (二), la structure \zh{对……很好} exige un nom : \zh{身体} convient (et \zh{健康} est réservé à la case 五, où l'adjectif suit \zh{更}) ; en (三), seul \zh{每天} peut précéder \zh{早上}.
\end{corrige}

\begin{corrige}[Exercice 8 — Remets les mots dans l'ordre]
\begin{enumerate}
\item \zh{运动对健康很有好处。} \emph{(Yùndòng duì jiànkāng hěn yǒu hǎochu.)} « Le sport est très bénéfique pour la santé. » Structure : \cle{A + 对 + B + 很有好处}.
\item \zh{我每天早上在学校跑步。} \emph{(Wǒ měitiān zǎoshang zài xuéxiào pǎobù.)} « Chaque matin, je cours à l'école. » Ordre : sujet + temps + lieu (\zh{在} + lieu) + verbe.
\item \zh{学生们应该多锻炼身体。} \emph{(Xuéshengmen yīnggāi duō duànliàn shēntǐ.)} « Les élèves devraient s'entraîner davantage. » \zh{应该} puis \zh{多} précèdent le verbe.
\item \zh{他最喜欢打篮球。} \emph{(Tā zuì xǐhuan dǎ lánqiú.)} « Ce qu'il préfère, c'est jouer au basket-ball. » \zh{最} se place directement avant \zh{喜欢}.
\end{enumerate}
\textbf{Point de vigilance :} en chinois, l'ordre est Sujet + Temps + Lieu + Verbe ; l'inverse du français (« je cours à l'école chaque matin »).
\end{corrige}

\begin{corrige}[Exercice 9 — Type évaluation : compréhension de texte]
\textbf{Barème indicatif (sur 10) :} 5 questions de compréhension × 1,5 pt = 7,5 pts ; 2 questions de langue × 1,25 pt = 2,5 pts. Toute réponse juste mais mal formulée en chinois perd la moitié des points.

\textbf{Questions de compréhension :}
\begin{enumerate}
\item Parce que le sport est très important pour la santé : \zh{运动对我们的健康非常重要。}
\item Ils font du sport chaque jour après seize heures : \zh{每天下午四点以后做运动。}
\item Le football (\zh{踢足球}), le basket-ball (\zh{打篮球}) et la course (\zh{跑步}).
\item Il pense que le sport est non seulement bon pour le corps, mais qu'il aide aussi les études : \zh{运动不但对身体好，而且对学习也有帮助。}
\item L'école organise une rencontre sportive chaque année : \zh{学校每年还有一次运动会。}
\end{enumerate}

\textbf{Questions de langue :}
\begin{enumerate}
\item \zh{学生们每天下午四点以后做运动。} \emph{(Xuéshengmen měitiān xiàwǔ sì diǎn yǐhòu zuò yùndòng.)}
\item \zh{运动对学生的身体好，对学习也有帮助。} \emph{(Yùndòng duì xuésheng de shēntǐ hǎo, duì xuéxí yě yǒu bāngzhù.)}
\end{enumerate}
\textbf{Points de vigilance :} répondre par une phrase complète (sujet + verbe), pas par un mot isolé ; repérer la structure corrélée \zh{不但……而且……} « non seulement\ldots mais aussi\ldots ».
\end{corrige}

\begin{corrige}[Exercice 10 — Type évaluation : thème et version]
\textbf{Barème indicatif (sur 10) :} thème 5 pts (environ 1,6 pt par phrase) ; version 5 pts. En thème, une faute de caractère ou d'ordre des mots coûte 0,5 pt.

\textbf{A. Thème :}
\begin{enumerate}
\item \zh{运动对健康非常重要。} \emph{(Yùndòng duì jiànkāng fēicháng zhòngyào.)}
\item \zh{每天早上，我和朋友们一起跑步。} \emph{(Měitiān zǎoshang, wǒ hé péngyoumen yìqǐ pǎobù.)}
\item \zh{学生们应该每天做运动。} \emph{(Xuéshengmen yīnggāi měitiān zuò yùndòng.)}
\end{enumerate}

\textbf{B. Version :}
\begin{enumerate}
\item « Le sport rend mon corps plus sain. » (\zh{让} = faire, rendre ; \zh{更} = davantage, plus.)
\item « Ce qu'il préfère, c'est jouer au football, et il aime aussi nager. » (\zh{最喜欢} = aimer le plus, préférer.)
\item « Nous devrions faire de l'exercice chaque jour. » (\zh{锻炼身体} = s'entraîner, littéralement « forger le corps ».)
\end{enumerate}
\textbf{Points de vigilance :} ne pas traduire \zh{非常} par « très » en le plaçant après l'adjectif ; en thème, le complément de temps \zh{每天早上} se met avant le sujet ou juste après, jamais à la fin.
\end{corrige}

\begin{corrige}[Exercice 11 — Type évaluation : expression écrite]
\textbf{Barème indicatif (sur 10) :} respect de la consigne (5 à 8 phrases, au moins 40 caractères) : 2 pts ; emploi d'au moins cinq mots du vocabulaire de la leçon : 2 pts ; correction grammaticale (ordre des mots, structures) : 3 pts ; exactitude des caractères et du pinyin : 2 pts ; richesse et cohérence : 1 pt.

\textbf{Proposition de rédaction modèle :}

\zh{我最喜欢的运动是踢足球。我每天下午和朋友们一起在学校踢足球。因为踢足球很有意思，也对我的身体很好，所以我很喜欢。运动让我更健康、更快乐。我们应该每天锻炼身体。运动对每个人都很重要。}

\emph{Wǒ zuì xǐhuan de yùndòng shì tī zúqiú. Wǒ měitiān xiàwǔ hé péngyoumen yìqǐ zài xuéxiào tī zúqiú. Yīnwèi tī zúqiú hěn yǒu yìsi, yě duì wǒ de shēntǐ hěn hǎo, suǒyǐ wǒ hěn xǐhuan. Yùndòng ràng wǒ gèng jiànkāng, gèng kuàilè. Wǒmen yīnggāi měitiān duànliàn shēntǐ. Yùndòng duì měi ge rén dōu hěn zhòngyào.}

« Mon sport préféré est le football. Chaque après-midi, je joue au football avec mes amis à l'école. Parce que le football est très amusant et qu'il est très bon pour mon corps, je l'aime beaucoup. Le sport me rend plus sain et plus heureux. Nous devrions faire de l'exercice chaque jour. Le sport est très important pour chacun. »

\textbf{Points de vigilance :}
\begin{itemize}
\item Utiliser la trame : \zh{我最喜欢的运动是……} / \zh{我每天……和……一起……} / \zh{因为……所以……} / \zh{运动对……很重要。}
\item Vérifier les collocations verbe + sport : \zh{踢足球}, \zh{打篮球}, \zh{游泳}, \zh{跑步}.
\item Placer le temps et le lieu avant le verbe ; ne pas oublier \zh{一起} après \zh{和朋友们}.
\item Écrire le pinyin sous chaque phrase avec les marques de tons, et soigner les caractères (clé de l'eau \zh{氵} dans \zh{游泳}, clé du pied \zh{⻊} dans \zh{踢}).
\end{itemize}
\end{corrige}

\newpage

\coursec{Fiche bilan}

\begin{popfigure}
\begin{tikzpicture}[mindmap, grow cyclic, every node/.style={concept, font=\small},
  root concept/.append style={concept color=popblue, font=\bfseries\small, minimum size=3.2cm},
  level 1/.append style={level distance=3.6cm, sibling angle=51.4, minimum size=2.5cm, font=\footnotesize},
  level 2/.append style={level distance=2.4cm, sibling angle=40, minimum size=1.9cm, font=\scriptsize}]
\node[root concept, align=center]{运动与健康\\ \emph{yùndòng yǔ jiànkāng}\\ Le sport et la santé}
  child[concept color=poporange]{ node {Lexique du sport}
    child { node[align=center]{运动 跑步 游泳\\ 打球 锻炼} }
    child { node[align=center]{身体 健康\\ 心脏 体重} }
  }
  child[concept color=popgreen]{ node {Lexique de la santé}
    child { node[align=center]{好处 重要\\ 精神 压力} }
    child { node[align=center]{生病 休息\\ 睡眠} }
  }
  child[concept color=poppurple]{ node {Grammaire clé}
    child { node[align=center]{对……有好处\\ \emph{duì... yǒu hǎochu}} }
    child { node[align=center]{越来越 + adj.\\ 应该 / 可以} }
  }
  child[concept color=poppink]{ node {Clés et traits}
    child { node[align=center]{⻊ 足 (pied)\\ 氵 (eau) 忄 (cœur)} }
    child { node[align=center]{haut $\to$ bas\\ gauche $\to$ droite} }
  }
  child[concept color=popgold]{ node {Compréhension}
    child { node[align=center]{repérer l'idée\\ principale} }
    child { node[align=center]{répondre en\\ phrase complète} }
  }
  child[concept color=popblue!60]{ node {Collocations}
    child { node[align=center]{做运动 锻炼身体\\ 参加比赛} }
    child { node[align=center]{对健康\\ 很重要} }
  }
  child[concept color=popgreen!70]{ node {Culture}
    child { node[align=center]{太极 乒乓球\\ en Chine} }
    child { node[align=center]{football\\ au Cameroun} }
  };
\end{tikzpicture}
\legende{Carte mentale de la leçon : le sport et la santé (运动与健康)}
\end{popfigure}

\begin{bilan}
\textbf{1. Lexique clé à connaître par cœur}
\begin{itemize}
  \item \zh{运动} \emph{yùndòng} : sport, exercice ; \zh{锻炼} \emph{duànliàn} : s'entraîner ; \zh{身体} \emph{shēntǐ} : le corps ; \zh{健康} \emph{jiànkāng} : santé, en bonne santé.
  \item \zh{跑步} \emph{pǎobù} : courir ; \zh{游泳} \emph{yóuyǒng} : nager ; \zh{打球} \emph{dǎ qiú} : jouer au ballon ; \zh{比赛} \emph{bǐsài} : match, compétition.
  \item \zh{好处} \emph{hǎochu} : avantage, bienfait ; \zh{重要} \emph{zhòngyào} : important ; \zh{压力} \emph{yālì} : stress ; \zh{精神} \emph{jīngshén} : esprit, énergie.
  \item \zh{每天} \emph{měitiān} : chaque jour ; \zh{坚持} \emph{jiānchí} : persévérer ; \zh{休息} \emph{xiūxi} : se reposer.
\end{itemize}

\textbf{2. Règles de grammaire à connaître par cœur}
\begin{center}
\begin{tabularx}{\linewidth}{|X|X|X|}
\hline
\rowcolor{popblueL} Structure & Exemple & Traduction \\
\hline
A \zh{对} B \zh{有好处} (A est bénéfique à B) & \zh{运动对健康有好处。} \emph{Yùndòng duì jiànkāng yǒu hǎochu.} & Le sport est bon pour la santé. \\
\hline
\zh{越来越} + adjectif (de plus en plus) & \zh{他的身体越来越好。} \emph{Tā de shēntǐ yuèláiyuè hǎo.} & Il est de mieux en mieux en forme. \\
\hline
\zh{应该} + verbe (devoir, conseil) & \zh{我们应该每天锻炼。} \emph{Wǒmen yīnggāi měitiān duànliàn.} & Nous devrions nous entraîner chaque jour. \\
\hline
\zh{可以} + verbe (possibilité, permission) & \zh{跑步可以减少压力。} \emph{Pǎobù kěyǐ jiǎnshǎo yālì.} & Courir peut réduire le stress. \\
\hline
\zh{因为……所以……} (parce que... donc...) & \zh{因为运动很重要，所以我每天跑步。} \emph{Yīnwèi yùndòng hěn zhòngyào, suǒyǐ wǒ měitiān pǎobù.} & Parce que le sport est important, je cours chaque jour. \\
\hline
\end{tabularx}
\end{center}

\textbf{3. Clés (radicaux) du thème}
\begin{itemize}
  \item \zh{跑} : clé du \cle{pied} ⻊ (足) à gauche — idée de mouvement des jambes.
  \item \zh{游泳} : clé de l'\cle{eau} 氵(trois gouttes) à gauche.
  \item \zh{情、快} : clé du \cle{cœur} 忄 à gauche — sentiments, état d'esprit.
  \item Ordre des traits : de haut en bas, de gauche à droite, l'horizontal avant le vertical, l'extérieur avant l'intérieur.
\end{itemize}

\textbf{4. Méthodes express}
\begin{itemize}
  \item \cle{Lire un texte} : 1) lire le titre et la première phrase ; 2) repérer les mots connus ; 3) deviner les mots inconnus par la clé et le contexte ; 4) reformuler l'idée principale en une phrase.
  \item \cle{Répondre aux questions} : reprendre les mots de la question, répondre par une phrase complète (Sujet + Verbe + Complément), jamais par un seul mot.
  \item \cle{Mémoriser le lexique} : apprendre chaque mot avec sa \cle{collocation} (做运动, 参加比赛, 锻炼身体) et non isolément.
\end{itemize}
\end{bilan}

\begin{aretenir}[Checklist avant le Bac]
\begin{itemize}
  \item $\square$ Je sais écrire de mémoire : 运动, 锻炼, 健康, 身体, 比赛, 跑步, 游泳.
  \item $\square$ Je connais le pinyin et les tons de tout le lexique du thème.
  \item $\square$ Je sais employer la structure \zh{对……有好处} sans oublier \zh{对}.
  \item $\square$ Je place correctement \zh{越来越} devant l'adjectif.
  \item $\square$ Je distingue \zh{应该} (conseil) et \zh{可以} (possibilité).
  \item $\square$ Je sais construire une phrase de cause avec \zh{因为……所以……}.
  \item $\square$ Je reconnais les clés 足, 氵, 忄 et je respecte l'ordre des traits.
  \item $\square$ Je sais répondre aux questions de compréhension par des phrases complètes.
  \item $\square$ Je connais les collocations : \zh{做运动}, \zh{锻炼身体}, \zh{参加比赛}, \zh{减少压力}.
  \item $\square$ Je peux dire en trois phrases pourquoi le sport est important (santé, stress, esprit d'équipe).
  \item $\square$ Je connais deux éléments culturels : le tai-chi et le ping-pong en Chine, le football au Cameroun.
  \item $\square$ Je vérifie la ponctuation chinoise pleine chasse (，。？) dans mes productions écrites.
\end{itemize}
\end{aretenir}

\findecours

\end{document}
